\documentclass[
    11pt,
    twoside,
    openright,
    numbers=noenddot, 
    bibliography=totoc,
    captions=tableheading
]{scrbook}

% --- 1. LINGUA E FONT ---
\usepackage[T1]{fontenc}
\usepackage[italian]{babel}
\usepackage{lmodern} 
\usepackage[stretch=10, shrink=10]{microtype} % Migliora micro-spaziature

% --- 2. LAYOUT (GEOMETRY) ---
\usepackage{geometry}
\geometry{ % FORMATO A4
    a4paper,         % Sostituisce paperwidth e paperheight impostando automaticamente 210x297mm
    top=30mm,        % Aumentato per bilanciare l'altezza maggiore dell'A4
    bottom=30mm,
    inner=25mm,      % Leggermente aumentato per mantenere le proporzioni della pagina
    outer=20mm,
    headsep=10mm,
    footskip=14mm,
    bindingoffset=6mm % Invariato: lo spazio per la colla/cucitura rimane standard
}

% Parametri per evitare righe orfane e vedove (Qualità tipografica massima)
\widowpenalty=10000
\clubpenalty=10000
\displaywidowpenalty=10000

% --- 3. COLORI FORMALI (Spostati in alto per evitare l'errore "color undefined") ---
\usepackage[table]{xcolor} % Caricato prima di tikz per evitare conflitti xcdraw
\definecolor{formalblue}{RGB}{0, 51, 102} 
\definecolor{formalbackground}{gray}{0.95} 
\definecolor{physicsblue}{RGB}{24, 76, 120}
\definecolor{physicsgray}{RGB}{245, 247, 250}
\definecolor{solutiongray}{RGB}{100, 110, 120}

% --- 4. MATEMATICA E FISICA ---
\usepackage{amsmath} 
\usepackage{amsthm}  % <--- Fondamentale per thmtools
\usepackage{mathtools}
\usepackage{amsfonts}
\usepackage{amssymb}
\usepackage{bm}      
\usepackage{physics} 
\usepackage{siunitx} 
\usepackage{dsfont}
\usepackage{pifont}

\sisetup{output-decimal-marker={,}} 
% RISOLUZIONE CONFLITTO PHYSICS VS SIUNITX
\AtBeginDocument{\RenewCommandCopy\qty\SI}

\numberwithin{equation}{section}

% --- TEOREMI OTTIMIZZATI CON THMTOOLS ---
\usepackage{thmtools}
\declaretheoremstyle[
    headfont=\sffamily\bfseries,  % Nome in sans-serif grassetto nero
    notefont=\sffamily\bfseries\upshape, % Titolo tra parentesi in grassetto dritto nero
    bodyfont=\itshape,            % Testo in corsivo (\emph risalterà in dritto)
    headpunct={.},
    postheadspace=0.5em,
    spaceabove=12pt,
    spacebelow=12pt
]{stileteorema}

\declaretheorem[style=stileteorema, name=Teorema, parent=chapter]{teorema}
\declaretheorem[style=stileteorema, name=Definizione, parent=chapter]{definizione}
\declaretheorem[style=stileteorema, name=Proprietà, parent=chapter]{proprietà}

% --- AMBIENTE DIMOSTRAZIONE (IN NERO) ---
\renewenvironment{proof}
  {\par\vspace{0.5em}\noindent\textbf{\sffamily Dimostrazione.}}
  {\hfill$\blacksquare$\par\vspace{1em}}

% Scorciatoie per Meccanica Quantistica
\newcommand{\m}[1]{\mathcal{#1}}
\newcommand{\h}[1]{\hat{\mathcal{#1}}}
\newcommand{\hd}[1]{\hat{\mathcal{#1}}^\dagger}
\newcommand{\hatd}[1]{\hat{{#1}}^\dagger}
\newcommand{\avg}[1]{\left\langle #1 \right\rangle}
\newcommand{\hb}[1]{\hat{\bm{#1}}}
\renewcommand{\Im}{\operatorname{Im}}
\renewcommand{\Re}{\operatorname{Re}}

% Macro per giustificare passaggi nelle equazioni
\newcommand{\giustifica}[2][=]{\overset{\text{\scriptsize (#2)}}{#1}}


% --- 5. GRAFICA E TABELLE ---
\usepackage{graphicx}
\usepackage{tikz}
\usepackage{pgfplots}
\usepackage{tikz-3dplot}
\usepackage{float}
\usepackage{booktabs}
\usepackage{caption}
\usepackage{subcaption}
\usepackage{multicol}
\usepackage{enumitem}

\usetikzlibrary{circuits.ee.IEC, arrows.meta, calc, positioning, 3d, decorations.markings, angles, quotes, decorations.pathmorphing, svg.path}
\pgfplotsset{compat=1.18}

\captionsetup{
    font=small,
    labelfont=bf,
    format=hang,
    width=0.9\textwidth,
    labelsep=colon
}

% --- 6. PERSONALIZZAZIONE KOMA-SCRIPT (Stile Titoli e Pagine) ---
\setkomafont{chapter}{\Huge\bfseries}
\setkomafont{section}{\Large\bfseries}
\setkomafont{subsection}{\large\bfseries}
\setkomafont{pageheadfoot}{\normalfont\small\sffamily\color{gray}} 
\setkomafont{pagenumber}{\normalfont\bfseries\sffamily\color{black}} 

\usepackage[automark, headsepline=0.4pt]{scrlayer-scrpage}
\clearpairofpagestyles
\lehead{\leftmark} 
\rohead{\rightmark} 
\lefoot*{\pagemark} 
\rofoot*{\pagemark} 

% --- 7. NAVIGAZIONE ---
\usepackage{hyperref}
\hypersetup{
    colorlinks=true,
    linkcolor=black,
    citecolor=black,
    urlcolor=formalblue, % Uniformato ai colori del documento
    pdftitle={Dispense di Meccanica Quantistica},
    pdfauthor={Carlo Alberto Riva}
}
\renewcommand{\contentsname}{Indice}
\renewcommand{\tablename}{Tabella}

% --- 8. AMBIENTI PERSONALIZZATI (Box e Problemi) ---
\usepackage[most]{tcolorbox}

% Box Approfondimento (con indicatore di chiusura)
% Box Approfondimento con Bordo Completo
\newtcolorbox{approfondimento}[1][]{
    enhanced, breakable,
    colback=formalbackground, % Sfondo grigio chiaro
    colframe=formalblue,       % Cornice completa blu formale
    boxrule=2pt,               % Spessore uniforme su tutti e 4 i lati
    coltitle=formalblue,                
    fonttitle=\bfseries\sffamily\large, 
    title={Approfondimento},            
    sharp corners, detach title,                       
    before upper={\tcbtitle\par\vspace{0.3em}}, 
    left=1em, right=1em, top=1em, bottom=1em, 
    parbox=false,
    overlay unbroken and last={% Quadratino di fine box in basso a destra
        \node[anchor=south east, outer sep=2pt, color=formalblue!40] at (frame.south east) {\tiny$\blacksquare$};
    },
    #1
}

% 2. Box Concetto Chiave
% Definizione del colore per i concetti chiave (es. un rosso mattone scuro o oro vecchio)
\definecolor{keygold}{RGB}{153, 102, 0}
\definecolor{keybackground}{RGB}{253, 251, 245}

% Box Concetto Chiave / Equazione Fondamentale
% Box Concetto Chiave / Equazione Fondamentale con Titolo Personalizzato
% Box Concetto Chiave / Equazione Fondamentale con Bordo Spesso Completo
\newtcolorbox{concettochiave}[2][]{
    enhanced, breakable,
    colback=keybackground,    % Colore dello sfondo interno
    colframe=keygold,         % Colore della cornice (linea spessa)
    boxrule=2pt,              % Spessore uniforme su tutti e 4 i lati
    coltitle=keygold,                
    fonttitle=\bfseries\sffamily\normalsize, 
    title={#2},               
    sharp corners, detach title,                       
    before upper={\tcbtitle\par\vspace{0.4em}}, 
    left=1em, right=1em, top=0.8em, bottom=0.8em, 
    parbox=false,
    overlay unbroken and last={
        \node[anchor=south east, outer sep=2pt, color=keygold!40] at (frame.south east) {\tiny$\blacksquare$};
    },                       
    #1                        
}

% Contatore Problemi integrato per capitolo
\newcounter{probcount}[chapter]
\renewcommand{\theprobcount}{\thechapter.\arabic{probcount}}

% Box Testo Problema 
\newtcolorbox{problemabox}[1]{
    enhanced, breakable,
    colback=physicsgray, colframe=physicsblue,
    attach boxed title to top left={yshift=-2mm, xshift=5mm},
    boxed title style={colback=physicsblue},
    fonttitle=\bfseries\sffamily,
    code={\refstepcounter{probcount}\addcontentsline{toc}{subsection}{Prob. \theprobcount\ -- #1}},
    title={Problema \theprobcount \quad #1},
    sharp corners, boxrule=0.8pt,
    top=4mm, bottom=3mm, left=3mm, right=3mm,
    before skip=15pt, after skip=10pt
}

% Ambiente Soluzione
\newenvironment{soluzione}{%
    \par\vspace{5mm}\noindent
    \begin{tcolorbox}[
        enhanced, breakable, frame hidden, colback=white,
        borderline west={2pt}{0pt}{solutiongray}, sharp corners,
        left=4mm, right=0mm, top=0mm, bottom=0mm,
        before skip=0pt, after skip=15pt
    ]
    {\small\sffamily\bfseries\color{physicsblue}SOLUZIONE:}\par\nopagebreak\vspace{2mm}
}{
    \end{tcolorbox}\par
}

% ==========================================
% INIZIO DOCUMENTO
% ==========================================
\title{Dispense di Meccanica Quantistica}
\author{Carlo Alberto Riva}
\date{Politecnico di Torino}

\begin{document}

\frontmatter 
\maketitle 
\tableofcontents

\mainmatter 

% Inclusione capitoli
%!TEX root = ../dispense_QP.tex

\chapter{La fine della fisica classica}
All'inizio del secolo scorso, una serie di esperimenti rese clamorosamente evidente l'inadeguatezza della fisica classica nel fornire una spiegazione soddisfacente a un numero crescente di fenomeni su scala microscopica e atomica. L'adozione di un'interpretazione classica portava a risultati contraddittori nello studio di processi che coinvolgevano la radiazione ad alta frequenza, le proprietà di conduzione dei solidi, l'interazione tra radiazione e materia, l'apparente comportamento ondulatorio di elettroni e atomi e molti altri aspetti della fisica atomica e della materia condensata. La necessità di una nuova e più efficace base teorica, in grado di fornire un'interpretazione coerente di questa ricca fenomenologia, portò alla formulazione della \textbf{Meccanica Quantistica}. Tra i molti esperimenti in questo senso \emph{pionieristici}, vanno ricordati in particolare: \\

\noindent
\emph{La quantizzazione della radiazione}
\begin{enumerate}
    \item La misura dell'energia (per unità di volume e di frequenza) della Radiazione di Corpo Nero, la quale mostrò come la legge (classica) di Rayleigh-Jeans fallisse nel descrivere lo spettro ad alta frequenza, mentre la legge (quantistica) di Planck si adattasse perfettamente ai dati sperimentali (Planck, 1900).
    \item L'Effetto Fotoelettrico, che portò al postulato secondo cui la luce è composta da fotoni, i quali rappresentano quanti di energia $E = \hbar\omega$ (Einstein, 1905).
\end{enumerate}

\noindent \emph{La quantizzazione dei livelli atomici}
\begin{enumerate}
    \setcounter{enumi}{2}
    \item Il carattere discreto esibito dagli spettri atomici di emissione e di assorbimento.
    \item L'esperimento di Franck-Hertz (1914), che fornì l'evidenza diretta dell'esistenza di livelli energetici discreti all'interno degli atomi.
\end{enumerate}

\noindent \emph{Il dualismo onda-corpuscolo per la materia}
\begin{enumerate}
    \setcounter{enumi}{4}
    \item L'esperimento di Davisson-Germer (1927), in cui un fascio di elettroni, incidendo su un bersaglio (un reticolo cristallino con una costante di reticolo opportunamente piccola), subisce un effetto di diffrazione. Questo fenomeno dà origine a figure di interferenza del tutto analoghe a quelle tipiche della luce visibile o dei raggi X.
    \item L'esperimento di Thomson (1927-1928), in cui un fascio di elettroni, attraversando una sottile lamina d'oro, produce su una lastra fotografica le ben note figure di diffrazione (rappresentate da una serie di anelli concentrici).
\end{enumerate}

\section{La riformulazione di Schrödinger}
In un contesto quanto mai effervescente come quello appena descritto, si percepisce sempre più crescente la richiesta impellente di un \emph{cambio radicale di paradigma} per descrivere accuratamente la realtà. Fra le varie esperienze sperimentali, la sopracitata esperienza di Davisson e Germer ai Bell Labs ha in particolare dimostrato come gli \textbf{elettroni}, in certi contesti, esibiscano comportamenti tipici delle \emph{onde}. L'intuizione geniale che ebbe Schrödinger è stata quella di tentare di trattare gli elettroni stessi come delle onde caratterizzate da una equazione caratteristica che prende in argomento la ben nota \emph{funzione d'onda} $\psi(\bm{r},t)$, che ne descrive l'evoluzione spazio temporale:
\begin{equation}
    \psi(\bm{r},t) = A e^{i(\bm{k}\cdot\bm{r}- \omega t)} \quad.
    \label{eq:func_onda}
\end{equation}
Cerchiamo ora di decifrare i "simboli" coinvolti nell'equazione:
\begin{itemize}[noitemsep]
    \item $\bm{r}$: la generica posizione nello spazio\footnote{A meno di diversa indicazione, assumeremo sempre lo spazio come tridimensionale e coincidente con $\mathbb{R}^3$.};
    \item $t$: il generico istante di tempo;
    \item $A$: l'\emph{ampiezza} dell'onda;
    \item $i$: l'\emph{unità immaginaria}\footnote{Si ricordi l'uguaglianza fondamentale $i = \sqrt{-1}$.};
    \item $\bm{k}$: il \emph{vettore d'onda}. Si definisce come
    \[
        \bm{k} = \frac{2\pi}{\lambda} \bm{{u}}_r \quad,
    \]
    con $\lambda$ la \emph{lunghezza d'onda di De Broglie} caratteristica della particella\footnote{La lunghezza d'onda di De Broglie è parte inegrante della cosiddetta \emph{Ipotesi di De Broglie}, fondamento della fisica moderna che associa alle particelle caratteristiche ondulatorie. Vale la relazione $\lambda = {h}/{p}$, dove $h$ è la costante di Planck e $p$ la quantità di moto del corpuscolo.} e $\bm{{u}}_r$ il versore orientato secondo il verso di propagazione dell'onda; 
    \item $\omega$: la pulsazione d'onda definita come
    \[
        \omega = 2\pi\nu = \frac{2\pi v}{\lambda} \quad,
    \]
    con $\nu$ la frequenza d'onda, inverso del periodo e $v$ la velocità di propagazione.
\end{itemize}
L'interpretazione di Einstein dell'effetto fotoelettrico e lo studio di Planck sulla quantizzazione dell'energia ci permette inoltre di assumere l'energia degli elettroni come $E = \hbar \omega$, con $\hbar$ la costante di Planck ridotta ($\hbar = h/2\pi$). 

Siamo pronti per derivare, almeno in termini euristici, la fantomatica equazione di Schrödinger. Per una discussione meno edulcorata si veda l'appendice. 

\subsection{L'ipotesi elettromagnetica e l'equazione di D'Alembert}
Il retaggio elettromagnetico figlio dall'incredibile lavoro di Maxwell potrebbe indurre a pensare che la funzione d'onda debba soddisfare, così come i campi elettrico e magnetico nel vuoto, l'equazione delle onde per eccellenza, quella di \emph{d'Alembert}. Si presti attenzione al fatto che l'intuizione non è assolutamente banale: il generico elettrone in moto descritto dalla funzione $\psi(\bm{r},t)$ analizzata produce infatti un campo elettromagnetico; è necessario pertanto capire se le equazioni che governano suddetti campi funzionano anche per il corpuscolo che li genera. 

Impostiamo quindi una relazione \emph{simil}  D'Alembert:
\begin{equation}
    \nabla^2 \psi = \frac{1}{\sigma^2} \pdv[2]{\psi}{t} \quad,
    \label{eq:d'ale}
\end{equation}
e cerchiamo di verificare se, sostituita la $\psi(\bm{r,t})$ al suo interno risulta un'uguaglianza \emph{vera}. Il parametro $\sigma$ corrisponde, in riferimento alla formulazione standard dell'equazione,  alla velocità di propagazione della radiazione mentre $\nabla^2$ è l'operatore di Laplace. Sfruttando la \eqref{eq:func_onda} e svolgendo le derivate otteniamo per la parte spaziale:
\[
\begin{aligned}
    \nabla^2 \psi = \nabla \cdot (\nabla\psi) = \nabla \cdot (i\bm{k}\psi) &= i\bm{k}\nabla\psi + \overbrace{\psi\nabla \cdot (i\bm{k})}^{0} \\ &\overset{(1)}{=} (i\bm{k})\cdot(i\bm{k})\psi = -|\bm{k}|^2 \psi = -k^2 \psi\\
\end{aligned}
\]
dove in $(1)$ abbiamo usato l'uniformità spaziale del vettore d'onda.
Per quanto riguarda la parte temporale abbiamo:
\[
    \pdv[2]{\psi}{t} = \pdv{}{t} \, \bigg(\pdv{\psi}{t}\bigg) = \pdv{}{t}\,(-i\omega\psi) = -\omega^2\psi \quad.
\]
Prima di procedere, è interessante analizzare il significato fisico di $k^2$. Vale infatti:
\[
    k^2 = \bm{k} \cdot \bm{k} = \left(\frac{2\pi}{\lambda} \bm{\hat{u}}_r\right) \cdot \left(\frac{2\pi}{\lambda} \bm{\hat{u}}_r\right) = \left(\frac{2\pi}{\lambda}\right)^2 \bm{u}_r \cdot \bm{u}_r \overset{(2)}{=} \frac{p^2}{\hbar^2}\quad,
\]
dove in $(2)$ abbiamo sfruttato l'Ipotesi di De Broglie e l'unitarietà del prodotto scalare di un versore con se stesso. Abbiamo quindi $p = k\hbar$. 

Siamo a questo punto pronti a sostituire le relazioni trovate nell'equazione di partenza:
\[
            -k^2\psi = -\frac{1}{\sigma^2} \omega^2 \psi \implies (\omega^2-\sigma^2 k^2)\psi = 0\quad.
\]
Perché l'uguaglianza valga per qualsiasi scelta della funzione d'onda serve che il termine fra parentesi sia identicamente nullo, da cui:
\[
    \omega = \sigma k \quad.
\]
Quanto appena ottenuto prende il nome di \emph{relazione di dispersione} e permette di stabilire come la lunghezza d'onda della radiazione si rapporta alla pulsazione. Facendo affidamento a De Broglie e a Einstein è ancora possibile espandere l'equazione e ottenere una relazione fra energia e impulso\footnote{Il termine \emph{impulso} è deliberatamente usato per indicare la quantità di moto.}:
\[
    \frac{E}{\hbar} = {\sigma} \frac{p}{\hbar} \implies
    E = {\sigma} p
    \quad .
\]
L'uguaglianza trovata è purtroppo fisicamente incoerente: per una particella libera e \emph{massiva} (come l'elettrone), infatti, l'energia è una funzione quadratica della quantità di moto $E \propto p^2$. La diretta proporzionalità appena ricavata funziona solo nel caso delle radiazioni: l'energia del fotone, infatti, è data da $E = cp$, con $c$ la velocità della luce.

È così necessario cambiare approccio: quello classico-elettromagnetico non risulta più in grado di fornire una corretta descrizione della realtà osservabile.

\subsection{L'ansatz di Schrödinger}
Durante una conferenza a Zurigo, venne fatto notare a Schrödinger l'esigenza di trovare un'equazione d'onda capace di descrivere in modo rigoroso l'evoluzione spazio temporale delle sìdette \emph{onde di materia} (gli elettroni analizzati), coerentemente con quanto postulato da de Broglie. L'aneddoto storico riporta che fu il fisico Peter Debye, al termine di un seminario sulle onde di fase tenuto dallo stesso Schrödinger alla fine del 1925, a obiettare che \guillemetleft\emph{Se si ha a che fare con delle onde, bisogna avere a disposizione un'equazione d'onda}\guillemetright.

Raccogliendo questa sfida, Schrödinger non "dimostrò" la sua equazione partendo da principi primi della meccanica classica, ma procedette per via euristica formulando, per l'appunto, un \textit{ansatz}. La sua intuizione fu quella di riprendere e ampliare la profonda analogia formale, già scoperta da William Rowan Hamilton nel secolo precedente, tra la meccanica e l'ottica geometrica. Così come l'ottica geometrica (basata sul concetto di raggio luminoso) rappresenta un limite approssimato dell'ottica ondulatoria, valido quando la lunghezza d'onda della luce è trascurabile rispetto alle dimensioni caratteristiche del sistema ($\lambda \to 0$), Schrödinger ipotizzò che la meccanica classica non fosse altro che un limite approssimato di una più generale e fondamentale \textbf{meccanica ondulatoria}. In questo quadro concettuale, l'austriaco postulò la seguente relazione:
\begin{equation}
    \nabla^2\psi = i\sigma\pdv{\psi}{t}  \quad.
    \label{eq:ansatz_schrodi}
\end{equation}
Per verificarne però la bontà è necessario verificare che l'equazione rispetti la dipendenza $E \propto p^2$. Calcolando le derivate e sostituendo otteniamo:
\[
        (ik)^2\psi - i\sigma(-i\omega)\psi = -(k^2 + \sigma\omega)\psi = 0 \implies k^2 = -\sigma \omega \quad. 
\]
Possiamo così ricavare la relazione fra energia e quantità di moto:
\[
    \frac{p^2}{\hbar^2} = - \sigma \frac{E}{\hbar} \implies E = -\frac{p^2}{\sigma \hbar}
\]
ed ecco la tanto agognata dipendenza quadratica dalla quantità di moto. Serve in ultima istanza scegliere $\sigma$ per riprodurre la usuale relazione per l'energia di una particella libera $E = p^2/2m$.
\[
    E = \frac{p^2}{2m} = -\frac{p^2}{\sigma \hbar} \implies \sigma = -\frac{2m}{\hbar}
\quad.
\]
Sostituendo il $\sigma$ appena trovato nella \eqref{eq:ansatz_schrodi}, osserviamo l'ipotesi di Schrödinger assumere una forma sempre più concreta:
\begin{equation}
        -\frac{2mi}{\hbar} \pdv{\psi}{t} = \nabla^2\psi \implies
        {i\hbar \pdv{\psi}{t} = -\frac{\hbar^2}{2m}\nabla^2\psi} \quad.  
    \label{eq:schrodi_cin}
\end{equation}

\begin{approfondimento}[title={Un ponte fra i due mondi: la relatività ristretta}]
Vale la pena notare come la relazione relativistica di Einstein per l'energia totale di una particella libera:
\[
    E = \sqrt{p^2c^2 + m^2c^4} \quad
\]
rappresenti il vero "ponte" teorico tra i due regimi analizzati in precedenza. L'equazione di d'Alembert e l'ansatz di Schrödinger descrivono infatti i due limiti asintotici di questa legge fondamentale. 

Nel caso della luce, costituita da fotoni con $m = 0$, la relazione si riduce all'esatta proporzionalità lineare $E = cp$ discussa in precedenza, coerentemente con la relazione di dispersione che deriva dall'equazione delle onde di d'Alembert. 

Al contrario, per particelle massive in moto a velocità non relativistiche ($p \ll mc$), è possibile espandere in serie di Taylor il termine sotto radice:
\[
    E = \sqrt{p^2c^2 + m^2c^4} = mc^2\sqrt{1 + \frac{p^2}{m^2c^2}} \approx mc^2\left( 1 + \frac{p^2}{2m^2c^2} \right) = mc^2 + \frac{p^2}{2m} \quad.
\]
Notiamo come, a meno della costante additiva $mc^2$ (proporzionale l'energia di riposo della particella), l'espressione coincida con l'energia cinetica classica $E_c = p^2/2m$, l'unica a soddisfare l'equazione di Schrödinger per una particella libera. 

Per descrivere quantisticamente l'energia relativistica completa senza ricorrere ad approssimazioni, sarà necessario formulare nuove equazioni, come l'equazione di Klein-Gordon e di Dirac.
\end{approfondimento}

\section{Un cambio di paradigma}
L'equazione \eqref{eq:schrodi_cin} ci permette quindi di descrivere con successo quelle particelle che esibiscono comportamenti coerenti con
un fenomeno ondulatorio. È istruttivo allora tentare di comprendere se le nuove grandezze coinvolte nell'equazione si leghino in qualche modo alle già conosciute quantità tipiche della fisica classica. Iniziamo dall'energia.

\subsection{L'energia totale}
Sviluppando il membro sinistro della \eqref{eq:schrodi_cin} si ottiene:
\begin{equation}
        i\hbar \pdv{\psi}{t} = i\hbar (-i\omega)\psi = \hbar \omega \psi = E\psi \quad.
        \label{eq:op_E}
\end{equation}
L'effetto complessivo della derivazione della funzione d'onda nel tempo, congiunto alla moltiplicazione per il fattore $i\hbar$ "produce" quindi l'energia totale della particella. Forti di una profonda analogia con l'algebra lineare e il concetto di \emph{problema agli autovalori}, possiamo allora interpretare l'azione del termine $i\hbar \partial_t$ sulla funzione d'onda come quella di un \emph{operatore} ($\implies$ matrice, applicazione lineare) che estrae da $\psi(\bm{r},t)$ ($\implies$ vettore) l'energia della particella ($\implies$ autovalore).

\subsection{L'energia cinetica}
Analizziamo ora il membro destro della \eqref{eq:schrodi_cin}:
\begin{equation*}
    -\frac{\hbar^2}{2m}\nabla^2 \psi = -\frac{\hbar^2}{2m} (i\bm{k}) \cdot (i\bm{k}) \psi = -\frac{\hbar^2}{2m} (-k^2) \psi = \frac{k^2\hbar^2}{2m} \psi \overset{(3)}{=} \frac{p^2}{2m} \psi \quad,
\end{equation*}
dove in $(3)$ abbiamo sfruttato l'ipotesi di de Broglie $p = \hbar k$. È immediato notare come il termine moltiplicativo della funzione d'onda sia l'energia totale della particella che, nel caso di un sistema libero (come quello trattato nella discussione dell'ansatz di Schrödinger), corrisponde all'energia cinetica. Infatti, forti di quanto appena ricavato:
\[
        i\hbar\pdv{\psi}{t} = E\psi = -\frac{\hbar^2}{2m}\nabla^2\psi = \frac{p^2}{2m}\psi = K\psi \implies
        E = K \quad.
\]

Il pattern inizia a questo punto ad emergere: abbiamo identificato due \emph{operatori} che, applicati alla funzione d'onda generano rispettivamente la sua energia totale e la sua energia cinetica. Definiamo allora \textbf{operatore energia cinetica} $\hat{K}$:
\begin{equation}
    {\h{K} = -\frac{\hbar^2}{2m}\nabla^2} \quad,
    \label{eq:operatore_K}
\end{equation}
caratterizzato dall'equazione agli autovalori\footnote{Nel prosieguo della trattazione sarà sempre più chiaro perché queste siano effettivamente equazioni agli autovalori. La capacità di astrazione è in ogni caso necessaria perché i vari operatori che abbiamo definito e che definiremo non saranno necessariamente ascrivibili al mondo delle \emph{matrici}. Una forte reminiscenza all'algebra lineare, in ogni caso, sarà particolarmente evidente durante l'analisi dell'approccio di Dirac alla meccanica quantisitica.}:
\[
    \h{K}\psi = \frac{p^2}{2m} \psi \quad.
\]
Da notare il fatto che l'uguaglianza vale \emph{solo perché} si sta analizzando una partcella descritta da un'onda piana a cui è associato un momento \emph{ben definito}. L'equazione perde di significato con stati caratterizzati da una quantità di moto non definibile. 

\subsection{L'operatore quantità di moto}
Anche se non direttamente coinvolta nell'equazione, la quantità di moto è una quantità fondamentale e pertanto ne proponiamo l'indagine in termini quantistici. L'idea è quindi quella di cercare di esprimere l'\textbf{operatore quantità di moto} $\hat{\bm{p}}$. Dobbiamo cambiare lievemante approccio però rispetto a quanto fatto con l'operatore energia cinetica, dal momento che $\hat{\bm{p}}$ \emph{deve} essere un operatore \emph{vettoriale}: la sua azione su $\psi$ deve infatti produrre la sua controparte classica, non uno scalare (come $K$) bensì il vettore quantità di moto $\bm{p}$. 

Dalla meccanica classica sappiamo che $K = p^2/2m = \bm{p}\cdot\bm{p}/{2m}$. Per analogia, in meccanica quantistica dovremo avere:
\begin{equation*}
    \h{K} = \frac{\hat{\bm{p}} \cdot \hat{\bm{p}}}{2m} \implies \hat{\bm{p}} \cdot \hat{\bm{p}} = -\hbar^2\nabla^2 \quad.
\end{equation*}
Ricordando che il Laplaciano è il prodotto scalare del gradiente con se stesso ($\nabla^2 = \nabla \cdot \nabla$), la relazione precedente suggerisce per $\hat{\bm{p}}$ una forma vettoriale proporzionale al gradiente:
\begin{equation*}
    \hat{\bm{p}} = \beta \nabla \quad,
\end{equation*}
con $\beta$ costante da determinare. Sostituendo questa ipotesi nella definizione dell'operatore energia cinetica abbiamo:
\begin{equation*}
    (\beta\nabla) \cdot (\beta\nabla) = \beta^2\nabla^2 = -\hbar^2\nabla^2 \implies \beta^2 = -\hbar^2 \implies \beta = \pm i\hbar \quad.
\end{equation*}
Siamo di fronte a un'ambiguità matematica: sia scegliendo il segno più che il segno meno otteniamo il corretto operatore energia cinetica. Per risolvere questo dubbio, dobbiamo imporre che l'operatore $\hat{\bm{p}}$, agendo sulla generica funzione d'onda $\psi(\bm{r},t) = A e^{i(\bm{k}\cdot\bm{r} - \omega t)}$, "estragga" il corretto valore della quantità di moto vettoriale, coerente con il verso di propagazione dell'onda, ovvero $\bm{p} = +\hbar\bm{k}$. Testiamo la due ipotesi, $\hat{\bm{p}} = \pm i\hbar\nabla$:
\[
\begin{cases}
    (+) \quad \hat{\bm{p}} = i\hbar\nabla \implies (i\hbar\nabla) \psi = i\hbar (i\bm{k}) \psi = -\hbar\bm{k}\psi = -\bm{p}\psi \quad  \ensuremath{\times}\\
    (-) \quad \hat{\bm{p}} = -i\hbar\nabla \implies (-i\hbar\nabla) \psi = -i\hbar (i\bm{k}) \psi = \hbar\bm{k}\psi = +\bm{p}\psi \quad \ensuremath{\checkmark} \quad .
\end{cases}
\]
L'ambiguità è a questo punto risolta e deduciamo in modo inequivocabile la forma dell'operatore quantità di moto:
\begin{equation}
    {\hat{\bm{p}} = -i\hbar\nabla} \quad.
    \label{eq:operatore_p}
\end{equation}

È ancora necessaria una postilla per poter dire conclusa l'analisi di $\hb{p}$. Dobbiamo infatti ancora chiarire il significato dell'espressione \emph{operatore vettoriale}. Si abbia ben chiaro che la sua azione su $\psi$ è quella di produrre $\bm{p}$, dunque una tripla di scalari $(p_1, p_2, p_3)$ nel generico spazio tridimensionale. In questo senso, intrepreteremo $\hat{\bm{p}}$ come un operatore "formato da tre componenti scalari" (pur non essendo concretamente un vettore), ciascuna delle quali agente lungo una direzione indipendente dello spazio in cui vive $\psi$:
\begin{equation}
    \hat{\bm{p}} = 
    \begin{pmatrix}
        \hat{p}_1 \\
        \hat{p}_2 \\
        \hat{p}_3 \\
    \end{pmatrix}
    = 
    \begin{pmatrix}
        -i\hbar \partial_{x_1} \\
        -i\hbar \partial_{x_2} \\
        -i\hbar \partial_{x_3} \\
    \end{pmatrix}
    \quad.
    \label{eq:op_p}
\end{equation}   

\section{L'equazione di Schrödinger}
Quanto fatto sino ad ora ci conferisce le basi necessarie per formulare nella sua forma più completa la fantomatica equazione di Schrödinger (S.E., dal'inglese \emph{Schrödinger's equation}). Si noti innanzitutto che la \eqref{eq:schrodi_cin} è stata ricavata imponendo che la funzione d'onda $\psi$ descrivesse un sistema libero, perciò non immerso in campi di forze. Assumendo però di trattare anche \emph{sistemi monogenici}\footnote{Ovvero sistemi in cui si può applicare il principio dell'\emph{azione stazionaria} di Hamilton, in cui le forze possono essere espresse tramite adeguati potenziali scalari.}, si deve poter tenere in considerazione di un generico campo scalare $V(\bm{r},t)$ che descrive l'energia potenziale di interazione. L'energia meccanica totale vale
\[
    E = K + V = H \quad,
\]
con $H$ l'\textbf{hamiltoniana} del sistema. Si presti estrema attenzione a quanto appena discusso: l'hamiltoniana di un sistema, così come definita entro la trattazione della meccanica analitica, corrisponde all'energia del sistema \emph{solo sotto precise condizioni}. Serve infatti che il potenziale che descrive l'interazione del sistema come l'esterno sia indipendente dalle velocità e che qualsiasi eventuale vincolo esistente sia \emph{scleronomo}, dunque indipendente dal tempo.

Per proporre un'adeguata riformulazione della \eqref{eq:ansatz_schrodi}, dobbiamo quindi definire un operatore che estragga energia meccanica del sistema, intesa come somma di energia cinetica e potenziale. Nell'ansatz iniziale avevamo banalmente $H = K$, da cui:
\[
    \h{H} = \h{K} = -\frac{\hbar^2}{2m}\nabla^2 \quad.
\]
Generalizziamo così per induzione al caso generale $H = K + V$, introducendo a tal proposito l'\textbf{operatore energia potenziale} $\h{V}$, tale che $\h{V}\psi = V(\bm{r}, t)\psi$. 
Possiamo adesso definire, in piena analogia con la meccanica classica, l'\textbf{operatore hamiltoniano} $\h{H}$ come:
\[
    \h{H} = \h{K} + \h{V} = -\frac{\hbar^2}{2m}\nabla^2 + V(\bm{r},t) \quad.
\]
Siamo in direttura d'arrivo. Notando che sia $\h{H}$ sia $i\hbar \partial_t$ producono l'energia totale del sistema se applicati alla funzione d'onda, abbiamo l'\textbf{equazione di Schrödinger}:
\[
    {i\hbar \pdv{\psi(\bm{r},t)}{t} = \h{H}\psi(\bm{r},t)} \quad,
\]
esprimibile equivalentemente riscrivendo gli operatori:
\[
    i\hbar \pdv{\psi(\bm{r},t)}{t} = -\frac{\hbar^2}{2m}\nabla^2\psi(\bm{r},t) + V(\bm{r},t)\psi(\bm{r},t) \quad.
\]

\begin{concettochiave}{Equazione di Schrödinger}
    \noindent
    L'evoluzione temporale di un generico corpuscolo di massa $m$ confinato in un volume $\Omega$ e immerso in un potenziale scalare $V(\bm{r},t)$ è dettata dall'\textbf{equazione di Schrödinger}:
    \begin{equation}
        {i\hbar \pdv{\psi(\bm{r},t)}{t} = \h{H}\psi(\bm{r},t)} \quad,
        \label{eq:schrodinger_completa}
    \end{equation}
    un'equazione differenziale lineare alle derivate parziali che richiede soluzioni almeno $\mathcal{C}^2(\Omega)$ nello spazio e $\mathcal{C}^1(\Omega)$ nel tempo.
\end{concettochiave}

\subsubsection{L'Hamiltoniana come energia totale del sistema}
Come accennato in precedenza, i dettami della meccanica analitica assicurano che l'uguaglianza tra la funzione Hamiltoniana ${H}$ e l'energia meccanica totale $E = K + V$ non è un'identità a priori, ma è garantita esclusivamente se il sistema è soggetto a vincoli {scleronomi} e potenziali velocità-indipendenti. Solo in questo caso, infatti, l'energia cinetica risulta essere una funzione quadratica omogenea delle velocità generalizzate e, per il teorema di Eulero sulle funzioni omogenee, si ottiene ${H} = E$. Se i vincoli fossero \emph{reonomi} (variabili nel tempo), si avrebbe ${H} \neq E$. 

È dunque lecito chiedersi: cosa garantisce in Meccanica Quantistica l'assenza di vincoli reonomi, tale da giustificare l'assunzione ${H} = {E}$ per postulare l'equazione di Schrödinger?

La giustificazione risiede nel dominio di validità fondamentale della teoria quantistica. A livello subatomico o microscopico, i vincoli macroscopici della meccanica classica (come binari, cerniere, o superfici in movimento) \textbf{non esistono}. I costituenti elementari della materia si muovono liberamente nello spazio tridimensionale $\mathbb{R}^3$, soggetti unicamente alle interazioni fondamentali (forze elettromagnetiche, nucleari, ecc.). 
In assenza di superfici vincolari, le coordinate generalizzate del sistema coincidono banalmente con le coordinate cartesiane delle singole particelle. L'energia cinetica assume quindi \emph{sempre} ed esclusivamente la forma di una forma quadratica omogenea:
\begin{equation*}
    K = \sum_{i=1}^N \frac{\bm{p}_i^2}{2m_i} \quad.
\end{equation*}
Essendo strutturalmente assente qualsiasi termine lineare o di grado zero nelle quantità di moto, la coincidenza formale tra l'Hamiltoniana e l'energia meccanica (${H} = E$) è intrinsecamente e sempre garantita a livello fondamentale.

Rimane però da chiarire il caso in cui il potenziale dipenda esplicitamente dal tempo, $V = V(\bm{r}, t)$ (ad esempio, un elettrone investito da un'onda elettromagnetica generata da un laser esterno). In questo caso:
\begin{equation*}
    \h{H}(t) = \h{K} + \h{V}(\bm{r}, t) \quad.
\end{equation*}
L'operatore $\h{H}(t)$ \emph{continua} a rappresentare l'osservabile "energia totale del sistema" istante per istante. Il fatto che l'Hamiltoniana dipenda esplicitamente dal tempo ($\partial_t \h{H} \neq 0$) implica semplicemente che \textbf{l'energia del sistema quantistico non si conserva}, poiché esso sta scambiando lavoro con l'ambiente esterno classico che genera il campo. L'equazione di Schrödinger $i\hbar\partial_t \psi = \hat{H}(t)\psi$ resta formalmente valida come generatrice della dinamica, ma lo stato del sistema sarà costretto a compiere transizioni tra livelli energetici diversi proprio a causa della forzante esterna tempo-dipendente.


\subsection{Il principio di sovrapposizione}
La \emph{linearità} della S.E., permette, almeno in linea teorica, di avere a disposizione \emph{infinite} soluzioni per l'equazioni. Dette $\psi_1(\bm{r},t), \dots, \psi_n(\bm{r},t)$ un numero $n$ di funzioni d'onda che soddisfano la S.E.,
\[
    \begin{aligned}
        i\hbar \partial_t\psi_1(\bm{r},t) &= \h{H}\psi_1(\bm{r},t)\\
        &\vdots  \\
        i\hbar \partial_t\psi_n(\bm{r},t) &= \h{H}\psi_n(\bm{r},t)
    \end{aligned}
    \quad,
\]
qualsiasi \textbf{combinazione lineare} $\Psi(\bm{r},t)$ delle $\psi_n(\bm{r},t)$ 
\begin{equation}
    \Psi(\bm{r},t) = \sum_{j=1}^n c_j\psi_j(\bm{r},t) \quad,
\end{equation}
con $c_j \in \mathbb{C}$ costanti dette \emph{ampiezze}, rimane soluzione della S.E. Infatti:
\begin{equation*}
\begin{aligned}
    i\hbar\pdv{\Psi(\bm{r},t)}{t} &= i\hbar \sum_{j=1}^n c_j\psi_j(\bm{r},t) \overset{(4)}{=} \sum_{j=1}^n c_j \left(i\hbar \pdv{\psi_j(\bm{r},t)}{t} \right)\\ &= \sum_{j=1}^n c_j \h{H}\psi_j(\bm{r},t) = \h{H}\sum_{j=1}^n c_j\psi_j(\bm{r},t) \overset{(5)}{=} \h{H}\Psi(\bm{r},t) \quad.
\end{aligned}
\end{equation*}
In $(4)$ e $(5)$ abbiamo sfruttato la linearità degli operatori differenziali $\partial_t$ e $\h{H}$ (costituito da una parte differenziale, $\h{K}$, e una moltiplicativa, $\h{V}$, quindi evidentemente lineare rispetto alla somma).

\section{L'interpretazione di Copenhagen}
La meccanica classica, con tutte le sue derive Newtoniana, Lagrangiana, Hamiltoniana, ecc., è una teoria \textbf{deterministica}. La conoscenza di posizione e impulso di una particella istante per istante durante il suo moto è \emph{sufficiente} per determinare qualsiasi altra informazione riguardo la stessa, compresa la sua traiettoria nello spazio (eventualmente delle fasi). Dobbiamo ora capire se questa riformulazione quantistica della fisica conserva queste proprietà prettamente deterministiche oppure se, anche in questo contesto, è necessario un cambio di approccio. 

\subsection{L'idea di Born e la densità di probabilità}
«\emph{Cosa rappresenta, concretamente, la funzione d'onda?}». Questa domanda ha aleggiato a lungo tra i pionieri della meccanica quantistica prima di trovare una risposta esplicativa. Se già risulta poco intuitivo accettare che un corpuscolo materiale sia descritto da un'onda $\psi(\bm{r},t)$, comprendere il significato fisico diretto di tale ampiezza appare ancora più arduo. Fu \textbf{Max Born} a intuire la svolta, introducendo l'\textbf{interpretazione probabilistica} che permette di estrarre informazione fisica osservabile a partire da $\psi$.

Per comprendere la logica di Born, risulta particolarmente illuminante un esperimento concettuale basato sulla classica \emph{esperienza della doppia fenditura di Young} per la radiazione elettromagnetica. 

\subsubsection{L'esperienza della doppia fenditura}
Si consideri una sorgente che emette radiazione verso una parete dotata di due fenditure. Sullo schermo di rilevazione posto a valle, l'intensità luminosa $I(P)$ registrata in un generico punto $P$ è nota essere proporzionale al modulo quadro del campo elettrico:
\begin{equation*}
    I(P) \propto |\bm{E}(P)|^2
\end{equation*}
Volendo reinterpretare il fenomeno in chiave quantistica (ovvero corpuscolare), l'intensità $I(P)$ misura la frazione di fotoni che incide in $P$ rispetto al totale di quelli emessi dalla sorgente, in ossequio all'interpetazione della radiazione del corpo nero fornita da Planck e dell'effetto fotoelettrico di Einstein. Di conseguenza, il rapporto tra l'intensità locale e quella totale definisce direttamente la \textbf{probabilità che un singolo fotone arrivi nel punto $P$}.

Estendendo questa intuizione alle onde di materia, Born ipotizzò che la funzione d'onda $\psi$ svolga per le particelle un ruolo del tutto analogo a quello del campo elettrico per i fotoni. Pertanto, il modulo quadro della funzione d'onda non descrive una massa o una carica "diluita" nello spazio, bensì una \textit{densità di probabilità}.    

\subsubsection{Definizione e proprietà della densità di probabilità}
Si definisce \textbf{densità di probabilità} $\rho(\bm{r},t)$ la probabilità per unità di volume di rinvenire una particella descritta dalla funzione d'onda $\psi(\bm{r},t)$ nell'infinitesimo $\dd^3{\bm{r}}$ all'istante $t$. In formule:
\begin{equation}
    \rho(\bm{r}, t) = |\psi(\bm{r}, t)|^2 = \psi(\bm{r}, t)^*\psi(\bm{r}, t) = \dv{\mathcal{P}}{^3\bm{r}} \quad,
    \label{eq:prob_dens}
\end{equation}
con $\psi^*$ il complesso coniugato di $\psi$ e con $\mathcal{P}$ la probabilità. Di conseguenza, se si confina il corpuscolo in un volume $\Omega$, la probabilità che si trovi entro $\Omega$ deve essere unitaria:
\begin{equation}
        \mathcal{P}(\Omega) = \int_\Omega |\psi(\bm{r}, t)|^2 \dd^3{\bm{r}} = 1
        \label{eq:unit_psi}
\end{equation}
e, preso $\Omega^{'} \subseteq \Omega$, deve essere:
\[
    \mathcal{P}\big(\Omega^{'}\big) = \int_{\Omega^{'}} |\psi(\bm{r}, t)|^2 \dd^3{\bm{r}} \leq 1 \quad.
\]
Quanto appena postulato nasconde non poche insidie. Infatti, ha senso aspettarsi che il confinamento di una particella entro un determinato volume (si immagini una scatola) sia \emph{definitivo}: il corpuscolo non può scappare dal suo contenitore. Come si può però costruire una teoria probabilistica in cui la probabilità totale deve essere costante e unitaria ma la sua densità è tempo-dipendente? In altre parole, \guillemetleft\emph{Perchè, se la densità di probabilità $\rho(\bm{r},t)$ dipende esplicitamente dal tempo, la probabilità totale associata è una costante?}\guillemetright.

\subsection{L'evoluzione temporale della densità di probabilità e l'equazione di continuità}
Partiamo da una delle poche certezze sin'ora a nostra disposizione, la \eqref{eq:prob_dens}. Presa una particella descritta dalla funzione d'onda $\psi(\bm{r},t)$, studiare come la probabilità totale ad essa associata si evolva nel tempo implica valutare:
\begin{equation}
    \dv{\mathcal{P}}{t} = \dv{}{t} \int_\Omega \rho(\bm{r}, t) \dd^3{\bm{r}} \overset{(6)}{=} \int_\Omega \pdv{\rho(\bm{r}, t)}{t} \dd^3{\bm{r}} \quad,
    \label{eq:studio_dens_prob}
\end{equation}
dove in $(6)$ si è sfruttata la regola di Leibniz per passare l'operatore di derivata sotto il segno d'integrale\footnote{Si noti che l'operazione d'integrale in questione \emph{satura} qualsiasi eventuale dipendenza spaziale della funzione integranda (nel caso in esame la densità di probabilità $\rho$). La funzione che sopravvive all'integrazione deve dipendere esclusivamente dal tempo e il segno di derivata totale è così giustificato. Nel momento in cui si fanno commutare gli operatori, azione possibile solo perché agiscono su variabili indipendenti, serve tenere presente che l'integranda può avare dipendenze sia esplicite sia implicite dal tempo. Per ovviare a questo cavillo si trasforma $\text{d}_t \mapsto \partial_t$.}. È quindi necessario caratterizzare la quantità integranda per carpire le proprietà di $\text{d}_t \mathcal{P}$. 

Invocando la definizione di densità di probabilità e muniti di pazienza si comincia la derivazione:
\[
    \begin{aligned}
        \pdv{\rho(\bm{r}, t)}{t} &= \pdv{|\psi(\bm{r}, t)|^2}{t} = \pdv{}{t}\, (\psi(\bm{r}, t)^*\psi(\bm{r}, t)) \\
        &= \psi(\bm{r}, t)^*\pdv{\psi(\bm{r}, t)}{t} + \psi(\bm{r}, t)\pdv{\psi(\bm{r}, t)^*}{t} \quad.
    \end{aligned}
\]
L'equazione di Schrödinger ci permette di riscrivere i termini contenenti le derivate parziali rispetto al tempo delle funzioni d'onda tramite gli operatori energia cinetica e potenziale. Nel primo caso abbiamo:
\[
    \pdv{\psi(\bm{r}, t)}{t} = \frac{1}{i\hbar} \left( -\frac{\hbar^2}{2m}\nabla^2 \psi(\bm{r}, t) + V(\bm{r}, t)\psi(\bm{r}, t) \right)
\]
e, prendendo il coniugato ad ambo i membri ricaviamo la seconda quantità utile:
\[
    \pdv{\psi(\bm{r}, t)^*}{t} = -\frac{1}{i\hbar} \left( -\frac{\hbar^2}{2m}\nabla^2 \psi(\bm{r}, t)^* + V(\bm{r}, t)\psi(\bm{r}, t)^* \right) \quad.
\]
Sostituendo otteniamo:
\[
    \begin{aligned}
        \pdv{\rho(\bm{r}, t)}{t} = \frac{\psi(\bm{r}, t)^*}{i\hbar}\bigg( &-\frac{\hbar^2}{2m}\nabla^2 \psi(\bm{r}, t) + V(\bm{r}, t)\psi(\bm{r}, t) \bigg) -\\ &- \frac{\psi(\bm{r}, t)}{i\hbar}\left( -\frac{\hbar^2}{2m}\nabla^2 \psi(\bm{r}, t)^* + V(\bm{r}, t)\psi(\bm{r}, t)^* \right) \quad.
    \end{aligned}
\]  
I termini contenenti i potenziali si elidono e, con le dovute semplificaizoni algebriche, si ha:
\begin{equation}
        \pdv{\rho(\bm{r}, t)}{t} = -\frac{\hbar}{2mi}\left(\psi(\bm{r}, t)^*\nabla^2\psi(\bm{r}, t) - \psi(\bm{r}, t)\nabla^2\psi(\bm{r}, t)^*\right) \quad.
        \label{eq:step_eq_continuità}
\end{equation}
D'ora in poi ometteremo gli argomenti delle funzioni analizzate per alleggerire la notazione. 
Possiamo riscrivere il fattore fra parentesi in maniera più conveniente. Per tentare di intuire quale sia la strada migliore, cerchiamo di capire \emph{davvero} cosa stiamo calcolando. 

Si immagini che la particella di cui stiamo cercando di tracciare il moto sia una carica $q$ in moto in un conduttore. Le leggi di Maxwell, tramite cui analizziamo il suo comportamento, non si arrogano il diritto di trattare la \emph{traiettoria} della singola carica, bensì ci danno informazioni interessanti riguardo il rapporto incrementale fra cariche e volume, ovvero riguardo la densità di carica 
\[
\rho_{(q)} = \frac{q}{\tau} \,, \quad q = nu
\]
dove con $\tau$ indichiamo il volume del conduttore, con $u = \SI{1}{\coulomb}$ la carica unitaria e con $n$ il numero di cariche elementari che compongono $q$\footnote{Per rendere più chiaro quanto appena descritto, si consideri il seguente esempio: sia $q = \SI{18}{\coulomb}$. Allora possiamo immaginare che $q$ non sia altro che l'agglomerato di $n = 18$ cariche elementari $u$.}. \guillemetleft\emph{Ma cosa rappresenta questa densità?}\guillemetright. Scriviamo:
\[
    \rho_{(q)} = \frac{n}{\tau} \, u\quad,
\] 
dove il termine $n/\tau$ non è altro che la probabilità di trovare una carica elementare nell'unità di volume, quasi come la densità di probabilità quantistica. Abbiamo così a disposizione una fortissima analogia fra la nostra $\rho(\bm{r},t)$ e la densità di carica classica, la cui evoluzione temporale è determinata dall'\emph{equazione di continuità}:
\[
    \pdv{\rho_{(q)}}{t} + \nabla \cdot \bm{J}_{(q)} = 0 \implies \pdv{\rho_{(q)}}{t} = -\nabla \cdot \bm{J}_{(q)} \quad.
\] 
Possiamo allora lavorare affinché la \eqref{eq:step_eq_continuità} sia quanto più possibile affine alla controparte classica: a tal fine si tenta di fare emergere l'operatore di divergenza "a fattor comune", così da avere $\partial_t \rho = \nabla \cdot (\text{qualcosa})$. L'operazione tuttavia non è assolutamente ovvia trattandosi di un operatore differenziale. Iniziamo allora con il formalizzare l'intuizione:
\[
    \psi^*\nabla^2\psi - \psi\nabla^2\psi^* \overset{?}{=} \nabla \cdot \big( \psi^*\nabla\psi - \psi\nabla\psi^* \big) \quad.
\]
Calcoliamo:
\[
\begin{aligned}
    \nabla \cdot \big( \psi^*\nabla\psi - \psi\nabla\psi^* \big) &= \nabla\psi\cdot\nabla\psi^* + \psi^*\nabla^2\psi - \nabla\psi^* \cdotp\nabla\psi - \psi\nabla^2\psi^* \\ &= \psi^*\nabla^2\psi - \psi\nabla^2\psi^* \quad.
\end{aligned}
\]
Appurata la fattibilità matematica del raccoglimento, sostituiamo e otteniamo la forma finale dell'\textbf{equazione di continuità} nel caso quantistico:
\[
    \pdv{\rho}{t} + \nabla \cdot \left[\frac{\hbar}{2mi}\big( \psi^*\nabla\psi - \psi\nabla\psi^* \big) \right] = 0 \quad,
\]
con il termine fra parentesi quadre che, in piena analogia con il caso elettromagnetico, prende il nome di \textit{densità di corrente di probabilità}:
\[
    \bm{J} = \frac{\hbar}{2mi}\big( \psi^*\nabla\psi - \psi\nabla\psi^* \big) \quad.
\]
L'equazione di continuità assume così la sua rinomata forma compatta:
\[
    \pdv{\rho(\bm{r}, t)}{t} + \nabla \cdot \bm{J}(\bm{r},t) = 0 \quad.
\]

\begin{concettochiave}{L'equazione di continuità e la densità di corrente di probabilità}
    \noindent
    In piena analogia con il caso elettromagnetico, la densità di probabilità quantistica evolve nel tempo in ossequio all'\textbf{equazione di continuità}:
    \begin{equation}
    \pdv{\rho(\bm{r}, t)}{t} + \nabla \cdot \bm{J}(\bm{r},t) = 0 \quad.
    \label{eq:continuità}
\end{equation}
dove $\bm{J}(\bm{r},t)$ è la \emph{densità di corrente di probabilità}, definita come:
\begin{equation}
    \bm{J} = \frac{\hbar}{2mi}\big( \psi^*\nabla\psi - \psi\nabla\psi^* \big) \quad.
\end{equation}
\end{concettochiave}

\subsubsection{Integrazione dell'equazione di continuità}
Ritorniamo alla questione iniziale, ovvero come evolve nel tempo la probabilità totale. Se integriamo nel generico volume di definizione $\Omega$ in è confinato il corpuscolo descritto dalla funzione d'onda $\psi$, ritroviamo la nozione di derivata della probabilità totale nel tempo:
\[
    \int_\Omega \pdv{\rho}{t} \, \dd^3{\bm{r}} = \dv{\mathcal{P}}{t} = -\int_\Omega \nabla \cdot \bm{J} \, \dd^3{\bm{r}} \quad.
\]
Sfruttando il Teorema della divergenza di Gauss, si riscrive il membro destro dell'equazione e si ottiene la \eqref{eq:continuità} in forma integrale:
\[
    \dv{\mathcal{P}}{t} = -\oint_{\partial \Omega} \bm{J} \cdot \dd{\bm{S}} \quad,
\]
che ha un'interpretazione fisica bellissima: la variazione temporale della probabilità totale $\mathcal{P}$ di trovare la particella all'interno del volume $\Omega$ è pari al flusso, cambiato di segno, della densità di corrente di probabilità $\bm{J}$ attraverso la superficie chiusa $\partial \Omega$ che delimita il volume stesso\footnote{Si ricordi l'analogia precedentemente discussa: la controparte classica dell'equazione è $\dv{q}{t} = -\oint_{\partial \Omega} \bm{J}_q \cdot \dd{\bm{S}}$. Per convenzione, un flusso positivo (uscente dalla superficie) indica cariche che \emph{escono} dal volume, viceversa un flusso negativo (entrante) comporta un aumento del numero di cariche entro il volume. La traduzione quantistica è quindi presto fatta.}. In altre parole, la probabilità quantistica \emph{non si crea e non si distrugge}, ma fluisce da una regione all'altra dello spazio in modo continuo.

È facile allora capire che affinché $\mathcal{P}$ sia costante (e unitaria) nel tempo, e quindi perché l'interpretazione di Copenaghen sia salva dal punto di vista matematico, serve che il flusso della densità di corrente sia nullo sul bordo di $\Omega$. La condizione si concretizza dunque come:
\[
    \oint_{\partial \Omega} \bm{J} \cdot \dd{\bm{S}} = \frac{\hbar}{2mi} \oint_{\partial \Omega} \big(\psi^*\nabla\psi - \psi\nabla\psi^*\big) \cdot \dd{\bm{S}} = 0 \quad.  
\]
Per comprendere in che condizioni questa richiesta è verificata è necessario discutere riguardo alcune caratteristiche della funzione d'onda.

\subsection{Il decadimento all'infinito della funzione d'onda}
Si può determinare, almeno qualitativamente, cosa succede ad una generica $\psi(\bm{r},t)$ quando ci si trova molto distanti dall'origine, ovvero nel caso in cui $|\bm{r}| = r \to \infty$. 

Una prima fondamentale assunzione che faremo è che, nel caso in cui il volume di definizione di una particella coincida con l'intero spazio tridimensionale ($\Omega \equiv \mathbb{R}^3$), la funzione d'onda debba essere: $(1)$ \emph{velocemente decadente a infinito}; $(2)$ \emph{sufficientemente regolare}. La condizione $(1)$ è necessaria per scongiurare casi patologici contraddistinti da probabilità divergenti mentre la $(2)$, benchè non sia un semplice vezzo matematico, è ridondante ma semplifica notevolmente i calcoli. 

È necessario a questo capire quanto rapidamente $\psi(\bm{r},t)$ debba convergere a $0$ in un intorno di infinito perchè il flusso della densità di corrente di probabilià associata sia anch'esso infinitesimo nel limite $r \to \infty$. La forma funzionale cercata per la funzione d'onda, per grandi distanze dall'origine, è del tipo: 
\begin{equation}
        \psi(\bm{r},t) \simeq \frac{f(r,t)}{r^\gamma} 
        \label{eq:psi_r_grande}
\end{equation} 
con $f:\mathbb{R} \to \mathbb{C}$ una tale che $|f(r)| < \infty$. Si noti che argomento di $\psi$ è il vettore $\bm{r}$ mentre della funzione $f$ è semplicemente la distanza $r = |\bm{r}|$. La trattazione si concentra quindi sul determinare $\gamma$ affinché l'ipotesi di Born sia ben posta. 

\subsubsection{Calcolo del flusso}
Si ricava preliminarmente $\bm{J}$ tramite la definizione:
\[
\begin{aligned}
        \bm{J}(\bm{r},t) &= \frac{\hbar}{2mi}\bigg[\underbrace{\frac{f(r,t)^*}{r^\gamma} \left( \frac{\nabla f(r,t)}{r^\gamma} - \frac{\gamma f(r,t)}{r^{\gamma+1}}\bm{u}_r \right)}_a - \, a^* \bigg] \\ &=  \frac{\hbar}{2mi}\left[\left( \frac{\nabla f(r,t)}{r^{2\gamma}} f(r,t)^* - \frac{\gamma|f(r,t)|^2}{r^{2\gamma+1}}\bm{u}_r \right) - a^*\right] \\
        &= \frac{\hbar}{2mi}\biggl[\frac{1}{r^{2\gamma}}\biggl( {f(r,t)^*}\nabla f(r,t) -  \underbrace{\frac{\gamma|f(r,t)|^2}{r}\bm{u}_r}_{\text{reale}}\biggr) - a^*\biggr] \quad.
\end{aligned}
\]
Dal momento che $a^*$ è il complesso coniugato del primo termine, la differenza uccide la parte puramente reale. Resta quindi una elegantissima dipendenza per la densità di corrente: 
\[
    \bm{J}(\bm{r},t) \propto \frac{1}{r^{2\gamma}}\left({f(r,t)^*}\nabla f(r,t) - {f(r,t)}\nabla f(r,t)^*  \right) \quad.
\]
Tornando al flusso, sappiamo che la generica superficie $\Sigma$ cresce come $r^2$. In coordinate sferiche, infatti, l'elemento  infinitesimo di superficie risulta $\dd{\bm{S}} = \bm{u}_r \, r^2 \dd{\Omega}$, con $\dd{\Omega}$ l'elemento infinitesimo di angolo solido. Scriviamo così l'integrale di flusso come:
\begin{equation}
    \oint_{\Sigma = \partial \Omega} \bm{J} \cdot \dd{\bm{S}} = \oint_{4\pi} (\bm{J} \cdot \bm{u}_r) \, r^2 \dd{\Omega} \overset{(7)}{\propto} \frac{r^2}{r^{2\gamma}} \oint_{4\pi} f(\Omega) \dd{\Omega} = \frac{C}{r^{2\gamma - 2}} \quad,
\end{equation}
dove $C$ è una costante reale che racchiude il risultato dell'integrazione sulle sole variabili angolari. In $(7)$ abbiamo sfruttato il fatto che il gradiente di $f$ è proporzionale al versore $\bm{u}_r$ e constatato che l'integrazione avviene \emph{solo} sull'angolo solido $4\pi$ e non interessa minimamente la parte radiale.
Avendo così separato analiticamente la dipendenza radiale, possiamo ora applicare rigorosamente il limite per $r \to \infty$ all'intera espressione:
\[
    \lim_{r \to \infty} \oint_{\partial \Omega} \bm{J} \cdot \dd{\bm{S}} \propto \lim_{r \to \infty} \frac{1}{r^{2\gamma-2}} \quad,
\]
che è infinitesimo solo nel caso in cui l'esponente a denominatore è strettamente positivo:
\begin{equation}
        2\gamma - 2 > 0 \implies \gamma > 1 \quad.
\end{equation}

\subsubsection{Considerazioni finali}
Abbiamo quindi scoperto che affinché la probabilità totale di trovare una particella entro il volume in cui è stata confinata sia costante nel tempo, serve che la funzione d'onda decada strettamente più velocemente di $r^{-1}$ per $r \to \infty$. Coerentemente, in questa condizione abbiamo: 
\[
    \dv{\mathcal{P}}{t} = 0 \implies \mathcal{P} = \text{cost.}
\] 

\subsection{Altre fondamentali proprietà della funzione d'onda}
Quanto discusso sino ad ora ci ha permsso di gettare le dovute basi necessarie per elevare la trattazione e analizzare due ulteriori caratterizzazioni della funzione d'onda: il suo essere parte dello \textit{spazio delle funzioni a quadrato sommabile} $\mathcal{L}^2(\Omega \subseteq \mathbb{R}^3)$ e la \textit{condizione di normalizzazione}.

\subsubsection{Lo spazio delle funzioni a quadrato sommabile $\bm{\mathcal{L}^2}$}
Proviamo ora a guardare con occhi diversi la condizione di unitarietà della probabilità totale prescritta dall'interpretazione di Copenaghen. Il fatto che l'integrale del modulo quadro di $\psi(\bm{r},t)$ sia costantemente 1 sul volume di definizione $\Omega$ della particella implica certamente che:
\[
    1 = \int_\Omega |\psi(\bm{r})|^2 \dd^3{\bm{r}} < +\infty \quad.
\]
Di conseguenza, perché una generica funzione d'onda rappresenti un \emph{sistema fisico}, serve che il suo modulo quadro sia integrabile e non divergente sul dominio di definizione: in una parola, $\psi(\bm{r},t)$ deve essere a \textbf{quadrato sommabile} sul dominio di definizione\footnote{Si noti come la nozione di sommabilità al quadrato può essere estesa all'intero spazio tridimensionale $\mathbb{R}^3$ oppure a un suo sottospazio $\Omega \in \mathbb{R}^3$ e, in quest'ultimo caso, si dovrebbe parlare di funzioni \emph{localmente} a quadrato sommabile $\psi \in \mathcal{L}^2_{loc}$. Per semplicità si omette la dicitura "{loc}" e si intenderà sempre \guillemetleft \emph{a quadrato sommabile sul dominio di definizione $\Omega$}\guillemetright.}. 

L'insieme delle funzioni a quadrato sommabile costituisce uno \emph{spazio vettoriale} identificato dal simbolo $\mathcal{L}^2(\Omega)$\footnote{Si anticipa che $\mathcal{L}^2(\Omega \subseteq \mathbb{R}^3)$ è anche uno \emph{spazio di Hilbert}. Questa peculiarità è di fondamentale importanza per la coerenza matematica della meccanica quantistica e verrà discussa più approfonditamente nel prosieguo della trattazione.}. Si ricorda che per dimostrare rigorosamente questa conclusione sarebbe necessario dimostrare le fantomatiche \emph{7 proprietà} che caratterizzano un generico spazio vettoriale. Dimostriamo di seguito, per semplicità, che $\mathcal{L}^2$ è necessariamente almeno un \emph{sottospazio vettoriale}: questo ci permetterà di sfruttare comodamente la nozione di combinazione lineare (e dunque il principio di sovrapposizione) durante l'analisi della S.E. di un sistema.
\begin{proprietà}[Esistenza dell'elemento neutro]
    L'insieme delle funzioni a quadrato sommabile $\mathcal{L}^2(\Omega)$ ammette l'\emph{elemento neutro} $\psi = 0$.
\end{proprietà}
\begin{proof}
    La dimostrazione è banale se si nota che $\psi(\bm{r},t) = 0$ è evidemtemente a quadrato sommabile su qualsiasi sottoinsieme di $\mathbb{R}^3$ (e quindi anche su $\mathbb{R}^3$).
\end{proof}
\begin{proprietà}[Chiusura rispetto alla somma]
    L'insieme delle funzioni a quadrato sommabile $\mathcal{L}^2(\Omega)$ è \emph{chiuso} rispetto all'operazione di combinazione lineare.
\end{proprietà}
\begin{proof}
    Dobbiamo dimostrare che una generica combinazione lineare di funzioni a quadrato sommabile è a quadrato sommabile. Si cosiderino a tal fine $n$ generiche funzioni a quadrato sommabile $\psi_i \in \mathcal{L}^2(\Omega)$ e altrettante costanti complesse arbitrarie $c_i \in \mathbb{C}$. Per ipotesi, deve valere:
    \[
        \int_{\Omega} \dd^3{\bm{r}} \, |\psi_i(\bm{r},t)|^2 < +\infty \quad.
    \]
    Si definisce allora una generica funzione $\Psi$ come combinazione lineare delle $\psi_i$:
    \[
        \Psi(\bm{r},t) = \sum_{i=1}^n c_i \psi_i(\bm{r},t) \quad.
    \]
    Dal momento che il dominio delle $\psi_i$ è $\Omega$ per ipotesi, deve valere $\text{dom}(\Psi) = \Omega$. Si stima dall'alto allora l'integrale sul volume di definizione del suo modulo quadro:
    \[
        \begin{aligned}
            \int_{\Omega} \dd^3{\bm{r}} \, |\Psi(\bm{r},t)|^2 &= \int_{\Omega} \dd^3{\bm{r}} \, \left|\sum_{i=1}^n c_i \psi_i(\bm{r},t)\right|^2 \overset{(8)}{\leq} \int_{\Omega} \dd^3{\bm{r}} \,\sum_{i=1}^n |c_i \psi_i(\bm{r},t)|^2 \\
            &= \sum_{i=1}^n |c_i|^2 \int_{\Omega} \dd^3{\bm{r}} \,|\psi_i(\bm{r},t)|^2 < +\infty 
        \end{aligned}
    \]
    perché ogni integrale è non divergente per ipotesi. In $(8)$ si è sfruttata la \emph{disuguaglianza di Minkowski} per $\mathcal{L}^2$. 
\end{proof}

\begin{proprietà}[Chiusura rispetto al prodotto per scalare]
    L'insieme delle funzioni a quadrato sommabile $\mathcal{L}^2(\Omega)$ è \emph{chiuso} rispetto all'operazione di prodotto per scalare.
\end{proprietà}
\begin{proof}
    Il prodotto per uno scalare non può modificare la proprietà di non divergenza dell'integrale di una funzione.  
\end{proof}

Le implicazioni che derivano da questa discussione sono molteplici\footnote{Si noti che, con il fine di alleggerire la notazione, si ometterà spesso il dominio di definizione, lasciando spazio al simbolo $\mathcal{L}^2$ per identificare il generico spazio delle funzioni a quadrato sommabile.}, in particolare otteniamo immediatamente un'ulteriore garanzia del fatto che assumere come soluzione dell'equazione di Schrödinger una combinazione lineare di diverse funzioni d'onda anch'esse soluzione della S.E. non è solo lecito dal punto di vista dell'equazione differenziale, ma anche dal profilo strettamente algebrico.

\subsubsection{La condizione di normalizzazione}
Abbiamo sin'ora discusso quale sia il tipo di funzioni d'onda che hanno senso fisico, ovvero che garantiscono un rigoroso fondamento matematico all'interpretazione probabilistica di Born. È necessario a questo punto imporre ancora che la probabilità totale di trovare una particella nel suo volume di definizione non sia semplicemente non divergente, bensì \textbf{unitaria}. 

Si consideri una particella confinata, per semplicità, nell'intero spazio tridimensionale $\Omega \equiv \mathbb{R}^3$. Richiediamo quindi che l'integrale del modulo quadro della funzione d'onda ad essa associata sia \emph{normalizzato}, ovvero che restituisca esattamente $1$:
\begin{equation}
    \int_{\mathbb{R}^3}  \dd^3{\bm{r}} \, |\psi(\bm{r},t)|^2 = 1 \quad.
    \label{eq:normalizzazione_ps}
\end{equation}
Per verificare la fattibilità di questa richiesta sfruttiamo il comportamento asintotico $\psi \sim f(r)/r^\gamma$  analizzato in precedenza e, con il fine di evitare patologie matematiche nell'origine dovute all'uso di un'espressione valida solo per grandi distanze, scindiamo il dominio di integrazione in una sfera di raggio $R$ arbitrariamente grande e nel dominio asintotico da $R$ a $+\infty$. All'interno della sfera $r \le R$, assumiamo che la funzione $\psi$ sia continua e limitata, restituendo un integrale certamente finito che chiameremo $I_0$. Passando in coordinate sferiche per il termine asintotico, otteniamo:
\begin{equation}
\begin{aligned}
    \int_{\mathbb{R}^3} |\psi|^2 \dd^3{\bm{r}} &= \int_{r \le R} |\psi|^2 \dd^3{\bm{r}} + \int_R^\infty \int_0^\pi \int_0^{2\pi} \frac{|f(r)|^2}{r^{2\gamma}} r^2 \sin \theta \dd{r} \dd{\theta} \dd{\phi} \\
    &= I_0 + \int_R^\infty \frac{|f(r)|^2}{r^{2\gamma-2}} \dd{r} {\int_0^\pi \sin\theta \dd{\theta}} {\int_0^{2\pi} \dd{\phi}} \\
    &= I_0 + 4\pi \int_R^\infty \frac{|f(r)|^2}{r^{2\gamma-2}} \dd{r} \quad.
    \label{eq:int_mod_quad}
\end{aligned}
\end{equation}
Ricordando che l'integrale del modulo quadro della funzione d'onda rappresenta la probabilità totale, che è necessariamente minore di infinito, l'ultimo integrale improprio deve necessariamente convergere a un valore finito. Essendo $|f(r)|^2$ limitata asintoticamente per definizione, il criterio di convergenza degli integrali impone che l'esponente a denominatore sia strettamente maggiore di $1$:
\begin{equation}
    2\gamma - 2 > 1 \implies \gamma > \frac{3}{2} \quad,
\end{equation}
risultato assolutamente coerente con il precedente $\gamma > 1$.


\begin{concettochiave}{Condizione di normalizzazione della funzione d'onda e costanza della probabilità totale nel tempo}
    \noindent
    Perché ad una generica funzione d'onda $\psi(\bm{r},t)$ ambientata in $\mathbb{R}^3$ sia associata una probabilità totale costante nel tempo, e quindi sia in accordo con l'interpretazione di Copenaghen, serve che 
    \[
        \psi = o\left(\frac{1}{r} \right)
    \] 
    e quindi che $\psi$ sia \emph{o-piccolo} di $1/r$ per $r \to \infty$.
    La condizione di normalizzazione invece, ottenuta imponendo che $\psi \in \mathcal{L}^2$, prescrive che 
    \[
        \psi = {o}\left(\frac{1}{r^{\frac{3}{2}}} \right) \quad.
    \]
    È evidente quindi che una generica funzione d'onda deve \textbf{necessariamente} essere a quadrato sommabile in quanto questa condizione garantisce automaticamente che la probabilità totale sia indipendente dal tempo.
\end{concettochiave}

\noindent
A questo punto, supponendo che la \eqref{eq:int_mod_quad} restituisca il valore $\alpha \in \mathbb{R}^+$, per ottenere la condizione di probabilità unitaria prescritta in \eqref{eq:normalizzazione_ps}, è sufficiente sfruttare la linearità dell'equazione di Schrödinger e ridefinire la funzione d'onda moltiplicandola per l'opportuna \emph{costante di normalizzazione} $A = 1/\sqrt{\alpha}$:
\[
    \psi_{\text{norm}}(\bm{r}, t) = \frac{1}{\sqrt{\alpha}} \psi(\bm{r}, t) \implies \int_{\mathbb{R}^3} |\psi_{\text{norm}}|^2 \dd^3{\bm{r}} = \frac{1}{\alpha} \underbrace{\int_{\mathbb{R}^3} |\psi|^2 \dd^3{\bm{r}}}_{\alpha} = 1 \quad.
\]

Da questo momento in poi, assumeremo sempre che le funzioni d'onda trattate siano state preventivamente normalizzate, garantendo la perfetta validità del postulato di Born.

\section{Il prodotto scalare di funzioni d'onda}
Avendo identificato il "mondo" delle funzioni d'onda come uno spazio vettoriale, è necessario sviscerare ancora una questione: come descrivere una forma di \textit{prodotto scalare} in $\mathcal{L}^2$.

Date due funzioni d'onda a quadrato sommabile $\psi$ e $\phi$, definiamo il loro \textbf{prodotto scalare} come:
\begin{equation}
    (\psi, \phi) \equiv \int_{\mathbb{R}^3} \dd^3{\bm{r}}\, \psi(\bm{r},t)^*\phi(\bm{r},t) \quad.
    \label{eq:prod_scal}
\end{equation}
Nel caso del prodotto scalare di una funzione per se stessa, recuperiamo la nozione di norma quadra:
\[
    \norm{\psi}^2 \equiv (\psi,\psi) = \int_{\mathbb{R}^3} \dd^3{\bm{r}}\, \psi(\bm{r},t)^*\psi(\bm{r},t) = \int_{\mathbb{R}^3}\dd^3{\bm{r}}\, |\psi(\bm{r},t)|^2 = 1 \quad.
\]
Si noti come la norma quadra di una genrica funzione d'onda associata ad uno \emph{stato fisico} non è altro che la probabilità totale di trovare la particella entro il suo spazio di esistenza ed è, di conseguenza, \emph{sempre unitaria} per l'interpretazione di Copenaghen\footnote{Si presti attenzione, in ogni caso, al fatto che la norma è una forma definita positiva e dunque, per definizione, può assumere qualsiasi valore in $\mathbb{R}^+$. La restrizione a un valore unitario vale solo, per l'appunto, nel caso di \emph{stati fisici}.}. 

\subsection{Proprietà del prodotto scalare di funzioni d'onda}
Sono evidenti alcune analogie con il prodotto scalare caratteristico degli spazi euclidei che ricordiamo essere definito come segue. 
Dati due vettori $\bm{a}, \bm{b} \in \mathbb{R}^n$, il loro prodotto scalare è dato da:
\[
    (\bm{a}, \bm{b}) \equiv \bm{a} \cdot \bm{b} = \sum_{i = 1}^n a_ib_i \quad,
\]
con $a_i$ e $b_i$ l'$i$-esima componente dei vettori. Il prodotto scalare euclideo è pertanto definito come una forma \emph{bilineare e simmetrica}.

Nel caso quantistico, invece, si perdono queste "comode" caratteristiche in luogo di due nuove peculiarità che ci accompagneranno nei calcoli. La \eqref{eq:prod_scal} restituisce un costrutto algebrico evidentemente non bilineare, bensì \textbf{sesquilineare}, ovvero lineare nel secondo argomento e antilineare (o lineare coniugato) nel primo. Abbiamo infatti:
\[
    (c_1\psi, c_2\phi) = \int_{\mathbb{R}^3} \dd^3{\bm{r}}\, c_1^*\psi(\bm{r},t)^*c_2\phi(\bm{r},t) = c_1^*c_2 \int_{\mathbb{R}^3} \dd^3{\bm{r}}\, \psi(\bm{r},t)^*\phi(\bm{r},t) = c_1^*c_2 (\psi, \phi), \quad c_i \in \mathbb{C} \quad.
\]
Il prodotto scalare complesso è inoltre \emph{non simmetrico} (essendo in generale $\psi^*\phi \neq \psi\phi^*$) ma \textbf{simmetrico in senso hermitiano}, ovvero:
\[
    (\psi, \phi) = \int_{\mathbb{R}^3} \dd^3{\bm{r}}\, \psi(\bm{r},t)^*\phi(\bm{r},t) = \biggl( \int_{\mathbb{R}^3} \dd^3{\bm{r}}\, \phi(\bm{r},t)^*\psi(\bm{r},t) \biggr)^* = (\phi, \psi)^* \quad.
\]


%!TEX root = ../dispense_QP.tex

\chapter{Gli operatori fisici, la quantizzazione canonica e il commutatore}
Il passaggio a una \emph{meccanica ondulatoria}, sancito dalle discussioni sull'interpretazione di Copenaghen della funzione d'onda e sul significato della densità di probabilità quantistica, segna uno spartiacque definitivo rispetto alla visione deterministica della meccanica classica, in cui le equazioni del moto definiscono una traiettoria esatta. Quest'ultima deve pertanto essere necessariamente, perlomeno, \emph{rivisitata} in favore di un framework intrinsecamente probabilistico. 

\guillemetleft\emph{Ma perché non abbandonarla direttamente}\guillemetright? La risposta risiede in quanto analizzato sino ad ora: dal momento che le densità di probabilità quantistica e "classica“ si comportano allo stesso modo, qualcosa deve necessaiamente essere conservato e cercheremo di capire \emph{cosa effettivamente non cambia} in quanto segue.  
Prima di procedere, è pedagogico capire come si trasformano le entità con cui \emph{qualsiasi} equazione fisica deve fare i conti: le \emph{grandezze fisiche}. Come si può intuire nella trattazione dell'equazione di Schrödinger, in questo nuovo nuovo paradigma le grandezze fisiche perdono la loro natura di semplici \emph{campi} (scalari o vettoriali) per assumere la veste matematica di \textbf{operatori lineari} che agiscono sulla funzione d'onda "producendo“ la controparte classica. Questa transizione strutturale porta con sé la conseguenza algebrica più dirompente dell'intera teoria: a differenza delle variabili classiche, in cui il prodotto di due coordinate appartenenti allo spazio delle fasi $q_\alpha$ e $p_\alpha$ è commutativo $q_\alpha p_\alpha = p_\alpha q_\alpha$, il prodotto tra operatori quantistici è generalmente \emph{non commutativo}.

Per quantificare e studiare rigorosamente questa asimmetria, introdurremo lo strumento cardine del formalismo quantistico: il \textbf{commutatore}. Esso non costituisce un mero artificio matematico, ma rappresenta il fulcro del \textbf{principio di quantizzazione canonica} formulato da Dirac. Questo principio stabilisce una mappa formale e rigorosa tra la struttura simplettica dello spazio delle fasi classico e l'algebra degli operatori, associando le \emph{parentesi di Poisson} tra due grandezze dinamiche $\{A, B\}$ al commutatore tra i rispettivi operatori, secondo la celebre prescrizione:
\begin{equation}
    \{A, B\} \rightarrow \frac{1}{i\hbar}[\hat{A}, \hat{B}] \quad.
\end{equation}


\section{Gli operatori fisici}
La discussione sull'equazione di Schrödinger ci ha mostrato che alcune grandezze della fisica classica, come l'energia meccanica o la quantità di moto, si traducono nel contesto quantistico in \textbf{operatori} agenti sulla funzione d'onda. Ci tocca fare, come sempre, un piccolo -\emph{o grande}- passo di astrazione: caratterizzeremo d'ora in poi qualsiasi \emph{grandezza fisica osservabile}, ovvero misurabile, tramite un operatore a questa associato, in ossequio ad un principio che discuteremo nella prossima sezione, ovvero la cosiddetta \textbf{regola della quantizzazione canonica}. In questo contesto, distinguiamo due macro-classi di operatori: gli operatori moltiplicativi e gli operatori derivativi.

\subsection{Operatori moltiplicativi}
L'azione di determinati operatori si riduce a un semplice prodotto algebrico. L'esempio cardine è l'\textbf{operatore posizione}. Denotato il vettore posizione con $\bm{r} = (x_1, x_2, x_3)$, le sue singole componenti $\hat{x}_i$ sono operatori che agiscono sulla funzione d'onda restituendola moltiplicata per la rispettiva coordinata:
\begin{equation}
    \hat{x}_i \psi(\bm{r}, t) = x_i \psi(\bm{r}, t) \quad \text{per } i = 1, 2, 3 \quad.
\end{equation}
Raggruppando le componenti in una notazione vettoriale compatta, possiamo formalizzare l'operatore vettore posizione $\hat{\bm{r}}$ come:
\begin{equation}
    \hat{\bm{r}} = 
    \begin{pmatrix}
        \hat{x}_1 \\
        \hat{x}_2 \\
        \hat{x}_3
    \end{pmatrix}
    \implies \hat{\bm{r}}\psi(\bm{r},t) = \bm{r}\psi(\bm{r},t) \quad.
\end{equation}
Questa logica si estende in modo naturale a qualsiasi funzione scalare dipendente unicamente dalla posizione: l'operatore a essa associato agirà anch'esso in via puramente moltiplicativa. Se consideriamo, ad esempio, l'energia potenziale $U(\bm{r})$:
\begin{equation*}
    \hat{U}(\bm{r}) \psi(\bm{r},t) = U(\bm{r})\psi(\bm{r},t) \quad.
\end{equation*}

\subsection{Operatori derivativi}
La seconda classe è formata dagli operatori che comportano una derivazione spaziale. Come abbiamo dedotto dall'ansatz di Schrödinger, l'\textbf{operatore quantità di moto} $\hat{\bm{p}}$ ha natura differenziale:
\begin{equation}
    \hat{\bm{p}} = -i\hbar\nabla \implies \hat{p}_j \psi(\bm{r},t) = -i\hbar \pdv{\psi(\bm{r},t)}{x_j} \quad.
\end{equation}
Di conseguenza, anche gli operatori costruiti a partire dalla quantità di moto erediteranno questa natura differenziale. L'operatore energia cinetica $\hat{K}$ ad esempio, essendo proporzionale a $\hat{\bm{p}}^2$, coinvolgerà le derivate seconde spaziali, tramutandosi nel Laplaciano:
\begin{equation*}
    \hat{K} = \frac{\hat{\bm{p}} \cdot \hat{\bm{p}}}{2m} = -\frac{\hbar^2}{2m}\nabla^2 \quad.
\end{equation*}

\subsubsection{Casi particolari}
Esistono ancora alcuni operatori contraddistinti da un'azione che non è imputabile unicamente ad alcuna delle due classi appena discusse. Uno fra tutti, l'operatore \textbf{momento angolare} è definito in similitudine all'espressione classica:
\begin{equation}
    \bm{L} = \bm{r} \times \bm{p} \to \hat{\bm{L}} \equiv \hat{\bm{r}} \times \hat{\bm{p}} =
    \begin{pmatrix}
        \hat{L}_1 \\
        \hat{L}_2 \\
        \hat{L}_3
    \end{pmatrix} 
    = 
    \begin{pmatrix}
        \hat{x}_2\hat{p}_3 - \hat{x}_3\hat{p}_2 \\
        \hat{x}_3\hat{p}_1 - \hat{x}_1\hat{p}_3 \\
        \hat{x}_1\hat{p}_2 - \hat{x}_2\hat{p}_1
    \end{pmatrix}
    = -i\hbar
    \begin{pmatrix}
        x_2 \partial_{x_3} - x_3 \partial_{x_2} \\
        x_3 \partial_{x_1} - x_1 \partial_{x_3} \\
        x_1 \partial_{x_2} - x_2 \partial_{x_1}
    \end{pmatrix} \quad.
\end{equation}
In modo molto più compatto, sfruttando il \emph{tensore totalmente antisimmetrico di Levi-Civita} (o \emph{simbolo di Ricci}) $\epsilon_{ijk}$ e la convenzione di Einstein sugli indici ripetuti, la generica componente $i$-esima dell'operatore momento angolare si scrive:
\[
    \hat{L}_i = \epsilon_{ijk} \hat{x}_j \hat{p}_k = -i\hbar \epsilon_{ijk} x_j \partial_{x_k} \quad,
\]
relazione che si rivelerà di vitale importanza per calcolare in modo sistematico i commutatori tra le componenti di $\hat{\bm{L}}$. Si noti che l'azione dell'operatore momento angolare è sia di carattere moltiplicativo sia derivativo.

\begin{approfondimento}[title={Proprietà del simbolo di Levi-Civita}]
    \noindent
Il simbolo $\epsilon_{ijk}$ è definito come:
\begin{equation*}
    \epsilon_{ijk} = \begin{cases} 
    +1 & \text{se } (i,j,k) \text{ è una permutazione pari di } (1,2,3) \\
    -1 & \text{se } (i,j,k) \text{ è una permutazione dispari} \\
    0  & \text{se due o più indici sono uguali} \quad .
    \end{cases}
\end{equation*}
Risulta particolamrmente utile per alleggerire la notazione del prodotto vettoriale. Dati infatti due generici vettori $\bm{a}$ e $\bm{b}$, abbiamo:
\[
    \bm{a} \times \bm{b} = \bm{c} = (c_1, c_2, c_3) \quad.
\]
La $i$-esima componente di $\bm{c}$ può essere ricavata tramite la notazione di Levi-Civita in maniera quasi immediata:
\[
    c_i = \epsilon_{ijk}a_j b_k \quad.
\]
Proviamo, a scopo di esercizio, a calcolare $c_1$, ponendo quindi $i = 1$. Dalla definizione del simbolo di Ricci, segue che se $j \, \text{e/o} \, k = i$ allora $\epsilon_{ijk} = 0$. Per avere termini non nulli, $j$ e $k$ devono essere diversi e uguali a 2 o 3: abbiamo due differenti casi:
\[
    \begin{aligned}
        j &= 2 \quad k = 3 \rightarrow \text{permutazione pari} \rightarrow \epsilon_{123} = 1 \\
        j &= 3 \quad k = 2 \rightarrow \text{permutazione dispari} \rightarrow \epsilon_{132} = -1 \quad.
    \end{aligned}
\]
Sommando i due contributi infine otteniamo:
\[
    c_1 = \epsilon_{1jk}a_jb_k = \epsilon_{123} a_2b_3 + \epsilon_{132}a_3b_2 = a_2b_3 - a_3b_2 \quad,
\]
come ci auspicavamo. 

N.B. In un certo tal senso, il simbolo di Ricci nasconde quindi una sommatoria sui termini diversi dal primo. 
\end{approfondimento}

\subsection{Alcune proprietà degli operatori}
Dobbiamo a questo punto sviscerare le caratteristiche degli operatori fisici. Segue dunque una fondamentale discussione di carattere fortemente matematico-teorico che ci permetterà di avere sufficienti basi per affrontare argomenti più avanzati. Si noti come molti dei risultati presentati sono "tratti" dai fondamenti dell'algebra lineare e, pertanto, non verranno dimostrati.

La prima domanda che ci facciamo è: -\emph{cosa restituisce l'applicazione di un operatore?}-. La trattazione introduttiva sull'equazione di Schrödinger può fornire dei suggerimenti ma adesso formalizzeremo l'intuizione nel seguente teorema:

\begin{teorema}[Operatori fisici come applicazioni lineari]
Sia $\hat{A}$ un operatore e $\psi, \phi \in \mathcal{L}^2$ funzioni d'onda. L'applicazione di $\hat{A}$ è \emph{lineare} e abbiamo:
\[
    \hat{A}(c_1 \psi + c_2 \phi) = c_1\hat{A}\psi + c_2\hat{A}\phi \quad c_1, c_2 \in \mathbb{C} \quad.
\]
\end{teorema}
Questa linearità è conservata anche nella somma di due operatori: abbiamo infatti la seguente proprietà:

\begin{proprietà}[Somma di operatori]
Siano $\hat{A}$ e $\hat{B}$ due opratori fisici. L'applicazione lineare definita dalla loro somma è pari alla somma delle loro applicazioni lineari:
\[
    (\hat{A} + \hat{B})\psi = \hat{A}\psi + \hat{B}\psi \quad \psi \in \mathcal{L}^2 \quad.
\]
\end{proprietà} 
Ancora, per quanto riguarda il prodotto fra operatori si ha:
\begin{definizione}[Prodotto di operatori]
    Siano $\hat{A}$ e $\hat{B}$ due opratori fisici. L'applicazione lineare definita dal loro prodotto si ottiene come:
\[
    (\hat{A}\hat{B})\psi = \hat{A} (\hat{B}\psi) \quad.
\]
\end{definizione}
Si noti come il prodotto di due operatori, anche se associativo non è assolutamente commutativo. Infatti $\hat{A}\hat{B} \neq \hat{B}\hat{A}$. Possiamo a questo punto facilmente definire l'operatore identità e l'operatore zero: 
\begin{enumerate}
    \item \textbf{Operatore identità}:
    $
        \mathds{1}\hat{A} = \hat{A}\mathds{1} = \hat{A}, \quad \mathds{1}\psi = \psi .
    $
    \item \textbf{Operatore zero}:
    $
        0\hat{A} = \hat{A}0 = 0, \quad 0\psi = 0 .
    $
\end{enumerate}


Un altro risultato fondamentale è dato dal teorema che segue. È di vitale importanza capire se l'applicazioni di un operatore restituisce una funzione d'onda a quadrato sommabile:

\begin{teorema}[Operatori fisici come endomorfismi]
Sia $\hat{A}$ un operatore e $\psi \in \mathcal{L}^2$ una funzione d'onda. (Sotto opportune condizioni di regolarità) l'applicazione lineare rappresentata da $\hat{A}$ è un \emph{endomorfismo} con $\hat{A}: \mathcal{L}^2 \to \mathcal{L}^2$ e abbiamo: 
\[
    \hat{A} \psi = \phi \quad \psi,\phi \in \mathcal{L}^2 \quad.
\]
\end{teorema}
Ancora un piccolo passo. Attribuiremo agli operatori fisici la caratteristica di essere \emph{hermitiani}\footnote{Che vede il suo corrispondente dell'algebra lineare elementare nelle matrici complesse hermitiane (o simmetriche in campo reale).}, assunzione che si rivelerà incredibilmente importante nel susseguirsi della trattazione. Definiamo innanzitutto un \emph{operatore aggiunto} e \emph{autoaggiunto}:
\begin{definizione}[Operatore aggiunto e autoaggiunto]
    Sia $\hat{A}$ un operatore. Si definisce \emph{operatore aggiunto} di $\hat{A}$ l'operatore $\hat{A}^{\dagger}$ per cui vale:
\[
    (\phi, \hat{A}\psi) \equiv (\hat{A}^{\dagger}\phi, \psi)\quad \phi, \psi \in \mathcal{L}^2 \quad.
\]
Dalle proprietà del prodotto scalare complesso, abbiamo poi che:
\[
    (\phi, \hat{A}\psi) = (\hat{A}^{\dagger}\phi, \psi) = (\psi, \hat{A}^{\dagger}\phi)^* \quad.
\]
\end{definizione}

Diciamo che $\hat{A}$ è \emph{autoaggiunto} se vale $\hat{A} = \hat{A}^{\dagger}$.
Siamo pronti per comprendere il seguente risultato:
\begin{teorema}[Operatori fisici come operatori autoaggiunti]
    Sia $\hat{A}$ un operatore fisico. Allora $\hat{A}$ è autoaggiunto, o \emph{hermitiano}, e vale:
    \begin{equation}
    (\phi, \hat{A}\psi) = (\hat{A}^{\dagger}\phi, \psi) = (\hat{A}\phi, \psi) = (\psi, \hat{A}\phi)^* \quad \phi, \psi \in \mathcal{L}^2 \quad.
    \label{eq:hermitiano}
    \end{equation}
\end{teorema}
Possiamo quindi provare un'ultima interessante proprietà che tornerà molto utile nei calcoli:
\begin{proprietà}[Operatori fisici]
    Siano $\hat{A}$ e $\hat{B}$ due operatori fisici. Vale la seguente relazionare
\[
    \big(\hat{A}\hat{B}\big)^\dagger = \hat{B}^\dagger \hat{A}^\dagger \quad.
\]
\end{proprietà}
\begin{proof}
Dalla definizione di prodotto scalare, prese $\psi, \phi \in \mathcal{L}^2$, abbiamo:
\[
    \big( \psi, (\hat{A}\hat{B})^\dagger\phi \big) \overset{(1)}{=} \big( \phi, (\hat{A}\hat{B})\psi \big)^* = ( (\hat{A}\hat{B})\psi,\phi  \big)\quad,
\]
dove in $(1)$ abbiamo sfruttato l'hermitianità del prodotto di operatori hermitiani. Continuiamo:
\[
    ( (\hat{A}\hat{B})\psi,\phi  \big) = ( \hat{A}(\hat{B}\psi),\phi  \big) = ( \hat{B}\psi,\hat{A}^\dagger\phi  \big) = (\psi,\hat{B}^\dagger\hat{A}^\dagger\phi  \big) = (\hat{B}^\dagger\hat{A}^\dagger\phi, \psi  \big)^*\quad.
\]
Essendo $\hat{A}\hat{B}$ hermitiano, concludiamo che $\big(\hat{A}\hat{B}\big)^\dagger = \hat{B}^\dagger \hat{A}^\dagger$.
\end{proof}

\section{La regola della quantizzazione canonica}
Per cogliere appieno la natura degli operatori fisici appena introdotti, dobbiamo ancora trovare risposta al seguente quesito: -\emph{come si costruisce l'operatore quantistico corrispondente a una determinata grandezza classica?}-. La risposta risiede nel \textbf{Principio di Quantizzazione Canonica}, il formidabile ponte logico-matematico gettato da Paul Dirac tra la Meccanica Analitica e la Meccanica Quantistica.

Nel formalismo Hamiltoniano della meccanica classica, lo stato di un sistema è individuato da un punto nello \emph{spazio delle fasi}, descritto dalle coordinate generalizzate $q_\alpha$ e dai rispettivi momenti coniugati $p_\alpha$. Qualsiasi osservabile fisica $A$ è una funzione reale di tali variabili, $A(q,p,t)$. La struttura dinamica di questo spazio è governata dalle \textbf{parentesi di Poisson}, definite per due generiche osservabili $A$ e $B$ come:
\begin{equation*}
    \{A, B\} = \sum_\alpha \left( \pdv{A}{q_\alpha}\pdv{B}{p_\alpha} - \pdv{A}{p_\alpha }\pdv{B}{q_\alpha} \right) \quad.
\end{equation*}
Dalla definizione, seguono immediatamente le relazioni fondamentali per le variabili canoniche stesse:
\begin{equation}
    \{q_i, q_j\} = 0, \quad \{p_i, p_j\} = 0, \quad \{q_i, p_j\} = \delta_{ij} \quad,
\end{equation}
dove $\delta_{ij}$ è la Delta di Kronecker.

Sveliamo un grande arcano: in realtà, non è necessario \emph{abbandonare} completamente la teoria classica, per cui Newton prima e Lagrange, Eulero, Hamilton e Jacobi poi, hanno duramente lavorato. Basta \emph{riadattarla}. Il colpo di genio di Dirac fu infatti notare che la struttura algebrica delle parentesi di Poisson (che formano una cosiddetta \emph{algebra di Lie}) è formalmente identica all'algebra dei cosiddetti \emph{commutatori}, protagonisti del discorso che verrà. Dirac postulò quindi che la transizione al regime quantistico non dovesse invalidare l'impalcatura della meccanica classica, ma piuttosto richiederne una "traduzione" in un nuovo linguaggio matematico. 

Il Principio di Quantizzazione Canonica si articola in due prescrizioni operative:
\begin{enumerate}
    \item Alle variabili classiche spaziali e d'impulso $(q_i, p_i)$ si sostituiscono operatori lineari autoaggiunti $(\hat{q}_i, \hat{p}_i)$ agenti sullo spazio di Hilbert degli stati del sistema.
    \item A ogni parentesi di Poisson tra due grandezze classiche si fa corrispondere il \textbf{commutatore} tra i rispettivi operatori quantistici, diviso per una costante immaginaria proporzionale alla costante di Planck ridotta:
    \begin{equation}
        \{A, B\} \longrightarrow \frac{1}{i\hbar}[\hat{A}, \hat{B}] \quad,
        \label{eq:quantizzazione_dirac}
    \end{equation}
    dove $[\hat{A}, \hat{B}]$ rappresenta appunto il commutatore fra gli operatori $\hat{A}$ e $\hat{B}$.
\end{enumerate}
Questa mappatura impone che gli operatori quantistici obbediscano a specifiche regole di non-commutatività. Andiamo quindi a definire e calcolare esplicitamente questi commutatori.

\section{Il commutatore}
Il primo step di questo processo è quello di prendere ciò che permette di relazionare due generiche grandezze nello spazio delle fasi, le \textbf{parentesi di Poisson}, e tradurlo nel suo corrispondente quantomeccanico: il \textbf{commutatore}. 

\subsection{Definizione operatoriale e proprietà}
Definiamo il commutatore come l'\emph{operatore} che prende in argomento due generici operatori $\hat{A}$ e $\hat{B}$ e resituisce la differenza del loro prodotto incrociato:
\begin{equation}
    [\hat{A}, \hat{B}] = \hat{A}\hat{B} - \hat{B} \hat{A} \quad.
\end{equation}
Per poter giocare serenamente con il commutatore, dobbiamo estrarre dalla definizione alcune proprietà algebriche fondamentali:
\begin{enumerate}
    \item \textbf{Antisimmetria}: $[\hat{A}, \hat{B}] = -[\hat{B}, \hat{A}]$. Abbiamo infatti:
    \[
        [\hat{A}, \hat{B}] = \hat{A}\hat{B} - \hat{B} \hat{A} = - \big( \hat{B} \hat{A} - \hat{A}\hat{B} \big) = -[\hat{B}, \hat{A}] \quad.
    \]
    Segue che $[\hat{A}, \hat{A}] = 0$. 
    \item \textbf{Bilinearità}: così come per il prodotto scalare euclideo vale:
    \[
        \begin{aligned}
        [\hat{A}, c_1\hat{B} + c_2\hat{C}] &= c_1[\hat{A}, \hat{B}] + c_2[\hat{A}, \hat{C}] \\
        [c_1\hat{A} + c_2\hat{B}, \hat{C}] &= c_1[\hat{A}, \hat{C}] + c_2[\hat{B}, \hat{C}]
        \end{aligned}
    \]
    \item \textbf{Distributività} (Regola di Leibniz): $[\hat{A}, \hat{B}\hat{C}] = [\hat{A}, \hat{B}]\hat{C} + \hat{B}[\hat{A}, \hat{C}]$. Infatti:
    \[
    \begin{aligned}
               [\hat{A}, \hat{B}\hat{C}] = \hat{A}(\hat{B}\hat{C}) - (\hat{B}\hat{C})\hat{A} &= (\hat{A}\hat{B})\hat{C} - \hat{B}\hat{A}\hat{C} + \hat{B}\hat{A}\hat{C} - (\hat{B}\hat{C})\hat{A} \\
               &= (\hat{A}\hat{B})\hat{C} - (\hat{B}\hat{A})\hat{C} + \hat{B}(\hat{A}\hat{C}) - \hat{B}(\hat{C}\hat{A}) \\
               &= (\hat{A}\hat{B}-\hat{B}\hat{A})\hat{C} + \hat{B}(\hat{A}\hat{C} - \hat{C}\hat{A}) \\
               &= [\hat{A}, \hat{B}]\hat{C} + \hat{B}[\hat{A}, \hat{C}] \quad.
    \end{aligned}
    \]
    \item \textbf{Identità di Jacobi}: $[\hat{A}, [\hat{B}, \hat{C}]] + [\hat{B}, [\hat{C}, \hat{A}]] + [\hat{C}, [\hat{A}, \hat{B}]] = 0$. Infatti:
    \begin{equation*}
    \begin{aligned}
        [\hat{A}, [\hat{B}, \hat{C}]] &= [\hat{A}, \hat{B}\hat{C} - \hat{C}\hat{B}] \\
        &= \hat{A}(\hat{B}\hat{C} - \hat{C}\hat{B}) - (\hat{B}\hat{C} - \hat{C}\hat{B})\hat{A} \\
        &= \hat{A}\hat{B}\hat{C} - \hat{A}\hat{C}\hat{B} - \hat{B}\hat{C}\hat{A} + \hat{C}\hat{B}\hat{A} \quad.
    \end{aligned}
\end{equation*}
Sfruttando la simmetria per permutazione ciclica $(\hat{A} \to \hat{B}, \hat{B} \to \hat{C}, \hat{C} \to \hat{A})$, otteniamo le espansioni degli altri due termini:
\begin{equation*}
    \begin{aligned}
        [\hat{B}, [\hat{C}, \hat{A}]] &= \hat{B}\hat{C}\hat{A} - \hat{B}\hat{A}\hat{C} - \hat{C}\hat{A}\hat{B} + \hat{A}\hat{C}\hat{B} \quad, \\
        [\hat{C}, [\hat{A}, \hat{B}]] &= \hat{C}\hat{A}\hat{B} - \hat{C}\hat{B}\hat{A} - \hat{A}\hat{B}\hat{C} + \hat{B}\hat{A}\hat{C} \quad.
    \end{aligned}
\end{equation*}
Sommando i tre risultati, osserviamo come i dodici termini si elidano a coppie:
\begin{itemize}[noitemsep]
    \item $\hat{A}\hat{B}\hat{C}$ si elide con $-\hat{A}\hat{B}\hat{C}$ (rispettivamente nel primo e terzo termine);
    \item $-\hat{A}\hat{C}\hat{B}$ si elide con $\hat{A}\hat{C}\hat{B}$ (primo e secondo);
    \item $-\hat{B}\hat{C}\hat{A}$ si elide con $\hat{B}\hat{C}\hat{A}$ (primo e secondo);
    \item $\hat{C}\hat{B}\hat{A}$ si elide con $-\hat{C}\hat{B}\hat{A}$ (primo e terzo);
    \item $-\hat{B}\hat{A}\hat{C}$ si elide con $\hat{B}\hat{A}\hat{C}$ (secondo e terzo);
    \item $-\hat{C}\hat{A}\hat{B}$ si elide con $\hat{C}\hat{A}\hat{B}$ (secondo e terzo).
\end{itemize}
Tutti i termini si annullano identicamente, dimostrando l'identità.
\end{enumerate}

\subsection{Esempi di commutatori notevoli}
Procediamo a calcolare alcuni commutatori "notevoli" che torneranno utili in computazioni che seguono. Per completare lo svolgimento dei conti che seguono sarà, volta per volta, necessario avanzare richieste più o meno "stringenti" sulla regolarità di $\psi$. 

Utilizzeremo varie tecniche per condurre i calcoli: in alcuni casi, con il fine di esplicitare l'azione differenziale (o moltiplicativa) applicheremo direttamente l'operatore alla generica $\psi$, altrimenti ci limiteremo a ottenere un risultato dipendente esplicitamente da un operatore. 

Sarà tutto in ogni caso più chiaro entro breve.

\subsubsection{Due componenti dell'operatore posizione}
Come prevedibile dalla meccanica hamiltoniana, il risultato è banalmente zero. Per svolgere il conto applichiamo il commutatore alla funzione d'onda e svolgiamo la differenza:
\[
\begin{aligned}
    [\hat{{x}}_i, \hat{{x}}_j]\psi &= \hat{x}_i (\hat{x}_j \psi) - \hat{x}_j (\hat{x}_i \psi) \\
    &= \hat{x}_i x_j\psi - \hat{x}_j x_i\psi \\
    &= {x}_i x_j\psi - {x}_j x_i\psi = 0 \quad.
\end{aligned}
\]
Poiché assumiamo $\psi \neq 0$, deve essere $[\hat{{x}}_i, \hat{{x}}_j] = 0$.

\subsubsection{Due componenti dell'operatore quantità di moto}
Anche se il risultato è ancora sufficientemente prevedibile, come nel caso precedentemente, dovremo fare un'assunzione più forte su $\psi$ a parte dal richiedere che non sia banalmente nulla per ottenere il nostro risultato. Procedendo come prima otteniamo:
\[
\begin{aligned}
    [\hat{{p}}_i, \hat{{p}}_j]\psi &= \hat{p}_i (-i\hbar\pdv{\psi}{x_j} ) - \hat{p}_j (-i\hbar\pdv{\psi}{x_i} ) \\
    &= -\hbar^2 \pdv{\psi}{x_j}{x_i} +\hbar^2 \pdv{\psi}{x_i}{x_j} \\
    &\overset{(1)}{=} -\hbar^2 \bigg( \pdv{\psi}{x_j}{x_i} - \pdv{\psi}{x_j}{x_i} \bigg)= 0 \quad.
\end{aligned}
\]
In $(1)$ abbiamo attribuito a $\psi$ la caratteristica di avere almeno il gradiente differenziabile, $\psi \in \mathcal{C}^2(\mathbb{C}, \Omega)$, con $\Omega$ il volume di definizione, per poter applicare il Teorema di Schwarz e commutare le derivate parziali. Abbiamo quindi $[\hat{{p}}_i, \hat{{p}}_j] = 0$.

\subsubsection{Una coordinata e il suo momento canonico coniugato}
Classicamente, le parentesi di Poisson di una coordinata generalizzata $q_\alpha$ e del suo momento coniugato $p_\alpha = \partial_{\dot{q_\alpha}} L$, con $L$ la Lagrangiana del sistema, restituiscono:
\[
\{q_i, p_j\} = \delta_{ij} \quad.
\]
Procediamo ora a verificare se vale lo stesso nel mondo quantistico. Scegliamo come coordinate quelle che definiscono lo spazio delle fasi: una componente della posizione e una della quantità di moto.
\begin{equation*}
    \begin{aligned}
    [\hat{x_i}, \hat{p_j}]\psi &= \biggl(-i\hbar x_i\pdv{}{x_j} + i\hbar \pdv{}{x_j}x_i \biggr)\psi   \\
    &= -i\hbar x_i\pdv{\psi}{x_j} + i\hbar \pdv{(x_i\psi)}{x_j} \\
    &= -i\hbar x_i\pdv{\psi}{x_j} + -i\hbar x_i\pdv{\psi}{x_j} + i\hbar \delta_{ij}\psi = i\hbar \delta_{ij}\psi \quad,
    \end{aligned}
\end{equation*}
dove $\delta_{ij}$ rappresenta la \emph{Delta di Kronecker} di indici $i,j$. Essendo in generale $\psi \neq 0$, abbiamo $[\hat{x_i}, \hat{p_j}] = i\hbar \delta_{ij}$

\begin{quote}
\emph{Nota: formalmente, il risultato di un commutatore è un operatore, per cui andrebbe esplicitato l'operatore identità $\mathds{1}$. Per semplicità notazionale, tale termine verrà spesso omesso nei calcoli successivi.}
\end{quote}

Riflettendo sulle relazioni che scaturiscono dall'applicazione delle parentesi di Poisson alle variabili posizione e momento coniugato, quanto ottenuto tramite il commutatore differisce per un fattore costante pari a $i\hbar$:
\[
    \{x_j, p_j\} \rightarrow \frac{1}{i\hbar}[\hat{x}_j, \hat{p}_i] \quad.
\]
Come ampiamente preannunciato dal postulato di quantizzazione di Dirac, il risultato quantistico ricalca esattamente la struttura simplettica classica, a meno del fondamentale fattore $i\hbar$.

\subsubsection{Una componente dell'operatore quantità di moto e una funzione della posizione}
Procediamo come al solito:
\[
\begin{aligned}
    [\hat{{p}}_i, \hat{f}(\bm{r})]\psi &= \hat{p}_i (f(\bm{r})\psi) - \hat{f}(\bm{r})(-i\hbar\pdv{\psi}{x_i}) \\
    &= -i\hbar f(\bm{r})\pdv{\psi}{x_i} - i\hbar\psi\pdv{f(\bm{r})}{x_i} + i\hbar {f}(\bm{r})\pdv{\psi}{x_i} = - i\hbar\pdv{f(\bm{r})}{x_i} \psi \quad.
\end{aligned}
\]
Con le solite assunzioni, otteniamo $[\hat{{p}}_i, \hat{f}(\bm{r})] = -i\hbar \partial_{x_i} f(\bm{r})$.

Per pura curiosità, proviamo a fare lo stesso calcolo nell'impalcatura classica:
\[
\begin{aligned}
    \{p_i, f(\bm{r})\} &= \sum_\alpha \left( \pdv{p_i}{q_\alpha}\pdv{f(\bm{r})}{p_\alpha} - \pdv{p_i}{p_\alpha}\pdv{f(\bm{r})}{q_\alpha} \right) \\
    &= \sum_\alpha \bigg(-\delta_{i\alpha}\pdv{f(\bm{r})}{q_\alpha}\bigg) = - \pdv{f(\bm{r})}{q_i} \quad,
\end{aligned}
\]
Ancora una volta, l'analogia suggerita dalla quantizzazione canonica si rivela fondata.

\subsubsection{Una componente dell'operatore momento angolare e una componente dell'operatore quantità di moto}
Continuiamo: 
\[
\begin{aligned}
    [\hat{L}_i, \hat{p}_n]\psi &= [\epsilon_{ijk}\hat{x}_j\hat{p}_k, \hat{p}_n]\psi = \epsilon_{ijk}[\hat{x}_j\hat{p}_k, \hat{p}_n] = - \epsilon_{ijk}[\hat{p}_n,\hat{x}_j\hat{p}_k]\\
    &= - \epsilon_{ijk} \big( \underbrace{[\hat{p}_n,\hat{x}_j]}_{-i\hbar \delta_{nj}}\hat{p}_k - \hat{x}_j \underbrace{[\hat{p}_n,\hat{p}_k]}_{0} \big) = i\hbar\epsilon_{ijk} \delta_{nj}\hat{p}_k \psi = i\hbar\epsilon_{ink} \hat{p}_k\psi \quad.
\end{aligned}
\]
Otteniamo così $[\hat{L}_i, \hat{p}_n] = i\hbar\epsilon_{ink}\hat{p}_k$. Non calcoliamo la controparte classica, nella certezza del fatto che il meccanismo è ormai chiaro.


\subsubsection{Una componente dell'operatore momento angolare e una componente dell'operatore posizione}
Seguendo il solito metodo:
\[
\begin{aligned}
    [\hat{L}_i, \hat{x}_n]\psi &= [\epsilon_{ijk}\hat{x}_j\hat{p}_k, \hat{x}_n]\psi = \epsilon_{ijk}[\hat{x}_j\hat{p}_k, \hat{x}_n] = - \epsilon_{ijk}[\hat{x}_n,\hat{x}_j\hat{p}_k]\\
    &= - \epsilon_{ijk} \big( \underbrace{[\hat{x}_n,\hat{x}_j]}_{0}\hat{p}_k - \hat{x}_j \underbrace{[\hat{x}_n,\hat{p}_k]}_{i\hbar \delta_{nk}} \big) = i\hbar\epsilon_{ijk} \delta_{nk}\hat{x}_j \psi = i\hbar\epsilon_{ijn}\hat{x}_j\psi \quad.
\end{aligned}
\]
Abbiamo quindi $[\hat{L}_i, \hat{x}_n] = i\hbar\epsilon_{ijn}\hat{x}_j$.

\subsubsection{Una componente dell'operatore momento angolare e l'operatore quantità di moto quadra}
Prima di imbastire i calcoli, si ricordi che l'operatore quantità di moto al quadrato è uno scalare definito come la somma dei quadrati delle sue componenti:
\[
    \hat{\bm{p}}^2 = \sum_k \hat{p}_k^2 \quad.
\]
Svolgendo i calcoli tramite la regola di Leibniz per il commutatore, abbiamo:
\[
\begin{aligned}
    [\hat{L}_i, \hat{\bm{p}}^2]\psi &= \sum_k [\hat{L}_i , \hat{p}_k^2]\psi = \sum_k [\hat{L}_i , \hat{p}_k \hat{p}_k]\psi \\
    &= \sum_k \biggl([\hat{L}_i , \hat{p}_k] \hat{p}_k + \hat{p}_k [\hat{L}_i , \hat{p}_k] \biggr)\psi \\
    &\overset{(1)}{=} \sum_{k} \bigl( i\hbar \epsilon_{iks}\hat{p}_s\hat{p}_k + \hat{p}_k i\hbar \epsilon_{iks}\hat{p}_s \bigr)\psi \\
    &= i\hbar \sum_{k,s} \epsilon_{iks} (\hat{p}_s\hat{p}_k + \hat{p}_k\hat{p}_s)\psi \quad.
\end{aligned}
\]
In $(1)$ abbiamo esplicitato il commutatore $[\hat{L}_i, \hat{p}_k] = i\hbar \epsilon_{iks}\hat{p}_s$ calcolato precedentemente.

Arrivati a questo punto, per capire perché l'intera espressione si annulla, analizziamo i termini della doppia sommatoria. Sappiamo che le diverse componenti dell'impulso commutano fra loro ($[\hat{p}_s, \hat{p}_k] = 0$), quindi l'ordine in cui applichiamo gli operatori è ininfluente: $\hat{p}_s\hat{p}_k = \hat{p}_k\hat{p}_s$. La quantità in parentesi si riduce perciò a $2\hat{p}_s\hat{p}_k$. La nostra somma diventa:
\[
    [\hat{L}_i, \hat{\bm{p}}^2]\psi = 2i\hbar \sum_{k,s} \epsilon_{iks} \hat{p}_s\hat{p}_k \psi \quad.
\]
Pensiamo ora a cosa succede quando esplicitiamo questa somma su tutti i possibili valori di $k$ ed $s$ (che spaziano sugli assi $1, 2, 3$):
\begin{itemize}
    \item Se prendiamo i termini con $k = s$, il simbolo di Levi-Civita $\epsilon_{ikk}$ è pari a zero per definizione. Dunque questi termini si annullano immediatamente.
    \item Se prendiamo i termini con $k \neq s$, ci accorgiamo che per ogni specifica coppia (ad esempio $k=1, s=2$), esisterà all'interno della somma anche il suo termine "gemello" con gli indici scambiati ($k=2, s=1$).
\end{itemize}

Confrontiamo questi due termini gemelli. Da un lato, il prodotto degli operatori è identico per via della commutatività ($\hat{p}_1\hat{p}_2 = \hat{p}_2\hat{p}_1$). Dall'altro lato, la regola di permutazione del simbolo di Levi-Civita impone un cambio di segno quando si scambiano due indici: $\epsilon_{i12} = -\epsilon_{i21}$.
Questo significa che il termine corrispondente alla coppia $(1,2)$ sarà esattamente uguale e opposto al termine $(2,1)$. Sommandoli nella sommatoria, essi si cancelleranno a vicenda:
\[
    \epsilon_{i12} \hat{p}_1\hat{p}_2 + \epsilon_{i21} \hat{p}_2\hat{p}_1 = \epsilon_{i12} \hat{p}_1\hat{p}_2 - \epsilon_{i12} \hat{p}_1\hat{p}_2 = 0 \quad.
\]
Poiché questo ragionamento vale per \emph{qualsiasi} coppia di indici incrociati nella sommatoria, tutti i termini si elidono a due a due. L'intero risultato collassa quindi a zero:
\begin{equation*}
    [\hat{L}_i, \hat{\bm{p}}^2] = 0 \quad.
\end{equation*}

Fisicamente, ricordando che l'operatore energia cinetica $\hat{K}$ per una particella è proporzionale a $\hat{\bm{p}}^2$, questo risultato garantisce che le componenti del momento angolare commutano con l'energia cinetica. Questa proprietà si rivelerà centrale per lo studio dei sistemi a simmetria sferica, in cui il momento angolare costituisce una costante del moto.
\chapter{L'evoluzione dell'informazione quantistica}
Nei capitoli precedenti abbiamo eretto le fondamenta assiomatiche e algebriche della teoria quantistica. Abbiamo compreso come la massima informazione accessibile su un sistema sia codificata all'interno di una funzione d'onda e come le grandezze fisiche della meccanica analitica debbano essere promosse a operatori hermitiani, governati dalle stringenti regole della quantizzazione canonica e dalla non-commutatività. 

Tuttavia, il formalismo finora sviluppato definisce unicamente la bsae cinematica della teoria. Per descrivere un sistema fisico concreto ed estrarre previsioni quantitative misurabili in laboratorio, è necessario concentrare l'attenzione sull'operatore principe che governa l'intero sistema: l'\textbf{operatore Hamiltoniano} $\hat{\m{H}}$.

L'equazione di Schrödinger, discussa in precedenza nelle sue vesti fenomenologiche, assume ora il suo ruolo strutturale definitivo. Riducendo la trattazione ai sistemi conservativi, lo studio della dinamica si trasformerà in un \emph{problema agli autovalori} per l'energia. Risolvere l'equazione di Schrödinger equivale, in un certo senso, a "diagonalizzare" l'Hamiltoniana e al ricercare quindi i suoi \textbf{autovalori} e i suoi \emph{autovettori} che, nel mondo quantistico, diventano \textbf{autostati}.

Questo capitolo segna dunque il passaggio dall'astrazione algebrica alla prassi differenziale. Analizzeremo l'equazione di Schrödinger per diverse configurazioni spaziali dell'energia potenziale $V(\bm{r})$, partendo dai casi unidimensionali analiticamente risolubili: la particella libera, i gradini, le barriere e le buche di potenziale.

Il risultato concettualmente più profondo che emergerà da questa trattazione è la natura stessa del concetto di "quanto". Come vedremo, la \textbf{quantizzazione dell'energia} - storicamente introdotta tramite postulati \emph{ad hoc} nei modelli semiclassici - scaturirà qui in modo del tutto spontaneo e rigoroso. Essa non sarà un assioma imposto a priori, bensì una diretta e inevitabile conseguenza analitica della richiesta che le funzioni d'onda siano fisicamente ammissibili e che rispettino le \emph{condizioni al contorno} topologiche e di regolarità imposte dai vari potenziali fisici.

\section{L'evoluzione temporale della funzione d'onda}
In un sistema conservativo, nel quale l'hamiltoniana si conserva\footnote{Si ricordi che $H = E$.} e l'energia potenziale è tempo-indipendente, è possibile risolvere la S.E. assumendo che $\Psi(\bm{r},t)$ che descrive lo stato analizzato sia scomponibile nel prodotto di una funzione dipendente \textbf{esclusivamente dal tempo} $f(t)$ e di una dipendente \textbf{esclusivamente dallo spazio} $\psi(\bm{r})$:
\begin{equation}
    \Psi(\bm{r},t) = f(t)\psi(\bm{r}) \quad.
    \label{eq:disaccop}
\end{equation}
Scriviamo la S.E.:
\[
\begin{aligned}
    i\hbar \pdv{\Psi(\bm{r},t)}{t} = \hat{\m{H}}\Psi(\bm{r},t) \implies
    i\hbar \pdv{f(t)\psi(\bm{r})}{t} &= \hat{\m{H}}\big(f(t)\psi(\bm{r})\big) \\
    i\hbar \psi(\bm{r})\dv{f(t)}{t} &= f(t)\hat{\m{H}}\psi(\bm{r}) \quad.
\end{aligned}
\]
Isolando i termini dipendenti da spazio e tempo rispettivamente a membro destro e sinistro dell'uguaglianza, otteniamo:
\[
\frac{i\hbar}{f(t)}\dv{f(t)}{t} = \frac{\hat{\m{H}}\psi(\bm{r})}{\psi{(\bm{r})}} \quad.
\] 
Abbiamo \emph{disaccoppiato} la S.E., trasformandola in un'equazione che presenta ai suoi membri oggetti dipendenti da variabili fra loro indipendenti. In questo senso possiamo giustificare, almeno euristicamente, l'ansatz formalizzato nella \eqref{eq:disaccop}: i due membri della S.E. appena riformulata dipendono da variabili \textbf{indipendenti} quali, appunto, spazio $\bm{r}$ e tempo $t$. L'unico modo che abbiamo per risolvere l'equazione è pertanto quello di imporre un'ulteriore uguaglianza ad una costante comune $E$: 
\[
\frac{i\hbar}{f(t)}\dv{f(t)}{t} = \frac{\hat{\m{H}}\psi(\bm{r})}{\psi{(\bm{r})}} = E \implies 
\begin{dcases}
    \frac{i\hbar}{f(t)}\dv{f(t)}{t} = E \\
    {\hat{\m{H}}\psi(\bm{r})} = E{\psi{(\bm{r})}} \quad.
\end{dcases}
\] 
La soluzione della parte temporale porta a:
\begin{equation}
    f(t) = Ce^{-i\frac{t}{\hbar}E}, \quad C \in \mathbb{C} \quad.
\end{equation}
La parte spaziale, invece, corrisponde ad una - non tanto classica\footnote{Si noti che la presenza dell'operatore differenziale energia cinetica in $\hat{\m{H}}$, costringe a identificare l'equazione agli autovalori in questione come un'equazione differenziale del secondo ordine alle derivate parziali.} - equazione agli autovalori per l'operatore $\hat{\m{H}}$ in cui, famosa ai più come \textbf{equazione di Schrödinger stazionaria}. In piena analogia con l'algebra matriciale, si cerca una $\psi(\bm{r})$ (\emph{autostato} dell'hamiltoniana) che viene lasciata invariata dall'applicazione di $\hat{\m{H}}$, a meno di uno scalamento dato da un fattore costante $E$, che in questo contesto assume il significato di \textbf{autovalore dell'energia}. 

La generica soluzione dell'equazione è quindi data da:
\begin{equation}
    \Psi_E(\bm{r},t) = f(t)\psi(\bm{r}) = Ce^{-i\frac{t}{\hbar}E}\psi_E(\bm{r}) 
\end{equation} 
dove è stato aggiunto il pedice $E$ per evidenziare che l'autostato $\psi(\bm{r})$ è proprio quello associato all'autovalore $E$ e dare così coerenza alla notazione. Si noti che $\Psi$ è conosciuta come \textbf{stato stazionario} o \emph{soluzione stazionaria} della S.E. 
L'aggettivo \emph{stazionario} deriva dal fatto che la norma quadra di $\Psi_E$, ovvero la densità di probabilità $\rho$, è indipendente dal tempo. Infatti:
\[
    |\Psi_E(\bm{r},t)|^2 = |C|^2|\psi(\bm{r})|^2 = \rho \quad,
\]
essendo unitario il modulo quadro di un esponenziale complesso.

\subsection{Soluzione generica tramite principio di sovrapposizione}
La linearità della S.E. consente, come visto, di sfruttare il principio di sovrapposizione e ottenere che la più generica soluzione dell'equazione prende la forma di:
\begin{equation}
    \Psi(\bm{r},t) = \sum_n \Psi_{E_n}(\bm{r},t) = \sum_n C(E_n)e^{-i\frac{t}{\hbar}E_n}\psi_{E_n}(\bm{r}) \quad,
    \label{eq:sovrapp_discr}
\end{equation}
nel caso di una distribuzione discreta dei livelli energetici, oppure di:
\begin{equation}
    \Psi(\bm{r},t) = \int_\Omega \dd{E} \Psi_{E}(\bm{r},t) = \int_\Omega \dd{E} C(E)e^{-i\frac{t}{\hbar}E}\psi_{E}(\bm{r}) \quad,
\end{equation}
nel caso di distribuzioni continue di energia. Con $\Omega$ abbiamo indicato il dominio di definizione di $\psi$.

Abbiamo così, in maniera relativamente semplice, scoperto un fatto fondamentale: un generico stato che descrive un sistema non è altro che la combinazione lineare del prodotto di "stati elementari" del sistema stesso, dipendenti dagli autostati dell'hamiltoniana e da una parte temporale. Chissà che non si possa definire una sorta di \emph{base} tramite cui scrivere un qualsiasi stato $\psi$...

\subsubsection{Check della soluzione generale}
Controlliamo, per rigore formale, che la soluzione generale trovata risponda effettivamente ai requeisiti della S.E. Calcoliamo innanzitutto la derivata parziale rispetto al tempo e moltiplichiamola per il fattore $i\hbar$:
\[
i\hbar\pdv{\Psi(\bm{r},t)}{t} \overset{(1)}{=} i\hbar\sum_n C(E_n) \pdv{\bigg(e^{-i\frac{t}{\hbar}E_n}\bigg)}{t} \psi_{E_n}(\bm{r}) = \sum_n C(E_n) E_n e^{-i\frac{t}{\hbar}E_n}\psi_{E_n}(\bm{r}) \quad,
\]
dove in $(1)$ abbiamo sfruttato la linearità della derivata. Il membro destro della S.E. lo otteniamo invece da:
\[
\begin{aligned}
    \hat{\m{H}} \Psi(\bm{r},t) = \hat{\m{H}}\bigg[ \sum_n C(E_n) {e^{-i\frac{t}{\hbar}E_n}} \psi_{E_n}(\bm{r}) \bigg] &= \sum_n C(E_n) {e^{-i\frac{t}{\hbar}E_n}} \hat{\m{H}}\psi_{E_n}(\bm{r}) \\
    &=  \sum_n C(E_n)E_n {e^{-i\frac{t}{\hbar}E_n}} \psi_{E_n}(\bm{r}) \quad,
\end{aligned}
\]
evidentemente uguale a quanto ottenuto per la derivata temporale. Questo chiude il nostro check confermando che la soluzione generica trovata è assolutamente valida.

Non viene proposto il conto per distribuzioni continue di energia essendo formalmente identico.

\subsection{La rappresentazione di Feynman: l'operatore di evoluzione temporale}
Possiamo stabilire un ennesimo ponte con la meccanica classica. Nella formulazione di Hamilton, l'\emph{operatore di evoluzione temporale} è legato all'hamiltoniana del sistema essendo quest'ultima la funzione \emph{generatrice di traslazioni nel tempo}. Supponiamo ora che in meccanica quantistica valga lo stesso e definiamo, in piena analogia con la teoria classica, l'\textbf{operatore di evoluzione temporale}
\begin{equation}
    \hat{\m{U}}(t) = e^{-i\frac{\hat{\m{H}}}{\hbar}t} \quad.
\end{equation} 
Tramitev $\hat{\m{U}}$ possiamo scoprire l'evoluzione temporale del generico stato $\Psi$ conoscendo \emph{unicamente} $\Psi(\bm{r}, t = 0)$. 

Per comprendere come sia possibile, valutiamo innanzitutto $\Psi(\bm{r}, 0)$ nel caso di una distribuzione discreta di energia:
\[
    \Psi(\bm{r},0) = \sum_n C(E_n)e^{-i\frac{t = 0}{\hbar}E_n}\psi_{E_n}(\bm{r}) = \sum_n C(E_n)\psi_{E_n}\psi_{E_n}(\bm{r}) \quad.
\]
Sappiamo inoltre che, con sufficienti condizioni di regolarità, una funzione che prende in argomento un operatore $f(\hat{\m{A}})$, può essere espansa tramite \emph{serie Taylor}. Espandiamo quindi $\hat{\m{U}}$\footnote{Ricordando che $e^z = \sum_{s = 0}^{\infty} \frac{z^s}{s!}, \quad z \in \mathbb{C}$.}:
\[
    e^{-i\frac{\hat{\m{H}}}{\hbar}t} = \sum_{s=0}^{\infty} \frac{1}{s!}\left(-i\frac{\hat{\m{H}}}{\hbar}t\right)^s = \sum_{s=0}^{\infty} \frac{1}{s!}\left(\frac{-it}{\hbar}\right)^s \hat{\m{H}}^s
\]
e applichiamolo a $\Psi(\bm{r},0)$:
\[
\begin{aligned}
    \hat{\m{U}}\Psi(\bm{r},0) = \hat{\m{U}}\sum_n C(E_n)\psi_{E_n}(\bm{r}) &\overset{(1)}{=} \sum_n C(E_n) \hat{\m{U}} \psi_{E_n}(\bm{r}) \\
    &= \sum_n C(E_n) \sum_{s} \frac{1}{s!}\left(\frac{-it}{\hbar}\right)^s \underbrace{\hat{\m{H}}^s \psi_{E_n}(\bm{r})}_{(a)} \quad,
\end{aligned}
\]
dove in $(1)$ abbiamo sfruttato la linearità dell'operatore $\hat{\m{U}}$ evidente dall'espansione in Taylor. Ci rimane da calcolare il termine $(a)$ e, per farlo, iteriamo ripetutamente l'applicazione dell'operatore hamiltoniano su $\psi_{E_n}$:
\[
\begin{aligned}
    &\hat{\m{H}}\psi_{E_n}(\bm{r}) = E_n\psi_{E_n}(\bm{r}) \\
    &\hat{\m{H}}^2\psi_{E_n}(\bm{r}) = \hat{\m{H}}(\hat{\m{H}}\psi_{E_n}(\bm{r})) = E_n\hat{H}\psi_{E_n}(\bm{r}) = E_n^2\psi_{E_n}(\bm{r}) \\
    &\hat{\m{H}}^3\psi_{E_n}(\bm{r}) = \hat{\m{H}}(\hat{\m{H}}^2\psi_{E_n}(\bm{r})) = E_n^2\hat{\m{H}}\psi_{E_n}(\bm{r}) = E_n^3 \psi_{E_n}(\bm{r}) \\
    &  \quad \vdots \\
    & \hat{\m{H}}^s \psi_{E_n}(\bm{r}) = E_n^s\psi_{E_n}(\bm{r}) \quad.
\end{aligned}
\]
Sostituendo nell'espressione per $\hat{\m{U}}\Psi(\bm{r},0)$,
\[
\begin{aligned}
    \hat{\m{U}}\Psi(\bm{r},0) &= \sum_n C(E_n) \sum_{s} \frac{1}{s!}\left(\frac{-it}{\hbar}\right)^s E_n^s\psi_{E_n}(\bm{r}) = \sum_n C(E_n) \underbrace{\sum_{s} \frac{1}{s!}\left(-i\frac{E_n}{\hbar}t\right)^s}_{e^{-i\frac{E_n}{\hbar}t}}\psi_{E_n}(\bm{r}) \\
    \hat{\m{U}}\Psi(\bm{r},0) &= \sum_n C(E_n)e^{-i\frac{t}{\hbar}E_n}\psi_{E_n}(\bm{r}) = \Psi(\bm{r}, t)\quad,
\end{aligned}
\]
otteniamo la stessa espressione della \eqref{eq:sovrapp_discr}, confermando la bontà del metodo e dell'analogia classica. 

\begin{approfondimento}[title={Altro sull'operatore di evoluzione temporale}]
    \noindent
    Sino ad ora ci siamo concentrati sul valutare cosa restituisse l'applicazione di $\hat{\m{U}}$ sul generico stato $\Psi$ valutato per $t = 0$. 

    Se però applicassimo l'operatore di evoluzione temporale a un generico \emph{autostato}? Facciamo il calcolo.
    \[
    \begin{aligned}
        \hat{\m{U}}\psi_{E_n}(\bm{r}) = e^{-i\frac{\hat{\m{H}}}{\hbar}t}\psi_{E_n}(\bm{r}) &= \sum_s \frac{1}{s!}\left(-i\frac{\hat{\m{H}}}{\hbar}t\right)^s \psi_{E_n}(\bm{r}) = \sum_s \frac{1}{s!}\left(\frac{-it}{\hbar}\right)^s \hat{\m{H}}^s\psi_{E_n}(\bm{r})\\
        &= \sum_s \frac{1}{s!}\left(-i\frac{{E_n}}{\hbar}t\right)^s \psi_{E_n}(\bm{r}) = e^{-i\frac{E_n}{\hbar}t}\psi_{E_n}(\bm{r}) \quad.
    \end{aligned}
    \]
    Come si poteva prevdere dal conto condotto in precedenza, l'applicazione dell'operatore di evoluzione temporale a un generico autostato restituisce lo \textbf{stato stazionario} associato all'autostato stesso.
\end{approfondimento}

\section{Il caso della particella libera}

\subsubsection{Una breve introduzione}
Siamo finalmente pronti a vedere la prima concreta applicazione della teoria vista sino ad ora: il caso della \emph{particella libera}. Classicamente, lo studio dell'argomento in esame è semplice e la meccanica newtoniana è sufficiente per avere abbastanza informazioni sul sistema. Ancor prima di complicarsi la vita con equazioni di Lagrange, Eulero-Lagrange, Hamilton \& Co., il \emph{secondo principio di Newton} ci permette di codificare tutto in maniera elegantissima: detta infatti $m$ la massa del corpuscolo, abbiamo la monolitica
\[
    \bm{F} = \dot{\bm{p}} = m\bm{a} \quad.
\]
Si presti attenzione a non sottovalutare la teoria newtoniana in quanto nasconde un sistema di equazioni differenziali del secondo ordine da cui si può estrarre tutto lo scibile sul sistema:
\[
    \bm{F} = m\bm{a} \implies 
    \begin{cases}
        m\ddot{x}_1 = \dot{\bm{p}} \cdot \bm{u}_{x_1} \\
        m\ddot{x}_2 = \dot{\bm{p}} \cdot \bm{u}_{x_2} \\
        m\ddot{x}_3 = \dot{\bm{p}} \cdot \bm{u}_{x_3} \quad,
    \end{cases}
\]
che si riducono a 
\[
\bm{p} = \text{costante}
\]
nel caso in cui $\bm{F} = \bm{0}$, ovvero nella situazione in cui la particella è \emph{libera}. Con il termine "libera" si intende infatti l'assenza di campi di forze nel sistema analizzato.

\subsection{La quantistica entra in gioco}
Quantisticamente, la situazione è assai più complessa. Innanzitutto, se la perdita di determinazione che discutevamo poc'anzi è vera, dovremmo avere la possibilità di constatarla anche in questo contesto. È così sarà. Facciamo però un passo alla volta. 

L'assenza di campi di forze che influenzano la particella impone per l'hamiltoniana la forma $\hat{\mathcal{H}} = \hat{\mathcal{K}}$. L'equazione di Schrödinger si scrive:
\[
    i\hbar\pdv{\psi(\bm{r},t)}{t} = \hat{\m{H}}\psi(\bm{r},t) = \hat{\m{K}}\psi(\bm{r},t) = \frac{\hat{\bm{p}}^2}{2m}\psi(\bm{r},t) = -\frac{\hbar^2}{2m}\nabla^2\psi(\bm{r},t) \quad.
\]
Per risolvere l'equazione, riproponiamo lo stesso schema risolutivo utilizzato in precedenza per risolvere la S.E. in un caso generale. Ipotizziamo per la generica soluzione $\Phi(\bm{r},t)$ la forma:
\[
    \Phi(\bm{r},t) = f(t)\phi(\bm{r}) \quad. 
\]
Sostituendo nell'equazione otteniamo ancora una forma esponenziale complessa per la parte temporale, $f(t) = e^{-i\frac{E}{\hbar}t}$, e la solita equazione agli autovalori, questa volta per l'operatore $\h{K}$:
\begin{equation}
    \h{K}\phi(\bm{r}) = E\phi(\bm{r}) \quad.
    \label{eq:eigen_K}
\end{equation}
Per dare forma agli autostati dell'operatore energia cinetica, lo riscriviamo nella forma più conveniente:
\[
    \h{K}\phi(\bm{r}) = -\frac{\hbar^2}{2m}\nabla^2\phi(\bm{r}) = E\phi(\bm{r}) \quad.
\]
Data la simmetria del problema, possiamo semplificare i calcoli e ricondurci a un'equazione differenziale del secondo ordine \emph{unidimensionale}, per poi estendere il risultato a procedimento ultimato. La S.E. stazionaria diventa:
\[
    \dv[2]{\phi(x_1)}{x_1} = -\frac{2mE}{\hbar^2}\phi(x_1) \quad.
\]
Le radici del polinomio caratteristico associato sono:
\[
    \lambda_{1,2} = \pm i\sqrt{\frac{2mE}{\hbar^2}}
\]
e la generica soluzione dell'equazione è:
\[
    \phi(x_1) = Ae^{x_1\lambda_1} + Be^{x_1\lambda_2}, \quad A,B \in \mathbb{C} \quad.
\]
Conviene a questo punto introdurre una grandezza fondamentale che snellirà la notazione e ci fornirà un'interpretazione fisica immediata. Definiamo il \textbf{numero d'onda} $k$ a partire dalla relazione classica fra energia cinetica e quantità di moto ($E = p^2/2m$), sfruttando la relazione di De Broglie $p = \hbar k$:
\begin{equation}
    k = \frac{p}{\hbar} = \sqrt{\frac{2mE}{\hbar^2}} \quad.
\end{equation}

Si noti che l'energia $E$ deve essere necessariamente non negativa ($E \ge 0$). Un'energia negativa porterebbe infatti a un $k$ immaginario puro e, conseguentemente, a soluzioni esponenziali reali divergenti per $x_1 \to \pm\infty$, fisicamente inaccettabili in quanto non normalizzabili in alcun senso.

Sostituendo $k$ ed eliminando per comodità tipografica il pedice $1$ dalla coordinata spaziale (ponendo $x_1 = x$), gli autostati dell'operatore $\h{K}$ assumono la forma:
\[
    \phi_k(x) = A e^{ikx} + B e^{-ikx} \quad.
\]
Imponendo per semplicità $B = 0$ e estendendo il risultato a tre dimensioni, otteniamo il generico autostato dell'operatore energia cinetica:
\[
    \phi_{\bm{k}}(\bm{r}) = Ae^{i\bm{k}\cdot \bm{r}}\quad,
\]
con $\bm{k}$ vettore d'onda. Per quanto riguarda gli autovalori, la computazione è semplice e richiede unicamente di applicare $\h{K}$ all'autostato appena trovato:
\[
    \h{K} \phi_{\bm{k}}(\bm{r}) = -\frac{\hbar^2}{2m}\nabla^2 \big(Ae^{i\bm{k}\cdot \bm{r}}\big) = -\frac{\hbar^2}{2m}  (i\bm{k})\cdot(i\bm{k})Ae^{i\bm{k}\cdot \bm{r}} = \underbrace{\frac{\hbar^2k^2}{2m}}_{\text{autovalori}}\phi_{\bm{k}}(\bm{r}) \quad.
\]

\subsubsection{La costante d'integrazione $\bm{A}$}
Comunque sia fatta $A$, siamo di fronte ad un apparente disastro matematico: gli autostati trovati non appartengono allo spazio delle funzioni a quadrato sommabile! Infatti:
\[
    \int_{\mathbb{R}^3} \dd^3{\bm{r}}\, |\phi_{\bm{k}}(\bm{r})|^2 =  \int_{\mathbb{R}^3} \dd^3{\bm{r}}\, |A|^2 \to \infty \quad.
\] 
La motivazione per questo bizzarro ritrovamento sarà chiara più avanti, ma nel mentre dobbiamo concentrarci sul tentare di dare una nuova nozione di normalizzazione per trovare $A$, perché quella "classica" evidentemente non funziona. 

L'idea geniale è da attribuirsi a Dirac: dalla discussione sulle proprietà degli operatori fisici, come $\h{K}$, abbiamo scoperto che gli operatori hermitiani ammettono un set di autostati ortogonali. Se per gli spettri discreti l'ortonormalità è garantita dalla Delta di Kronecker (ovvero $(\phi_i, \phi_j) = \delta_{ij}$), per uno \emph{spettro continuo} come quello del vettore d'onda $\bm{k}$, questa condizione si generalizza utilizzando la \textbf{Delta di Dirac}. Proviamo a imporre questa condizione, detta \emph{normalizzazione nel senso delle distribuzioni}, su due generici autostati $\phi_{\bm{k}}$ e $\phi_{\bm{k}'}$:
\[
    (\phi_{\bm{k}}, \phi_{\bm{k}'}) = \delta^{(3)}(\bm{k} - \bm{k}') \quad.
\]
Sostituendo le espressioni delle onde piane e sviluppando il prodotto scalare:
\[
    (\phi_{\bm{k}}, \phi_{\bm{k}'}) = \int_{\mathbb{R}^3} \dd^3{\bm{r}}\, \phi_{\bm{k}}^* \phi_{\bm{k}'} = \int_{\mathbb{R}^3} \dd^3{\bm{r}}\, |A|^2 e^{-i\bm{k}\cdot \bm{r}} e^{i\bm{k}'\cdot \bm{r}} = |A|^2 \int_{\mathbb{R}^3} \dd^3{\bm{r}}\, e^{i(\bm{k}'-\bm{k})\cdot \bm{r}} \quad.
\]
Dall'analisi matematica e dalla teoria delle trasformate di Fourier, riconosciamo immediatamente che l'ultimo integrale coincide con la rappresentazione integrale esatta della Delta di Dirac tridimensionale, a meno di un fattore moltiplicativo. Abbiamo infatti:
\[
    \int_{\mathbb{R}^3} \dd^3{\bm{r}}\, e^{i(\bm{k}'-\bm{k})\cdot \bm{r}} = (2\pi)^3 \delta^{(3)}(\bm{k}' - \bm{k}) = (2\pi)^3 \delta^{(3)}(\bm{k} - \bm{k}') \quad,
\]
avendo sfruttato la parità della funzione Delta. Uguagliando questo risultato alla condizione di normalizzazione iniziale, otteniamo:
\[
    |A|^2 (2\pi)^3 \delta^{(3)}(\bm{k} - \bm{k}') = \delta^{(3)}(\bm{k} - \bm{k}') \quad.
\]
Da cui ricaviamo il valore del modulo quadro della costante:
\[
    |A|^2 = \frac{1}{(2\pi)^3} \implies A = \frac{1}{(2\pi)^{3/2}} e^{i\gamma} \quad.
\]
Scegliendo convenzionalmente nulla la fase globale ($\gamma = 0$), giungiamo all'espressione definitiva per le autofunzioni normalizzate nello spazio tridimensionale:
\begin{equation}
    \phi_{\bm{k}}(\bm{r}) = \frac{1}{(2\pi)^{3/2}} e^{i\bm{k}\cdot \bm{r}} \quad.
\end{equation}

\subsection{Il generico stato stazionario}
Componendo il generico stato stazionario mediante moltiplicazione della parte temporale e spaziale appena trovate si ottiene:
\begin{equation}
    \Phi_{\bm{k}}(\bm{r},t) = \phi_{\bm{k}}(\bm{r})f(t) \overset{(1)}{=} \frac{1}{(2\pi)^{3/2}} e^{i\bm{k}\cdot \bm{r}}  e^{-i\frac{E_{\bm{k}}}{\hbar}t} \overset{(2)}{=} \frac{1}{(2\pi)^{3/2}} e^{i(\bm{k} \cdot \bm{r}-\omega_{\bm{k}} t)} \quad,
    \label{eq:onda-Piana_lib}
\end{equation}
dove in $(1)$ abbiamo aggiunto il pedice $\bm{k}$ all'energia per specificare il legame con l'autovalore $\propto |\bm{k}|^2 = k^2$ di $\h{K}$ e in $(2)$ abbiamo sfuttato la relazione di Planck $E_{\bm{k}} = \hbar \omega_{\bm{k}}$. L'interpretazione della forma appena trovata è cristallina: una particella libera è descritta, quantisticamente, da un'\textbf{onda piana monocromatica}. Possiamo anche suggerire una motivazione per la scelta di porre $B = 0$: il termine $A$ era legato alla radice positiva del polinomio caratteristico dell'equazione differenziale associata al problema degli autovalori di $\h{K}$, portando all'onda \emph{progressiva} ottenuta nella \eqref{eq:onda-Piana_lib}. Il coefficiente $B$, invece, avrebbe generato un'onda \emph{regressiva}, scomoda ai fini della trattazione\footnote{Si noti che anche scegliendo di mantenere $B \neq 0$ avremmo ottenuto $B = (2\pi)^{-3/2}$.}. In questo senso, possiamo formulare una forma ancora più \emph{completa} per il generico stato $\Phi_{\bm{k}}(\bm{r},t)$, che includa anche il termine associato a $B$:
\[
    \Phi_{\bm{k}}(\bm{r},t) = \frac{1}{(2\pi)^{3/2}} \left[ e^{i(\bm{k} \cdot \bm{r}-\omega_{\bm{k}} t)} + e^{-i(\bm{k} \cdot \bm{r}-\omega_{\bm{k}} t)} \right]
\]
che corrisponde, come prevedibile, alla combinazione lineare di due onde, una progressiva e una regressiva. 

\begin{approfondimento}[title={Non solo questione di comodità}]
    \noindent
    Anche se più completo, questo modo di scrivere lo stato legato all'autovalore $\propto k^2$ è sconveniente perché non permette di avere l'autostato $\phi_{\bm{k}}(\bm{r})$, tramite cui abbiamo composto $\Phi_{\bm{k}}(\bm{r},t)$, come un autostato dell'operatore quantità di moto. Sfruttando un fatto scoperto in precedenza, ovvero $\phi_{\bm{k}}(\bm{r}) = \Phi_{\bm{k}}(\bm{r}, 0)$, applichiamo $\hat{\bm{p}}$ e otteniamo:
    \[
        \hat{\bm{p}}\phi_{\bm{k}} = -i\hbar \nabla \phi_{\bm{k}}(\bm{r}) = \frac{-i\hbar}{(2\pi)^{3/2}} \nabla\left[ e^{i\bm{k} \cdot \bm{r}} + e^{-i\bm{k} \cdot \bm{r}}\right] = \frac{\hbar\bm{k} }{(2\pi)^{3/2}} \left(e^{i\bm{k} \cdot \bm{r}} - e^{-i\bm{k} \cdot \bm{r}} \right) \quad,
    \]
    evidentemente non proporzionale a $\phi_{\bm{k}}(\bm{r})$ e quindi non autostato di $\hat{\bm{p}}$. L'unico modo per ricavare un autostato è, come fatto nella trattazione, \emph{eliminare} un esponenziale, ovvero porre uguale a zero il coefficiente a lui associato. Nel tal caso si otterrebbe:
    \[
        \hat{\bm{p}}\phi_{\bm{k}} = \pm \hbar \bm{k}\phi_{\bm{k}} 
    \]
    in base alla scelta di quale esponenziale far sopravvivere. 

    Nel continuo della trattazione sarà lampante il fatto che ottenere autostati "comuni" a più operatori sia qualcosa di formidabile. In questo senso, porre $B = 0$, preserva il fatto che $\hat{\bm{p}}$ e $\h{K}$ preservino \emph{lo stesso set di autostati}.
\end{approfondimento}

\subsection{La soluzione generale del problema della particella libera}
Possiamo costruire un qualsiasi stato per un corpuscolo non sottoposto a un campo di forze e in moto nell'intero spazio tridimensionale invocando il principio di sovrapposizione. In questo caso, però, bisogna notare che lo spettro dei vari valori che può assumere $k$ è \emph{continuo} e, di conseguenza, bisogna utilizzare la rappresentazione integrale:
\begin{equation}
    \Phi(\bm{r}, t) = \int_{\mathbb{R}^3} \dd^3{\bm{k}} \, a(\bm{k})\Phi_{\bm{k}}(\bm{r},t) = \int_{\mathbb{R}^3} \dd^3{\bm{k}} \, \frac{a(\bm{k})}{(2\pi)^{3/2}} e^{i(\bm{k} \cdot \bm{r}-\omega_k t)} \quad,
    \label{eq:phi_four}
\end{equation}
dove $a(\bm{k})$ rappresenta una sorta di \emph{coefficiente di combinazione lineare nel continuo}. Quanto appena trovato non è altro che la \textbf{Trasformata di Fourier} di $\Phi(\bm{r}, t)$.

\subsubsection{Check della soluzione}
Come sempre fatto, verifichiamo che la \eqref{eq:phi_four} soddisfi alla S.E. di partenza. Iniziamo calcolando il membro sinistro:
\[
\begin{aligned}
    i\hbar \pdv{\Phi(\bm{r}, t)}{t} &= i\hbar \pdv{}{t} \int_{\mathbb{R}^3} \dd^3{\bm{k}} \, \frac{a(\bm{k})}{(2\pi)^{3/2}} e^{i(\bm{k} \cdot \bm{r}-\omega_k t)} \overset{(1)}{=} i\hbar \int_{\mathbb{R}^3} \dd^3{\bm{k}} \, \frac{a(\bm{k})}{(2\pi)^{3/2}} \pdv{\big( e^{i(\bm{k} \cdot \bm{r}-\omega_k t)} \big)}{t} \\
    &= \int_{\mathbb{R}^3} \dd^3{\bm{k}} \, a(\bm{k}) (i\hbar )(-i\omega_{\bm{k}}) \Phi_{\bm{k}}(\bm{r}, t) = \int_{\mathbb{R}^3} \dd^3{\bm{k}} \, a(\bm{k}) E_k \Phi_{\bm{k}}(\bm{r}, t) \quad,
\end{aligned}
\]
dove in $(1)$ abbiamo sfruttato la commutazione di derivata temporale e integrale agente su $\bm{k}$. Usiamo $E_k$ come modo compatto per scrivere l'energia associata al vettore d'onda $\bm{k}$:
\[
    E_{\bm{k}} = \frac{\hbar^2 k^2}{2m} \quad.
\]

Il membro destro risulta invece:
\[
\begin{aligned}
    \h{H}\Phi(\bm{r}, t) = \h{K}\Phi(\bm{r}, t) &= -\frac{\hbar^2}{2m}\nabla^2 \int_{\mathbb{R}^3} \dd^3{\bm{k}} \, \frac{a(\bm{k})}{(2\pi)^{3/2}} e^{i(\bm{k} \cdot \bm{r}-\omega_{\bm{k}} t)}\\
    & = \int_{\mathbb{R}^3} \dd^3{\bm{k}} \, \frac{a(\bm{k})}{(2\pi)^{3/2}} \left[ -\frac{\hbar^2}{2m}\nabla^2 \big( e^{i(\bm{k} \cdot \bm{r}-\omega_{\bm{k}} t)} \big) \right] \\
    &= \int_{\mathbb{R}^3} \dd^3{\bm{k}} \, a(\bm{k}) E_{\bm{k}} \Phi_{\bm{k}}(\bm{r}, t) \quad,
\end{aligned}
\]
esattamente coincidente con quanto calcolato in precedenza.

\subsection{Gli stati localizzati: il pacchetto d'onda gaussiano}
Per considerare definitivamente archiviata questa analisi, rimane da risolvere l'integrale di Fourier ottenuto in precedenza e l'unica forma funzionale ancora da determinare sono i coefficienti $a(\bm{k})$.

Conducendo un'analisi semiclassica, forti dei risultati di De Broglie, abbiamo che una particella libera deve avere la quantità di moto costante e pari a:
\[
    \bm{p} = \hbar \bm{q} \quad,
\]
con $\bm{q}$ il vettore d'onda associato. Assumiamo allora che la trasposizione quantistica di questo sistema sia contraddistinta dai coefficienti della trasformata di Fourier così definiti:
\begin{equation}
    a(\bm{k}) = A e^{-(\bm{k}-\bm{q})^2 d^2}, \quad A \in \mathbb{C} \quad.
\end{equation}
Questo particolare ansatz ci permette di interpretare nella maniera più "classica" possibile la situazione. I maggiori contributi alla \eqref{eq:phi_four} sono infatti dati da tutti quegli stati con quantità di moto affini a quella assunta classicamente dalla particella mentre, per valori di $\bm{k}$ molto diversi da $\bm{q}$, la gaussiana garantisce un contributo prossimo allo zero. Globalmente, si registra che gli unici termini rilevanti ai fini del generico stato $\Phi$ derivano logicamente dagli stati stazionari $\Phi_{\bm{k}}$ caratterizzati da $\bm{k} \simeq \bm{q}$. 

Sostituiamo nell'integrale:
\[
    \Phi(\bm{r}, t) = \int_{\mathbb{R}^3} \dd^3{\bm{k}} \, \frac{a(\bm{k})}{(2\pi)^{3/2}} e^{i(\bm{k} \cdot \bm{r}-\omega_{\bm{k}} t)} = \int_{\mathbb{R}^3} \dd^3{\bm{k}} \, \frac{A e^{-(\bm{k}-\bm{q})^2 d^2}}{(2\pi)^{3/2}} e^{i(\bm{k} \cdot \bm{r}-\omega_{\bm{k}} t)} \quad.
\]
Per continuare, notiamo che l'integrale esteso su $\mathbb{R}^3$ non è altro che un integrale triplo sulle tre direzioni spaziali $x_{1,2,3}$. Allegerendo la notazione passando alle componenti vettoriali, abbiamo i seguenti passaggi intermedi:
\[
    (\bm{k}-\bm{q})^2 = (k_1 - q_1)^2 + (k_2 - q_2)^2 + (k_3 - q_3)^2, \quad \bm{k} \cdot \bm{r} = k_1x_1 + k_2x_2 + k_3x_3 \quad.
\]
Sostituendo e sfruttando le proprietà degli esponenziali otteniamo:
\[
\begin{aligned}
    \Phi(\bm{r}, t) &= \frac{A}{(2\pi)^{3/2}} \int_{\mathbb{R}} \dd{{k}_1} \, \int_{\mathbb{R}} \dd{{k}_2} \, \int_{\mathbb{R}} \dd{{k}_3} \,  e^{-(\bm{k}-\bm{q})^2 d^2} e^{i(\bm{k} \cdot \bm{r}-\omega_{\bm{k}} t)} \\
    &= \frac{A}{(2\pi)^{3/2}} \int_{\mathbb{R}} \dd{{k}_1} \, \int_{\mathbb{R}} \dd{{k}_2} \, \int_{\mathbb{R}} \dd{{k}_3} \, \prod_{j=1}^{3} e^{-({k}_j-{q}_j)^2 d^2} e^{i({k}_j {x}_j-\omega_{{k}_j} t)} \quad.
    \end{aligned}
\]
Data la simmetria del problema, assumiamo ora che è possibile scrivere la funzione d'onda come il prodotto di tre componenti indipendenti, ciascuna funzione di $k_{1,2,3}$:
\[
    \begin{aligned}
        \Phi(\bm{r}, t) &= \Phi(x_1, t) \Phi(x_2, t) \Phi(x_3, t)  = \prod_{j=1}^{3} \Phi(x_j, t) \\
        &= \prod_{j=1}^{3} \frac{A_j}{\sqrt{2\pi}} \underbrace{\int_{\mathbb{R}} \dd{{k}_j} \, e^{-({k}_j-{q}_j)^2 d^2} e^{i({k}_j {x}_j-\omega_{{k}_j} t)}}_I \quad,
    \end{aligned}
\]
dove abbiamo considerato, ai fini del disaccoppiamento, $A_1 A_2 A_3 = A$. Introduciamo, per semplicità di notazione, il fattore 
\[
    s(t) = \frac{\hbar t}{2m} \implies \omega_{\bm{k}} t = \frac{E_{\bm{k}}}{\hbar} t = \frac{\hbar k^2}{2m} t = s(t) k^2 \quad.
\]
Iniziamo, con pazienza, a lavorare sull'integrale:
\[
\begin{aligned}
    I &= \int_{\mathbb{R}} \dd{{k}_j} \, e^{-({k}_j-{q}_j)^2 d^2} e^{i({k}_j {x}_j-sk_j^2)} = \int_{\mathbb{R}} \dd{{k}_j} \, \exp\big[ -(k_j^2 + q_j^2 - 2q_jk_j)d^2 + ik_jx_j - isk_j^2 \big] \\
    &= e^{-q_j^2 d^2} \int_{\mathbb{R}} \dd{{k}_j} \, \exp\big[ -k_j^2d^2  + 2q_jk_jd^2 + ik_jx_j - isk_j^2 \big] \\
    &= e^{-q_j^2 d^2} \int_{\mathbb{R}} \dd{{k}_j} \, \exp\big[ -k_j^2\underbrace{(d^2 + is)}_{\nu}  +k_j \underbrace{( 2q_jd^2 + ix_j)}_{\eta}\big] = e^{-q_j^2 d^2} \int_{\mathbb{R}} \dd{{k}_j} \, \exp\big[ - \nu k_j^2 + \eta k_j \big] \\
    &\overset{(1)}{=} e^{-q_j^2 d^2} \int_{\mathbb{R}} \dd{{k}_j} \, \exp\left\{ -\nu\bigg[ \bigg( k_j - \frac{\eta}{2\nu} \bigg)^2 - \frac{\eta^2}{4\nu^2}  \bigg] \right\} = e^{-q_j^2 d^2 + \frac{\eta^2}{4\nu}} \int_{\mathbb{R}} \dd{{k}_j} \, e^{-\nu\left( k_j - \frac{\eta}{2\nu} \right)^2} \quad.
\end{aligned}
\]
In $(1)$ si è completato il quadrato. L'integrale è adesso quanto mai affine al classico gaussiano se non fosse per lo shift complesso del massimo dell'integranda. Proviamo a fare un cambio di variabile e vediamo che succede:
\[
    y = k_j - \frac{\eta}{2\nu} \rightarrow \dd{y} = \dd{k_j}, \quad \mathbb{R} \to \mathbb{R} - \frac{\eta}{2\nu}, \quad e^{-\nu\left( k_j - \frac{\eta}{2\nu} \right)^2} \to e^{-\nu y^2} \quad.
\]
Benchè la variabile d'integrazione originale sia reale e corra lungo $\mathbb{R}$, la necessaria sostituzione che serve per risolvere il calcolo impone di trattare con una variabile complessa $y \in \mathbb{C}$ integrata lungo una retta traslata rigidamente dall'asse reale di una quantità pari al fantomatico shift... come fare? 

Immaginiamo di tracciare una sorta di rettangolo sul piano complesso contraddistinto dai vertici $A = \left( -R, 0 \right)$, $B = \left( R, 0 \right)$, $C = \left( -R, -\frac{\eta}{2\nu}\right)$ e $D = \left( R, -\frac{\eta}{2\nu}\right)$. Il lato $AB$ giace su $\mathbb{R}$, $CD$ sulla retta su cui dovremmo integrare. Essendo $e^{-\nu y^2}$, possiamo sfruttare il \emph{Teorema di Cauchy-Goursat} per concludere che l'integrale sul rettangolo è nullo. Osserviamo però che, quando si fa tendere $R \to \infty$, l'integranda tende a zero! I contributi dovuti ai "lati verticali" del rettangolo sono quindi trascurabili e otteniamo che l'integrale lungo la retta reale traslata è uguale a quello calcolato lungo $\mathbb{R}$. Possiamo quindi scrivere:
\[
I = e^{-q_j^2 d^2 + \frac{\eta^2}{4\nu}} \int_{-\infty}^{+\infty}e^{-\nu y^2} = e^{-q_j^2 d^2 + \frac{\eta^2}{4\nu}} \sqrt{\frac{\pi}{\nu}} \quad.
\]
La funzione d'onda risultante è:
\[
\begin{aligned}
    \Phi(\bm{r},t) &= \prod_{j=1}^{3} \frac{A_j}{\sqrt{2\pi}} \sqrt{\frac{\pi}{\nu}} \exp[-q_j^2d^2 + \frac{(2q_jd^2 + ix_j)^2}{4\nu}]  = \prod_{j=1}^{3} \frac{A_j}{\sqrt{2\nu}} \exp[-q_j^2d^2 + \frac{(2q_jd^2 + ix_j)^2}{4\nu}] \\
    &= \prod_{j=1}^{3} \frac{A_j}{\sqrt{2(d^2 + is)}} \exp[-q_j^2d^2 + \frac{(2q_jd^2 + ix_j)^2}{4(d^2+is)}] \quad.
\end{aligned}
\]  

\subsubsection{La densità di probabilità gaussiana}
Se all'inizio del paragrafo abbiamo introdotto in qualche modo il concetto di "localizzazione" per la funzione d'onda, richiedendo che i coefficienti dello sviluppo di Fourier tenessero conto della differenza fra il momento del genrico stato $\phi$ e quello reale, quest'ultimo emerge in maniera quanto mai significativa quando si valuta la densità di probabilità associata a $\Phi$. 

Ricordando che $\rho = |\Phi(\bm{r},t)|^2 = |\Phi(x_1,t)|^2|\Phi(x_2,t)|^2|\Phi(x_3,t)|^2$, per i nostri fini è dunque sufficienti calcolare $|\Phi(x_j,t)|^2$. Iniziamo:
\[
\begin{aligned}
    |\Phi(x_j,t)|^2 = \Phi(x_j,t)^*\Phi(x_j,t) &= \frac{|A_j|^2}{{2|d^2 + is|}}e^{-2q_j^2d^2}\exp[\frac{(2q_jd^2 + ix_j)^2}{4(d^2+is)}]\exp[\frac{(2q_jd^2 + ix_j)^2}{4(d^2+is)}]^* \\
    &= \frac{|A_j|^2}{{2\sqrt{d^4+s^2}}} e^{-2q_j^2d^2} \exp{2\Re\left[\frac{(2q_jd^2 + ix_j)^2}{4(d^2+is)}\right]} \\
    &= \frac{|A_j|^2}{{2\sqrt{d^4+s^2}}} e^{-2q_j^2d^2} \exp[2q_j^2d^2 -d^2\frac{(2sq_j-x_j)^2}{2(d^4+s^2)}] \\
    &= \frac{|A_j|^2}{{2\sqrt{d^4+s^2}}} \exp[-d^2\frac{(2sq_j-x_j)^2}{2(d^4+s^2)}] \quad.
\end{aligned}
\]
Ci siamo quasi. Espandiamo ancora $2sq_j$:
\[
    2sq_j = 2\left(\frac{\hbar t}{2m}\right)q_j \overset{(1)}{=} v_j t \quad,
\]
dove in $(1)$ abbiamo usato la relazione \textbf{classica} per la quantità di moto $q_j = mv_j$. Sostituendo nella densità di probabilità otteniamo un'espressione definitiva:
\begin{equation}
    \rho_{x_j} = |\Phi(x_j,t)|^2 = \frac{|A_j|^2}{{2\sqrt{d^4+s^2}}} \exp[-d^2\frac{(v_jt-x_j)^2}{2(d^4+s^2)}] \quad.
\end{equation}
Il risultato è {sorprendente}. La gaussiana che descrive la densità di probabilità associata alla particella libera ha il suo massimo per $x_j = v_jt$, esattamente la posizione assunta classicamente dalla particella. Ecco trovata la tanto agognata, e semiclassica, \emph{localizzazione}. 

Definiamo in questo senso, e ora sarà chiaro il perché, $\Phi(\bm{r}, t)$ come \textbf{stato localizzato} o \textbf{pacchetto d'onda gaussiano}.

\subsubsection{Una progressiva delocalizzazione}
Abbiamo trovato un risultato specificatamente \emph{semiclassico}. \emph{Semi}-classico, per l'appunto. L'assurdo, almeno classicamente, lo ritroviamo analizzando e il massimo del pacchetto gaussiano e la sua ampiezza. Quest'ultima, nota come \emph{varianza} $\sigma$, corrisponde all'inverso del coefficiente moltiplicativo della variabile, eventualmente traslata, della funzione:
\[
    \sigma = 2\left( d^2 + \frac{s^2}{d^2} \right) = 2\left( d^2 + \frac{\hbar^2k^4}{4m^2d^2}t^2 \right) 
\]
che tende evidentemente a infinito quando $t \to \infty$. 

Il massimo, ottentuo per $x_j = v_jt$, assume come valore:
\[
    \max{\rho_{x_j}} = \frac{|A_j|^2}{{2\sqrt{d^4+s^2}}} = \frac{|A_j|^2}{2} \left( d^4 + \frac{\hbar^2k^4}{4m^2}t^2 \right)^{-\frac{1}{2}} 
\]
che diventa invece infinitesimo per tempi grandi. La densità di probabilità tende ad \emph{appiattirsi} con il passare del tempo, con il picco della gaussiana sempre meno "concentrato" nel suo massimo e quanto mai \emph{delocalizzato}. Ecco, in ultimo, un irreale effetto quantistico che si presenta ai nostri occhi: l'informazione riguardo la probabilità di trovare la particella in un determinato punto dello spazio diventa via via meno \emph{precisa} con il passare del tempo. A tutti gli effetti, il corpuscolo non è deterministiscamente confinato in un punto dello spazio, bensì risulta \textbf{delocalizzato}.

\section{Il Teorema di Ehrenfest}
Per chiudere il discorso riguardo l'evoluzione temporale dell'informazione quantistica, dobbiamo aggiungere ancora un importantissimo tassello che ci permetterà di estendere la nostra analisi all'evoluzione di quantità legate agli operatori e di porre ulteriori ponti (o distruggerne altri) con la meccanica classica. Discuteremo infatti in questa sezione un risultato fondamentale della meccanica quantistica, il \textit{Teorema di Ehrenfest}, fornendo preliminarmente la definizione di \textbf{valore di aspettazione} di un operatore fisico.

\subsection{Valore di aspettazione di un operatore}
Definiamo il {valore di aspettazione} di un operatore, o \emph{media}, relativo ad un generico stato $\psi(\bm{r},t)$ (definito sul volume $\Omega \subseteq \mathbb{R}^3$) la quantità reale:
\begin{equation}
        \langle \h{A} \rangle \equiv (\psi, \h{A} \psi) = \int_\Omega \dd^3{\bm{r}}\,\psi(\bm{r},t)^* \h{A}\psi(\bm{r},t) \quad. 
        \label{eq:exp_value}
\end{equation}
Il significato del prodotto scalare appena formulato è intrinsecamente probabilistico e al concetto di \emph{misura}\footnote{Si noti anche che, benchè sotto altra veste, sia nascosta nella definizione appena fornita la caratteristica di ogni operatore fisico di essere hermitiano. Se così non fosse, $\langle \h{A} \rangle \neq \langle \h{A} \rangle^*$ implicando $\langle \h{A} \rangle \notin \mathbb{R}$, impossibile essendo una misura necessariamente reale.} ed è profondamente connesso alla nozione di set completo di autofunzioni di un operatore hermitiano, ma tratteremo queste peculiarità più avanti. Per ora ci limitiamo a calcolare completamente un interessante valore di aspettazione di un operatore ormai ampiamente noto: l'operatore posizione $\hat{\bm{r}}$.

\subsubsection{Valore di aspettazione della posizione}
Ricordando che $\hat{\bm{r}}$ è un operatore moltiplicativo, usiamo la \eqref{eq:exp_value} e otteniamo:
\[
    \langle \hat{\bm{r}} \rangle = (\psi, \hat{\bm{r}} \psi) = \int_\Omega \dd^3{\bm{r}}\,\psi(\bm{r},t)^* \hat{\bm{r}}\psi(\bm{r},t) = \int_\Omega \dd^3{\bm{r}}\, |\psi(\bm{r},t)|^2 {\bm{r}}\quad.
\]
Il risultato appena trovato ha una limpida interpretazione. Possiamo infatti riscrivere l'integrale come: 
\[
    \langle \hat{\bm{r}} \rangle = \int_\Omega \dd^3{\bm{r}}\, \rho(\bm{r},t) {\bm{r}}
\]
e vedere che stiamo ottenendo il \emph{valore medio} dell'operatore posizione come "somma ponderata" (tramite densità di probabilità) di tutte le possibili posizioni assunte dal sistema. Stiamo quindi ricostruendo la controparte quantistica del \textbf{centro di massa} di un sistema, infatti:
\[
    \langle \hat{\bm{r}} \rangle_{\text{classico}} = \bm{r}_{\text{CDM}} = \frac{\sum_i m_i \bm{r}_i}{\sum_i m_i} = \sum_i \underbrace{\frac{m_i}{M}}_{\rho_{\text{classica}}} \bm{r}_i 
\]

\subsection{Considerazioni preliminari}
Nella formulazione classica fornitaci da Hamilton, la derivata rispetto al tempo di una generica grandezza $A$ che vive nello spazio delle fasi è data da:
\[
    \dv{A}{t} = \{A,H\} + \pdv{A}{t} \quad.
\]
L'idea è quella di cercare di capire se esiste un modo speculare in meccanica quantistica per calcolare l'evoluzione temporale del valore di aspettazione di un operatore. E Ehrenfest ci viene in soccorso.


\subsection{L'hamiltoniano come operatore hermitiano}
Per poter giungere formalmente al risultato ricercato, dobbiamo preliminarmente dimostrare che l'operatore $\h{H}$ è hermitiano. Per farlo, invochiamo la definizione e, armati di pazienza, svolgiamo i calcoli:
\[
\begin{aligned}
    (\psi, \h{H}\phi)^* &= \int_\Omega \dd^3{\bm{r}} \, \left[\psi(\bm{r},t)^* \h{H}\phi(\bm{r},t)\right]^* \\
    &= \int_\Omega \dd^3{\bm{r}} \, \psi(\bm{r},t) \left[ -\frac{\hbar^2}{2m}\nabla^2\phi(\bm{r},t) + U(\bm{r})\phi(\bm{r},t)\right]^* \\
    &\overset{(1)}{=} \int_\Omega \dd^3{\bm{r}} \, \psi(\bm{r},t) U(\bm{r}) \phi(\bm{r},t)^*  -\frac{\hbar^2}{2m}\int_\Omega \dd^3{\bm{r}} \, \psi(\bm{r},t) \nabla^2 \phi(\bm{r},t)^*\\
    &= (\phi, U\psi) - \frac{\hbar^2}{2m} \underbrace{\int_\Omega \dd^3{\bm{r}} \, \psi(\bm{r},t) \nabla^2 \phi(\bm{r},t)^*}_{(a)} \quad,
\end{aligned}
\]
dove in $(1)$ abbiamo sfruttato il fatto che l'energia cinetica è un operatore reale e abbiamo assunto che il potenziale sia puramente reale ($U = U^*$). Per continuare, integriamo $(a)$ per parti:
\[
\begin{aligned}
    (a) &= \int_\Omega \dd^3{\bm{r}} \, \left\{ \nabla \cdot[\psi(\bm{r},t)\nabla\phi(\bm{r},t)^*] - \nabla \psi(\bm{r},t) \cdot \nabla\phi(\bm{r},t)^* \right\} \\
    &= \underbrace{\oint_{\partial\Omega} \psi(\bm{r},t)\nabla\phi(\bm{r},t)^* \cdot \dd{\bm{S}}}_{0} - \int_\Omega \dd^3{\bm{r}} \, \nabla \psi(\bm{r},t) \cdot \nabla\phi(\bm{r},t)^* \quad.
\end{aligned}
\] 
Il primo integrale si annulla poiché le funzioni d'onda fisiche devono essere a quadrato sommabile, decadendo a zero all'infinito (oppure annullandosi sulla frontiera $\partial\Omega$). Integriamo ancora il secondo termine per parti:
\[
    \begin{aligned}
       (a) &= - \int_\Omega \dd^3{\bm{r}} \, \{\nabla\cdot [\phi(\bm{r},t)^*\nabla\psi(\bm{r},t)] - \phi(\bm{r},t)^*\nabla^2\psi(\bm{r},t) \} \\
       &= -\underbrace{\oint_{\partial\Omega} \phi(\bm{r},t)^*\nabla\psi(\bm{r},t) \cdot \dd{\bm{S}}}_0 + \int_\Omega \dd^3{\bm{r}} \, \phi(\bm{r},t)^*\nabla^2\psi(\bm{r},t) \quad.
    \end{aligned}
\]
Ancora una volta, per le medesime ragioni discusse in precedenza, il termine di flusso di frontiera è nullo. Sostituendo il risultato ottenuto, possiamo finalmente concludere che:
\begin{equation}
    \begin{aligned}
        (\psi, \h{H}\phi)^* &= (\phi, U\psi) -\frac{\hbar^2}{2m} \int_\Omega \dd^3{\bm{r}} \, \phi(\bm{r},t)^*\nabla^2\psi(\bm{r},t) \\
        &= (\phi, U\psi) + (\phi, \h{K}\psi) \\
        &= (\phi, \h{H}\psi) \quad. 
    \end{aligned}
\end{equation}

\subsection{La formulazione di Ehrenfest}
Siamo finalmente pronti per comprendere il teorema che stiamo per dimostrare. Partiamo dal chiederci: come varia rispetto al tempo il valore di aspettazione del generico operatore $\h{A}$? Ossia, quanto vale la \emph{derivata totale} rispetto al tempo di $\langle \h{A} \rangle$? Calcoliamo:
\[
\begin{aligned}
    \dv{\langle \h{A} \rangle}{t} = \dv{(\psi, \h{A}\psi)}{t} &= \dv{}{t}\int_\Omega \dd^3{\bm{r}} \, \psi(\bm{r},t)^* \h{A} \psi(\bm{r},t) \\
    &=  \int_\Omega \dd^3{\bm{r}} \, \pdv{}{t}\left[ \psi(\bm{r},t)^* \h{A} \psi(\bm{r},t) \right] \\
    &= \int_\Omega \dd^3{\bm{r}} \, \left[ \pdv{\psi(\bm{r},t)^*}{t}\h{A}\psi(\bm{r},t) + \psi(\bm{r},t)^*\pdv{\h{A}}{t}\psi(\bm{r},t) + \psi(\bm{r},t)^*\h{A}\pdv{\psi(\bm{r},t)}{t} \right] \\
    &= \left(\partial_t\psi, \h{A}\psi\right) + \left(\psi, \partial_t\h{A} \, \psi \right) + \left(\psi, \h{A} \, \partial_t\psi\right) \quad.
    \end{aligned}
\]
Dalla S.E. abbiamo che $i\hbar \partial_t \psi = \h{H}\psi$. Implementiamo questa relazione nel nostro calcolo e otteniamo:
\[
    \begin{aligned}
        \dv{\langle \h{A} \rangle}{t} &= \left( \frac{1}{i\hbar}\h{H}\psi, \h{A}\psi \right) + \left\langle \pdv{\h{A}}{t} \right\rangle + \left(\psi, \h{A}\left(\frac{1}{i\hbar}\h{H}\psi\right) \right) \\
        &= -\frac{1}{i\hbar} \left( \h{H}\psi, \h{A}\psi \right) + \frac{1}{i\hbar} \left(\psi, \h{A}\h{H}\psi\right) + \left\langle \pdv{\h{A}}{t} \right\rangle \\
        &\overset{(1)}{=} -\frac{1}{i\hbar} \left( \psi, \h{H}\h{A}\psi \right) + \frac{1}{i\hbar} \left(\psi, \h{A}\h{H}\psi\right) + \left\langle \pdv{\h{A}}{t} \right\rangle \\
        & = \frac{1}{i\hbar}\left( \psi, \left( \h{A}\h{H} - \h{H}\h{A} \right) \psi \right) + \left\langle \pdv{\h{A}}{t} \right\rangle = \frac{1}{i\hbar}\left( \psi, \left[ \h{A},\h{H} \right] \psi \right) + \left\langle \pdv{\h{A}}{t} \right\rangle \quad,
    \end{aligned}
\]
dove in $(1)$ abbiamo sfruttato l'hermitianità appena dimostrata dell'operatore hamiltoniano. Riconoscendo la definizione di valore di aspettazione, otteniamo quindi la formulazione generale del Teorema di Ehrenfest:
\begin{equation}
    \dv{\langle \h{A} \rangle}{t} = \frac{1}{i\hbar}\left\langle \left[ \h{A},\h{H} \right] \right\rangle + \left\langle \pdv{\h{A}}{t}\right\rangle\quad.
    \label{eq:teo_ehrenfest}
\end{equation}

\subsubsection{Interpretazione del risultato ottenuto}
In pieno accordo con il principio di quantizzazione canonica, la formulazione hamiltoniana dell'evoluzione di una generica grandezza nello spazio delle fasi rispecchia in maniera pressoché perfetta la teoria quantistica, a meno del fantomatico fattore $1/i\hbar$ predetto da Dirac. 

Si noti che, così come classicamente l'indipendenza della generica variabile $A$ esplicita dal tempo e dalla funzione di Hamilton garantisce che $A$ sia una costante del moto del sistema, lo stesso avviene in termini quantistici. L'evoluzione del valore di aspettazione di un operatore quantistico è infatti, a tutti gli effetti, \textit{governata dalle stesse leggi che interessano la sua controparte classica}.

\subsection{Esempi notevoli di applicazioni del Teorema}
Proviamo a questo punto a calcolare l'evoluzione temporale di $\langle  \hat{\bm{r}}\rangle $ e $\langle  \hat{\bm{p}}\rangle $, ricordando che classicamente si ha $\dot{\bm{r}} = \bm{v}$ e $\dot{\bm{p}} = \bm{F}$.

\subsubsection{Dalla posizione ...}
Sfruttiamo la \eqref{eq:teo_ehrenfest} e otteniamo:
\[
    \dv{\langle \hat{\bm{r}} \rangle}{t} = \frac{1}{i\hbar}\left\langle \left[ \hat{\bm{r}},\h{H} \right] \right\rangle + \left\langle \pdv{\hat{\bm{r}}}{t}\right\rangle \quad.
\]
Assumendo il sistema conservativo, il commutatore vale:
\[
\begin{aligned}
    \left[ \hat{\bm{r}},\h{H} \right] = \left[ \hat{\bm{r}},\h{U} + \h{K} \right] &= \underbrace{\left[ \hat{\bm{r}},\h{U} \right]}_{0} + \left[ \hat{\bm{r}},\h{K} \right] \\
    &= \frac{1}{2m}\left[ \hat{\bm{r}},\hat{\bm{p}}^2 \right] = \frac{1}{2m}\sum_i \left(\hat{p}_i\left[ \hat{x}_i,\hat{p}_i \right] +\left[ \hat{x}_i,\hat{p}_i \right] \hat{p}_i\right)\bm{u}_i \\
    &= \frac{1}{2m}\sum_i 2i\hbar \hat{p}_i \bm{u}_i = i\hbar\frac{\hat{\bm{p}}}{m} \quad.
\end{aligned}
\]
Sostituendo, abbiamo:
\[
    \dv{\langle \hat{\bm{r}} \rangle}{t} = \frac{1}{i\hbar} \left\langle i\hbar\frac{\hat{\bm{p}}}{m}\right\rangle + \left\langle \pdv{\hat{\bm{r}}}{t}\right\rangle = \left\langle \frac{\hat{\bm{p}}}{m}\right\rangle + \left\langle \pdv{\hat{\bm{r}}}{t}\right\rangle \quad.
\]
Nella maggior parte dei casi, l'operatore posizione non dipende esplicitamente dal tempo e otteniamo la traduzione della arcinota relazione classica:
\begin{equation}
    \dv{\langle \hat{\bm{r}} \rangle}{t} = \frac{1}{m} \langle \hat{\bm{p}}\rangle \quad.
\end{equation}


\subsubsection{... Alla quantità di moto}
Ripetiamo gli stessi calcoli fatti in precedenza:
\[
    \dv{\langle \hat{\bm{p}} \rangle}{t} = \frac{1}{i\hbar}\left\langle  \left[ \hat{\bm{p}},\h{H} \right] \right\rangle + \left\langle \pdv{\hat{\bm{p}}}{t}\right\rangle \quad.
\]
Si noti che $[\hat{\bm{p}}, \h{K}] = 0$, da cui:
\[
\begin{aligned}
    \left[ \hat{\bm{p}},\h{H} \right] = \left[ \hat{\bm{p}},\h{U} + \h{K} \right] &= \left[ \hat{\bm{p}},\h{U} \right] + \underbrace{\left[ \hat{\bm{p}},\h{K} \right]}_0 \\
    &= \sum_i -i\hbar\pdv{U}{x_i} \, \bm{u}_i \overset{(1)}{=} i\hbar \bm{F} \quad,
\end{aligned}
\]
dove in $(1)$ abbiamo definito l'$i$-esima componente della forza come $F_i = -\partial_{x_i}U$, mimando la formulazione classica. Ricostruendo, otteniamo:
\[
    \dv{\langle \hat{\bm{p}} \rangle}{t} = \frac{1}{i\hbar}\left\langle i\hbar \bm{F} \right\rangle + \left\langle \pdv{\hat{\bm{p}}}{t}\right\rangle = \left\langle  \bm{F} \right\rangle + \left\langle \pdv{\hat{\bm{p}}}{t}\right\rangle\quad.
\]
Ancora una volta, assumiamo $\hat{\bm{p}}$ non dipendente esplicitamente dal tempo e ritroviamo la controparte quantomeccanica della seconda legge di Newton:
\begin{equation}
    \dv{\langle \hat{\bm{p}} \rangle}{t} = \left\langle  \bm{F} \right\rangle \quad.
\end{equation}

\begin{approfondimento}[title={Eccezioni alla regola}]
    \noindent
    In realtà, la "perfetta" corrispondenza classica-quantistica tende a rompersi in alcune particolari situazioni come, ad esempio, casi in cui il potenziale in cui è immerso il sistema non è una forma al più quadratica delle posizioni, bensì presenta una dipendenza cubica, quartica, ecc. dalle coordinate. 
    
    Scelto il semplice potenziale unidimensionale $U(x) = \alpha x^4/4$, abbiamo:
    \[
        \left\langle  \bm{F} \right\rangle = -\left\langle  \pdv{U}{x} \right\rangle = -\left\langle  \alpha x^3 \right\rangle = -\alpha \left\langle   x^3 \right\rangle
    \] 
    che, in generale, è diverso da $-\alpha \left\langle   x \right\rangle^3$. Per \emph{energie sufficientemente grandi}\footnote{Sarà chiaro più avanti come questa assunzione si traduca formalmente nell'ipotesi di avere un sistema caratterizzato da \emph{numeri quantici} sufficientemente grandi.}, comunque, si può dimostrare che $\langle AB \rangle = \langle A \rangle \langle B \rangle$, concludendo che l'uguaglianza è verificata anche nel caso di $\langle x^3 \rangle = \langle x \rangle^3$. 

    In questo senso, il Teorema di Ehrenfest si inserisce in un quadro più ampio completato dal \textbf{Principio di Corrispondenza} di Bohr, che stabilisce che nel limite semiclassico ("grandi energie") le leggi della meccanica quantistica collassano su quelle della fisica classica.
    
\end{approfondimento}

%!TEX root = ../dispense_QP.tex

\chapter{Gli spazi di Hilbert}
In questo capitolo tenteremo di dare una caratterizzazione più matematicamente esaustiva dello spazio delle funzioni a quadrato sommabile $\mathcal{L}^2$, identificandolo a tutti gli effetti come uno \textbf{spazio di Hilbert}. Discuteremo quindi quali sono le peculiarità di uno spazio del genere e capiremo \emph{perché} ci serve che le funzioni d'onda vivano proprio in $\mathcal{L}^2$. 

Introduciamo preliminarmente la notazione di Dirac, con il fine di rendere più agevole la trattazione.

\section{Alcune nozioni basilari di algebra}
Iniziamo rispolverando rapidamente alcune nozioni fondamentali di algebra che saranno necessarie ai fini della comprensione dei concetti discussi. 

\begin{definizione}[Spazio vettoriale]
    Sia $\mathbb{K}$ un campo. Si definisce \emph{spazio vettoriale} $V$ su $\mathbb{K}$ un insieme di elementi, detti vettori:
    \[
        \ket{a}, \quad \ket{b}, \quad \ket{c}, \quad \dots
    \]
    su cui sono definite le operazioni di \textbf{somma} e \textbf{prodotto per scalare}. È inoltre necessario che siano rispettate le seguenti caratteristiche: $(1)$ commutatività e chiusura rispetto alla somma; $(2)$ associatività rispetto alla somma; $(3)$ associatività rispetto al prodotto; $(4)$ distributività rispetto al prodotto; $(5)$ chiusura rispetto al prodotto per scalare; $(6)$ esistenza dell'elemento nullo; $(7)$ esistenza dell'inverso nel senso della somma.
\end{definizione}
Nel nostro caso e nelle trattazioni di base, il campo $\mathbb{K}$ coincide con il campo complesso o quello reale $\mathbb{K} \in \{ \mathbb{R}, \mathbb{C} \}$. Abbiamo bisogno ora della nozione di base di uno spazio vettoriale.

\begin{definizione}[Base di uno spazio vettoriale]
    Sia $V$ uno spazio vettoriale su campo $\mathbb{K}$ e $\mathcal{B} \coloneq \{ \ket{\bm{w}_i} \in V : i \in \mathbb{N}^+ \} $ un insieme di vettori\footnote{Si denota con $\mathbb{N}^+$ l'insieme dei numeri naturali strettamente maggiori di zero.}. Si dice che $\mathcal{B}$ è una base di $V$ se: $(1)$ $\mathcal{B}$ è \textbf{libero} (ovvero è formato da elementi linearmente indipendenti); $(2)$ qualsiasi vettore dello spazio in questione può essere scritto come combinazione lineare degli elementi di $\mathcal{B}$. 

    Equivalentemente, si dice che $\mathcal{B}$ è un insieme \textbf{completo}.
\end{definizione}
Dal momento che abbiamo già introdotto la nozione di prodotto scalare (tanto nello spazio euclideo $\mathbb{E}^n$ quanto in $\mathcal{L}^2$ come forma sesquilineare), ci limitiamo a segnalare che la notazione di Dirac permette di scrivere in maniera compatta:
\[
    \bm{a} \cdot \bm{b} = \bra{\bm{a}} \ket{\bm{b}} = (\bm{a}, \bm{b}) \quad.
\]
In questo senso, le componenti dei due vettori si scrivono come:
\[
    \ket{\bm{b}} = \sum_i b_i \ket{\bm{e}_i}, \quad \bra{\bm{a}} = \sum_j a_j^* \bra{\bm{e}_j},
\]
dove abbiamo indicato con $\bm{e}_i$ l'$i$-esimo versore di base. Mettendo tutto insieme, otteniamo nella nuova scrittura del prodotto scalare:
\[
    \bra{\bm{a}} \ket{\bm{b}} = \sum_i \sum_j a_j^*b_i \bra{\bm{e}_j}\ket{\bm{e}_i}. 
\]
Una caratteristica fondamentale del prodotto scalare è il fatto che induce una particolare nozione di metrica, detta \textbf{norma} (il cui corrispondente negli spazi euclidei è il \emph{modulo} del vettore in questione). 

\begin{definizione}[Norma di un vettore]
    Sia $\ket{\bm{v}} = \sum_j v_j \ket{\bm{e}_j}$ un vettore. Si definisce \emph{norma} di $\ket{\bm{v}}$ la radice quadrata del prodotto scalare del vettore per se stesso:
    \[
        \norm{\bm{v}} \coloneq \sqrt{\bra{\bm{v}} \ket{\bm{v}}} = \sqrt{\sum_j \sum_i v_i^* v_j \bra{\bm{e}_i}\ket{\bm{e}_j}} \quad.
    \]
\end{definizione}
Un qualsiasi spazio vettoriale in cui è definita una {norma} e che risulta \emph{completo}\footnote{Si ricorda che uno spazio è completo se e solo se qualsiasi successione di elementi che stanno al suo interno converge ad un elemento dello spazio di partenza.} viene detto \textbf{spazio di Banach}. Ancora, se uno spazio è dotato di prodotto scalare si dice invece \textbf{spazio pre-hilbertiano}. Uno spazio pre-hilbertiano che è anche di Banach viene detto di \textbf{Hilbert}.

Per definire, almeno operativamente, questa fantomatica entità, dobbiamo ancora iquadrare una base ortonormale.

\begin{definizione}[Base ortonormale]
    Sia $V$ uno spazio vettoriale su campo $\mathbb{K}$ e $\mathcal{B} \coloneq \{ \ket{\bm{w}_i} \in V : i \in \mathbb{N}^+ \}$ una base in $V$. Si dice che $\mathcal{B}$ è una \emph{base ortonormale} se vale:
    \[
        \bra{\bm{w}_i} \ket{\bm{w}_j} = \delta_{ij} \quad \forall i,j \in \mathbb{N} \quad.
    \]
\end{definizione}
In questi casi, la norma di un vettore si semplifica in:
\[
    \norm{\bm{v}} = \left(\sum_j |v_j|^2 \right)^{\frac{1}{2}} \quad.
\]
Siamo pronti per dare la definizione di spazio di Hilbert. 
\begin{definizione}[Spazio di Hilbert]
    Sia $H$ uno spazio vettoriale infinitamente generato su campo complesso $\mathbb{C}$. $H$ è uno \textit{spazio di Hilbert} se: $(1)$ $H$ è uno spazio normato; $(2)$ $H$ è uno spazio completo; $(3)$ in $H$ esiste una nozione di prodotto scalare. Se queste condizioni sono soddisfatte, vale:
    \[
        \ket{\bm{v}} = \sum_{i = 0}^\infty v_i \ket{\bm{e}_i}, \quad \bra{\bm{v}}\ket{\bm{v}} = \sum_i |v_i|^2 < \infty \quad.
    \]
    In questo senso, uno spazio di Banach in cui la nozione di metrica metrica è garantita da una norma indotta da un prodotto scalare è detto di Hilbert.
\end{definizione}
Si noti come la seconda proprietà sia di vitale importanza per assicurare che il prodotto scalare fra due generici vettori non diverga. Infatti possiamo scrivere:
\[
    \left|\bra{\bm{a}}\ket{\bm{b}}\right|^2 = \left| \sum_i a_i^* b_i \right|^2 \overset{(1)}{\leq} \left( \sum_i |a_i|^2 \right) \left( \sum_i |b_i|^2 \right) = \norm{\bm{a}}^2 \norm{\bm{b}}^2 < \infty \quad,
\]
dove in $(1)$ abbiamo utilizzato la disuguaglianza di Cauchy-Schwarz e, alla fine, il fatto che qualsiasi vettore abbia norma $\in \mathbb{R}$.

\section{Un interessante spazio di Hilbert}
Detto tutto questo, a cosa ci è servito? La grande verità a cui ci stiamo per esporre è che la meccanica quantistica esiste solo grazie alla nozione matematica dello spazio di Hilbert. L'occhio attento avrà già notato, infatti, che lo spazio delle funzioni a quadrato sommabile $\mathcal{L}^2$ è \textbf{isomorfo} all'$H$ appena definito e, dunque, ne eredita tutte le proprietà. In questa sezione dimostreremo l'intuizione in questione. 

\subsection{Caratteristiche di \texorpdfstring{$\bm{\mathcal{L}^2}$}{L2}}
È immediato verificare che gli assiomi che caratterizzano uno spazio vettoriale valgono per $\mathcal{L}^2$: lo spazio delle funzioni a quadrato sommabile è dunque uno spazio vettoriale. In particolare, vale la nozione di combinazione lineare che garantisce che, in un certo volume $\Omega$:
\[
    \Psi = \sum_i c_i \psi_i \in \mathcal{L}^2(\Omega), \quad \psi \in \mathcal{L}^2, \quad c \in \mathbb{C} 
\]
suggerendo, in qualche modo, il già sfruttato \textbf{principio di sovrapposizione}. 
Abbiamo inoltre definito un prodotto scalare in $\mathcal{L}^2$ per due funzioni a quadrato sommabile $\psi, \phi$:
\[
    \bra{\psi}\ket{\phi} = (\psi,\phi) = \int_{\Omega \subseteq \mathbb{R}^3} \dd^3{\bm{r}} \, \psi^*\phi \quad.
\]
In particolare, gli stati quantistici soddisfano la \emph{condizione di normalizzazione} $\bra{\psi}\ket{\psi} = (\psi,\psi) = 1$, la cui validità è ciò su cui si fonda l'interpretazione probabilistica della teoria ed è dunque fondamentale. 

Della prossima affermazione bisognerà "fidarsi", anche se verranno presentati esempi che renderanno l'atto di fede meno astratto. Tra le funzioni d'onda $\psi \in \mathcal{L}^2$, possiamo \emph{sempre} estrarre un set che forma una base ortonormale per lo spazio delle funzioni a quadrato sommabile. Sosteniamo quindi che, esiste sempre almeno un insieme $\mathcal{B} \coloneq \{ \psi_i \in \mathcal{L}^2(\Omega) : i \in \mathbb{N}^+, \bra{\psi_i}\ket{\psi_j} = \delta_{ij} \}$. Si noti che la normalità è garantita dalla condizione di normalizzazione, infatti $\bra{\psi_i}\ket{\psi_i} = 1$. Diciamo che $\mathcal{B}$ è \textbf{completo} se possiamo esprimere un qualsiasi stato quantistico $\Psi$ come combinazione lineare degli elementi della base:
\begin{equation}
    \Psi = \sum_i c_i{\psi_i}, \quad c_i \overset{(1)}{=} \bra{\psi_i}\ket{\Psi} \in \mathbb{C} \quad.
    \label{eq:psi_vec}
\end{equation}
Questa rappresentazione corrisponde a una generalizzazione della \textbf{serie di Fourier} e ci riferiremo dunque ad essa in tal senso. Sfruttando la nozione di base ortonormale, poi, possiamo giustificare $(1)$, infatti:
\[
    \bra{\psi_i}\ket{\Psi} = \int_{\mathbb{R}^3} \dd^3{\bm{r}} \, \psi_i^* \Psi = \int_{\mathbb{R}^3} \dd^3{\bm{r}} \,\psi_i^* \sum_j c_j \psi_j = \sum_j c_j \int_{\mathbb{R}^3} \dd^3{\bm{r}} \, \psi_i^* \psi_j = \sum_j c_j \underbrace{\bra{\psi_i}\ket{\psi_j}}_{\delta_{ij}} = c_i \quad,
\]
con i vari coefficienti di combinazione che prendono il nome di \textbf{ampiezze}. L'interpetazione fisica associata al concetto di "ampiezza di una funzione" di base risiede nell'ideologia di Born: $|c_i|^2$ corrisponde infatti ad una probabilità, in particolare quella di ottenere lo stato $\psi_i$ associato a $c_i$ eseguendo una misura di qualche tipo su $\Psi$\footnote{Se quest'ultima considerazione non risulta chiara non c'è da preoccuparsi. La motivazione dell'assunzione è infatti da ricercarsi in argomenti più avanzati, trattati nell'ambito della teoria della misura.}

La \eqref{eq:psi_vec} ha un'interpretazione peculiare: ad ogni stato $\Psi$ è associata un'\textit{unica} $n$-upla di valori complessi $c_i$ che possono essere interpretati come le componenti dello stato stesso. Questo fatto suggerisce che esiste una corrispondenza 1:1 fra funzione d'onda in $\mathcal{L}^2$ e vettore dello spazio di Hilbert $\ket{c} = [c_1, c_2, \dots, c_1, \dots]^T$:
\[
    \psi = \sum_i c_i \psi_i \iff \ket{c} = \sum_j c_j\ket{j} \quad,
\]
con la nuova nozione di "versori di base" $\ket{j}$ per indicare la $j$-esima componente\footnote{Si noti come non si usa più il grassetto entro le "bra" e le "ket della notazione di Dirac per indicare i vettori. Questo, seppur inusuale sino ad ora, non è un errore e sarà prassi per snellire la notazione.}. La formula che definisce la \textbf{norma} quadra di una funzione d'ondarende ancora più evidente questa corrispondenza:
\[
    \begin{aligned}
        \bra{\Psi}\ket{\Psi} &= \int_{\mathbb{R}^3} \dd^3{\bm{r}} \, \Psi^* \Psi = \int_{\mathbb{R}^3} \dd^3{\bm{r}} \, \left(\sum_i c_i \psi_i \right)^*\sum_j c_j \psi_j \\
        &= \sum_i \sum_j c_i^* c_j \int_{\mathbb{R}^3} \dd^3{\bm{r}} \, \psi_i^* \psi_j = \sum_i \sum_j c_i^* c_j \bra{\psi_i}\ket{\psi_j} \overset{(1)}{=} \sum_i \sum_j c_i^* c_j \delta_{ij} = \sum_i |c_i|^2 \quad,
    \end{aligned}
\] 
dove in $(1)$ abbiamo sfruttato l'ortonormalità della base tramite cui abbiamo \emph{rappresentato} $\Psi$. Quanto ottenuto corrisponde alla cosiddetta \textbf{identità di Parseval} che rispecchia in tutto e per tutto la formulazione della norma quadra di un vettore nello spazio di Hilbert. 

\begin{approfondimento}[title={Limiti di validità dell'identità di Parseval}]
    \noindent
    Si può dimostrare che l'identità di Parseval è un'uguaglianza se e solo se la base tramite cui si rappresenta lo stato $\Psi$ di cui si calcola la norma è \textbf{completa}. Altrimenti vale la cosiddetta \emph{disuguaglianza di Bessel}:
    \[
        \norm{\Psi}^2 \geq \sum_i |c_i|^2 \quad.
    \]
\end{approfondimento}
Abbiamo sino ad ora mostrato che $\mathcal{L}^2$ è uno spazio dotato di norma (normato) in cui è possibile definire una base ortonormale completa di funzioni d'onda. Per attribuire però allo spazio delle funzioni a quadrato sommabile la caratteristica di essere uno spazio di Hilbert dobbiamo ancora dimostrare il fatto che la norma di una generica funzione d'onda sia strettamente convergente, qualsiasi essa sia. Enunciamo a tal fine una formulazione operativa del \textbf{Teorema di Fisher-Riesz}, del quale non sarà proposta la dimostrazione in quanto eccessivamente complicata.

\begin{teorema}[Fisher-Riesz - formulazione operativa]
    Sia $\Omega \subseteq \mathbb{R}^3$ e $\mathcal{B} \coloneq \{ \psi_i \in \mathcal{L}^2(\Omega) : i \in \mathbb{N}^+ \}$ una base per lo spazio delle funzioni a quadrato sommabile $\mathcal{L}^2(\Omega)$. Allora disuguaglianza:
    \[
        \sum_{i = 1}^\infty |c_i|^2 < \infty, \quad c_i \in \mathbb{C} \quad,
    \]
    con i coefficienti $c_i$ scelti arbitrariamente, è \textbf{condizione necessaria e sufficiente} perché esista una funzione d'onda $\phi \in \mathcal{L}^2(\Omega)$ la cui espansione di Fourier risulti:  
    \[
        \phi = \sum_{i= 1}^\infty c_i \psi_i \quad.
    \]
\end{teorema}
Cosa ci sta comunicando il teorema? Essenzialmente, abbiamo che ad ogni successione di numeri complessi tali che la serie dei loro moduli quadrati converge è associata una e una sola funzione d'onda in $\mathcal{L}^2$. Abbiamo quindi costruito una corrispondenza biunivoca tra lo spazio delle funzioni a quadrato sommabile e quello delle successioni a quadrato sommabile: $\mathcal{L}^2 \iff l^2$. In particolare, dal momento che qualsiasi successione a quadrato sommabile converge a un elemento di $\mathbb{C}$, evidentemente a quadrato sommabile, lo spazio vettoriale $l^2$ è \textbf{completo}. L'isomorfismo con $\mathcal{L}^2$ ci permette allora di attribuire anche a quest'ultimo la nozione di completezza, soddisfando una volta per tutte le nostre richieste.

\begin{approfondimento}[title={Il vero Teorema di Fisher-Riesz}]
    \noindent
    Il Teorema di Fisher-Riesz, enunciato come sui libri di analisi funzionale, recita:
    \begin{teorema}[Fisher-Riesz]
        Tutti gli spazi $\mathcal{L}^p$ sono \textbf{metricamente completi}.
    \end{teorema}
    Questo risultato è potentissimo se colto nella sua immensa profondità. Ricordiamo innanzitutto che la nozione di integrale usata negli spazi $\mathcal{L}^p$ è quella data da Lebesgue (da cui la "L"), così come quella di misura (in luogo dell'integrale di Riemann e della misura di Peano-Jordan). L'integrale di Lebesgue ci permette di trattare un set di funzioni che si comportano meno bene di quelle integrabili secondo Riemann (si pensi alla funzione di Dirichlet oppure a tutte quelle funzioni che differiscono da una continua quasi-ovunque), risultando perfetto ai fini della meccanica quantistica. Nel nostro caso, la completezza garantita dal teorema è proprio quanto serviva per dimostrare che $\mathcal{L}^2$ non è solo uno spazio pre-hilbertiano ma anche di Banach. 
\end{approfondimento}

\subsubsection{Riassunto delle proprietà di $\bm{\mathcal{L}^2}$}
Abbiamo appena ottenuto le seguenti, fondamentali, caratterizzazioni per lo spazio delle funzioni a quadrato sommabile:
\begin{enumerate}
    \item È possibile definire una nozione di norma.
    \item La norma in questione è indotta da un prodotto scalare.
    \item La nozione di completezza è garantita da Fisher-Riesz.
\end{enumerate}
Combinando tutte queste proprietà abbiamo quanto desideravamo: possiamo finalmente trattare $\mathcal{L}^2$ a tutti gli effetti come uno spazio di Hilbert!

\subsection{La relazione di completezza}
Continuiamo con un po' di matematica. In spazi vettoriali finitamente generati (o, più in generale, per basi a \emph{spettro discreto}), la conoscenza di una base permette di esprimere l'operatore identità $\mathds{1}$ in maniera "speciale". Sfruttando il fatto che $\bra{i}\ket{v} = \sum_j v_j \bra{i}\ket{j} = \sum_j v_j \delta_{ij} = v_i$, abbiamo:
\[
    \ket{v} = \sum_i v_i \ket{i} = \sum_i \bra{i}\ket{v} \ket{i} = \underbrace{\left( \sum_i \ket{i} \bra{i} \right)}_{\mathds{1}} \ket{v} = \ket{v} \quad.
\]
L'equazione che rappresenta l'operatore identità in funzione della base a nostra disposizione:
\begin{equation}
    \mathds{1} =  \sum_i \ket{i} \bra{i}  
\end{equation}
prende il nome di \textbf{relazione di completezza} (per lo spettro discreto) ed è di fondamentale utilità quando si cerca, ad esempio, di eseguire un cambiamento di base. Il nome si deve al fatto che è proprio la completezza della base che permette di scrivere l'operatore identità tramite i suoi proiettori elementari. 

Una grande domanda può sovvenire a questo punto: è possibile formulare un'analoga relazione di completezza quando passiamo a una base caratterizzata da uno \emph{spettro continuo} (come la base delle posizioni per studiare le funzioni in $\mathcal{L}^2$)? La risposta è sì e, per farlo, sfruttiamo la definizione di prodotto scalare. Sia $\mathcal{B} = \{ \phi_i \in \mathcal{L}^2 \}$ una base ortonormale discreta per $\mathcal{L}^2$; il generico stato $\psi$ si può espandere come:
\[
\begin{aligned}
    \psi(\bm{r}) = \sum_i c_i \phi_i(\bm{r}) &= \sum_i \bra{\phi_i}\ket{\psi} \phi_i(\bm{r}) \\
    &= \sum_i \phi_i(\bm{r}) \int_\Omega \dd^3{\bm{s}} \, \phi_i(\bm{s})^* \psi(\bm{s}) = \int_\Omega \dd^3{\bm{s}} \, \underbrace{\left(\sum_i  \phi_i(\bm{s})^* \phi_i(\bm{r})\right)}_{f} \psi(\bm{s}) \quad.
\end{aligned}
\]
Perché l'ultimo integrale restituisca $\psi(\bm{r})$, serve che $f$ svolga il ruolo della Delta di Dirac centrata in $\bm{r}$. Così effettivamente fa. Abbiamo infatti:
\[
    \psi(\bm{r}) = \int_\Omega \dd^3{\bm{s}} \, \underbrace{\left(\sum_i  \phi_i(\bm{s})^* \phi_i(\bm{r})\right)}_{\delta^{(3)}(\bm{r}-\bm{s})} \psi(\bm{s}) \quad.
\]
Definiamo in questo senso una \textbf{relazione di completezza spaziale} per le funzioni a quadrato sommabile:
\begin{equation}
    \delta^{(3)}(\bm{r}-\bm{s}) = \sum_i  \phi_i(\bm{s})^* \phi_i(\bm{r}) \quad.
    \label{eq:tua_equazione_sopra}
\end{equation}
È fondamentale notare la differenza concettuale con il caso discreto: la Delta di Dirac \emph{non} è l'operatore identità, bensì ne costituisce il nucleo integrale (o \emph{kernel}) quando questo viene proiettato nella rappresentazione delle coordinate\footnote{Per gli amanti del rigore matematico, si noti che la Delta di Dirac è una distribuzione a supporto compatto (un singolo punto) e, per definire la sua azione come funzionale lineare, le basta applicarsi a funzioni semplicemente continue in quel punto. Tuttavia, in meccanica quantistica, per garantire che l'applicazione iterata di operatori differenziali (come $\hat{\bm{p}}$) o moltiplicativi (come $\hat{\bm{r}}$) non ci porti mai fuori dallo spazio funzionale di lavoro, si è soliti restringere il dominio e richiedere che $\psi(\bm{r}) \in \mathcal{S}(\mathbb{R}^3)$, dove $\mathcal{S}$ è lo spazio di Schwartz delle funzioni test a decrescita rapida.}.

\begin{approfondimento}[title={La completezza nel continuo e la Delta di Dirac}]
    \noindent
    L'eleganza della notazione di Dirac permette di ricavare la relazione \eqref{eq:tua_equazione_sopra} in modo estremamente naturale, chiarendo definitivamente il ruolo della Delta di Dirac non come "operatore", ma come elemento di matrice. 
    
    Per uno spettro continuo come quello delle posizioni, l'operatore identità si scrive promuovendo la sommatoria a integrale:
    \[
        \mathds{1} = \int_{\mathbb{R}^3} \dd^3{\bm{r}} \, \ket{\bm{r}}\bra{\bm{r}} \quad.
    \]
    Se applichiamo questo operatore a una generica funzione di base, dobbiamo per definizione riottenere $\ket{\bm{s}}$ stesso:
    \[
        \ket{\bm{s}} = \mathds{1}\ket{\bm{s}} = \left( \int_{\mathbb{R}^3} \dd^3{\bm{r}} \, \ket{\bm{r}}\bra{\bm{r}} \right) \ket{\bm{s}} = \int_{\mathbb{R}^3} \dd^3{\bm{r}} \, \ket{\bm{r}} \braket{\bm{r}}{\bm{s}} \quad.
    \]
    Affinché l'integrale restituisca esattamente il vettore $\ket{\bm{s}}$, per le proprietà formali delle distribuzioni, il prodotto scalare al suo interno deve essere necessariamente la Delta di Dirac:
    \[
        \braket{\bm{r}}{\bm{s}} = \delta^{(3)}(\bm{r}-\bm{s}) \quad.
    \]
    A questo punto, l'espansione della Delta nella base discreta $\mathcal{B} = \{ \ket{i} \}$ diventa una banale applicazione dell'identità $\mathds{1} = \sum_i \ket{i}\bra{i}$. Calcoliamo infatti l'elemento di matrice inserendo l'identità discreta "in mezzo" al prodotto scalare:
    \[
        \delta^{(3)}(\bm{r}-\bm{s}) = \braket{\bm{r}}{\bm{s}} = \bra{\bm{r}} \mathds{1} \ket{\bm{s}} = \bra{\bm{r}} \left( \sum_i \ket{i}\bra{i} \right) \ket{\bm{s}} = \sum_i \braket{\bm{r}}{i}\braket{i}{\bm{s}} \quad.
    \]
    Ricordando infine che la funzione d'onda associata allo stato $\ket{i}$ valutata nel punto $\bm{r}$ è proprio definita come la proiezione $\braket{\bm{r}}{i} = \phi_i(\bm{r})$ (e analogamente $\braket{i}{\bm{s}} = \braket{\bm{s}}{i}^* = \phi_i(\bm{s})^*$), si ritrova esattamente la formula:
    \[
        \delta^{(3)}(\bm{r}-\bm{s}) = \sum_i \phi_i(\bm{r}) \phi_i(\bm{s})^* \quad.
    \]
    Una conseguenza interessante della relazione appena trovata è il fatto che, potendo costruire in linea di principio infinite basi ortonormali per $\mathcal{L}^2$, è possibile esprimere la delta di Dirac in \emph{infiniti modi diversi}.
\end{approfondimento}

\section{Alcuni esempi di basi fondamentali}
Il paragrafo che si accinge a cominciare mira a fornire alcuni esempi più concreti della discussione, sin'ora molto astratta, sugli spazi di Hilbert. Analizzeremo a tal fine alcuni semplici esempi di basi per lo spazio $\mathcal{L}^2$ e tenteremo di darne una caratterizzazione il più completa possibile. Seguono la base delle già analizzate \textbf{onde piane} in una \emph{scatola} e nel generico \emph{spazio tridimensionale}.

\subsection{Base delle onde piane in una scatola}
Definiamo innanzitutto una scatola \emph{cubica} di lato $L$ come il sottoinsieme di $\mathbb{R}^3$ così caratterizzato: $\mathbb{R}^3 \supseteq \Omega \coloneq \{ \bm{r} = (x_1,x_2,x_3) \in \mathbb{R}^3 : 0 \leq x_i \leq L \}$. Passiamo ora alla ricerca della base... ma cosa stiamo \emph{effettivamente} facendo? L'idea è quella di trovare un set di funzioni d'onda tramite cui sia possibile esprimere un \textbf{qualunque} stato $\Psi$. Riprendiamo quanto fatto risolvendo la S.E. per il caso della particella libera: la soluzione ci ha portato a trovare gli autostati dell'hamiltoniana che, composti con la parte temporale per l'ansatz dellam separazione delle variabili, ci hanno permesso di definire il generico stato $\Psi$. L'insieme degli autostati di $\h{H}$ è stato per noi, quindi, una base per $\mathcal{L}^2$. Procediamo allora con l'analisi della S.E. in una scatola. 

\subsubsection{Il caso della scatola}
Per una particella descritta dallo stato $\psi$ che vive in $\Omega$, non sottoposta a campi di forze, l'equazione di Schrödinger stazionaria si scrive:
\[
    \h{K} \psi(\bm{r}) = -\frac{\hbar^2}{2m} \nabla^2 \psi(\bm{r}) = \frac{\hbar^2 k^2}{2m} \psi(\bm{r}) \quad.
\]
Imponendo ora che il generico stato sia scomponibile come prodotto di funzioni d'onda indipendenti associate ai tre assi $\psi(\bm{r}) = \prod_j \psi(x_j)$, procediamo a risolvere l'equazione in termini monodimensionali:
\[
    \pdv[2]{\psi(x)}{x} = -k^2 \psi(x) \quad,
\]
dove abbiamo sott'inteso il pedice $x_j$ per la coordinata coinvolta nella derivazione. La soluzione dell'equazione è data dalla combinazione lineare di due onde piane, una progressiva e una regressiva:
\[
    \psi(x) = Ae^{ikx} + Be^{-ikx} \quad.
\]
Si noti che, fino ad ora, non abbiamo menzionato il fatto che la particella è confinata in una scatola o, declinandolo al caso monodimensionale studiato, che deve valere $0 \leq x \leq L$. Poiché il corpuscolo si trova \emph{entro} $\Omega$, serve che $\psi(\Omega^c) = 0$\footnote{Con $\Omega^c$ si è indicato il complementare di $\Omega$.}. Per la continuità quindi, $\psi(0) = \psi(L) = 0$, da cui:
\[
    \begin{cases}
        \psi(0) = A + B = 0 \\
        \psi(L) = Ae^{ikL} + Be^{-ikL} = 0
    \end{cases}
    \implies
    \begin{cases}
        B = -A \\
        A(e^{ikL} - e^{-ikL}) = 2i A \sin(kL) = 0 \quad.
    \end{cases}
\] 
Poiché il fattore $2iA$ è costante, affinché l'ultima uguaglianza sia verificata serve che il seno si annulli:
\[
\begin{aligned}
    kL &= n \pi \\
    k &= \frac{\pi}{L} n, \quad n \in \mathbb{Z} \implies \psi(x) = 2iA \sin(kx) \quad.
\end{aligned}
\]
Il numero d'onda non può assumere valori continui, bensì discreti: $k$ risulta \textbf{quantizzato}! Prima di ricostruire la generica soluzione tridimensionale, dobbiamo ricavare la costante $A$ e, per farlo, imponiamo la normalizzazione:
\[
\begin{aligned}
    1 &= \int_0^L \dd{x} |\psi(x)|^2 = \int_0^{L} \dd{x} 4|A|^2 \sin^2(kx)  \\
    1 &= 2L|A|^2 \\
    A &= \frac{1}{\sqrt{2L}} e^{i\theta}, \quad \theta \in [0, 2\pi] \quad.
\end{aligned}
\]
Imponendo per comodità il fattore di fase pari a $\theta = \frac{3}{2}\pi \implies e^{i\theta} = -i$, otteniamo che 
\[
    \psi(x) = \sqrt{\frac{2}{L}}\sin(kx) \quad,
\]
da cui la funzione tridimensionale risulta:
\begin{equation}
    \psi_{\bm{n}}(\bm{r}) = \left( \frac{2}{L} \right)^{3/2} \prod_{j=1}^3 \sin\left( \frac{n_j \pi x_j}{L} \right) \quad,
\end{equation}
dove si è introdotto il vettore dei numeri quantici $\bm{n} = (n_1, n_2, n_3)$ con $n_j \in \mathbb{N}^+$ per le tre direzioni spaziali. 

\subsubsection{Ortonormalità della base}
Verifichiamo ora rigorosamente che gli autostati trovati formino una base ortonormale per $\mathcal{L}^2(\Omega)$. Per farlo, calcoliamo il prodotto scalare tra due generici stati caratterizzati dai vettori quantici $\bm{n}$ e $\bm{m}$. Sfruttando la separabilità degli integrali lungo i tre assi, abbiamo:
\[
\begin{aligned}
    \braket{\psi_{\bm{n}}}{\psi_{\bm{m}}} &= \int_{\Omega} \dd^3{\bm{r}} \, \psi_{\bm{n}}(\bm{r})^* \psi_{\bm{m}}(\bm{r}) \\
    &= \prod_{j=1}^3 \underbrace{\left[ \frac{2}{L} \int_0^L \dd{x_j} \sin\left( \frac{n_j \pi x_j}{L} \right) \sin\left( \frac{m_j \pi x_j}{L} \right) \right]}_{I_j} \quad.
\end{aligned}
\]
Risolviamo l'integrale monodimensionale $I_j$ utilizzando le formule di prostaferesi per trasformare il prodotto in somma:
\[
    I_j = \frac{1}{L} \int_0^L \dd{x} \left[ \cos\left( \frac{(n_j - m_j)\pi x}{L} \right) - \cos\left( \frac{(n_j + m_j)\pi x}{L} \right) \right] \quad.
\]
Si presentano ora due casi distinti:
\begin{itemize}
    \item Se $n_j \neq m_j$, l'integrazione di funzioni puramente armoniche su un numero intero di periodi restituisce banalmente zero.
    \item Se $n_j = m_j$, il primo termine dell'integranda diventa $\cos(0) = 1$. L'integrale si riduce quindi a:
    \[
        I_j = \frac{1}{L} \int_0^L \dd{x} \left[ 1 - \cos\left( \frac{2n_j\pi x}{L} \right) \right] = \frac{1}{L} \left[ x - \frac{L}{2n_j\pi}\sin\left( \frac{2n_j\pi x}{L} \right) \right]_0^L = \frac{1}{L} [L - 0] = 1 \quad.
    \]
\end{itemize}
Essendo questo risultato valido e indipendente per ogni asse $j$, possiamo ricombinare i termini per ottenere la condizione di ortonormalità globale:
\begin{equation}
    \braket{\psi_{\bm{n}}}{\psi_{\bm{m}}} = \delta_{n_1, m_1} \delta_{n_2, m_2} \delta_{n_3, m_3} \equiv \delta_{\bm{n},\bm{m}} \quad.
\end{equation}
Il fatto che questo set ortonormale costituisca anche un sistema \emph{completo} (e sia dunque una base per $\mathcal{L}^2(\Omega)$) è garantito dal Teorema di Sturm-Liouville per i problemi agli autovalori con condizioni al contorno di Dirichlet.  

\begin{approfondimento}[title={Il Teorema di Sturm-Liouville e la Meccanica Quantistica}]
    \noindent
    Nel discutere la completezza degli autostati dell'Hamiltoniana, abbiamo dato per scontato che le soluzioni dell'equazione differenziale formassero una base per $\mathcal{L}^2$. La solida giustificazione matematica a questa assunzione risiede nel \textbf{Teorema di Sturm-Liouville}.

    Un problema di Sturm-Liouville regolare è definito da un'equazione differenziale lineare del secondo ordine della forma:
    \[
        -\dv{x} \left[ p(x) \dv{y(x)}{x} \right] + q(x)y(x) = \lambda w(x) y(x) \quad,
    \]
    definita su un intervallo $[a,b]$, con opportune condizioni al contorno omogenee (come quelle di Dirichlet $y(a)=y(b)=0$ della nostra scatola).
    
    L'equazione di Schrödinger stazionaria monodimensionale rientra esattamente in questa immensa classe di problemi. Riscrivendola come:
    \[
        -\frac{\hbar^2}{2m} \dv[2]{\psi(x)}{x} + U(x)\psi(x) = E \psi(x) \quad,
    \]
    identifichiamo immediatamente i termini: il coefficiente $p(x) = \frac{\hbar^2}{2m}$ (costante), la funzione $q(x) = U(x)$ (il potenziale), la funzione peso $w(x) = 1$ e il parametro $\lambda = E$ (l'autovalore dell'energia).

    La teoria di Sturm-Liouville assicura tre risultati analitici di portata enorme, che tradotti in fisica si leggono così:
    \begin{enumerate}
        \item \textbf{Spettro reale e limitato inferiormente:} Gli autovalori $\lambda_n$ (ovvero le energie $E_n$) ammessi dal sistema sono puramente reali, numerabili (poiché il dominio è confinato) e ammettono un limite inferiore (esiste sempre uno stato fondamentale), ma non un limite superiore.
        \item \textbf{Ortonormalità degli autostati:} Le autofunzioni $\psi_n(x)$ corrispondenti ad autovalori distinti sono intrinsecamente ortogonali rispetto alla funzione peso $w(x)$. Essendo $w(x)=1$ nel caso di Schrödinger, si ritrova esattamente la definizione del prodotto scalare standard in $\mathcal{L}^2$: 
        \[ \int_a^b \psi_n(x)^* \psi_m(x) \dd{x} = \delta_{nm} \quad. \]
        \item \textbf{Completezza della base:} È il risultato cardine. L'insieme delle autofunzioni costituisce una \emph{base completa} per lo spazio delle funzioni a quadrato sommabile nel dominio considerato. Qualsiasi stato fisico $\Psi(x)$, purché sufficientemente regolare e rispettoso delle condizioni al contorno, può essere espanso in serie in modo unico.
    \end{enumerate}
    Ecco perché, ogni volta che confiniamo una particella quantistica, l'equazione differenziale ci fornisce "gratuitamente" lo spazio vettoriale perfetto in cui far operare la meccanica matriciale.
\end{approfondimento}

\subsubsection{Relazione di completezza}
Avendo stabilito che $\mathcal{B}_{box} = \{ \psi_{\bm{n}} \}$ è una base ortonormale, possiamo scriverne esplicitamente la relazione di completezza spaziale. Richiamando la formula generale discussa in precedenza, ovvero $\sum_i \phi_i(\bm{r}) \phi_i(\bm{r}')^* = \delta^{(3)}(\bm{r}-\bm{r}')$, e sostituendovi i nostri autostati (notando che gli indici corrono ora sulle terne di interi positivi), otteniamo:
\begin{equation}
    \sum_{n_1=1}^\infty \sum_{n_2=1}^\infty \sum_{n_3=1}^\infty \left( \frac{2}{L} \right)^3 \prod_{j=1}^3 \sin\left( \frac{n_j \pi x_j}{L} \right) \sin\left( \frac{n_j \pi x'_j}{L} \right) = \delta^{(3)}(\bm{r}-\bm{r}') \quad.
\end{equation}
Questa espressione analitica asserisce un concetto fisicamente potentissimo: sommando tutte le infinite onde stazionarie quantizzate della scatola, opportunamente pesate sulle posizioni, si ricostruisce esattamente l'operatore identità (sotto forma di distribuzione Delta di Dirac) per ogni punto confinato nel dominio $\Omega$.

\subsubsection{Espansione di Fourier}
Come conseguenza del fatto che abbiamo identificato nello spettro delle autofunzioni di $\h{K}$ una base completa e ortonormale per $\mathcal{L}^2(\Omega)$, possiamo esprimere una qualsiasi funzione d'onda $\Psi$ come combinazione lineare degli elementi di base. Sfruttando la relazione di completezza, abbiamo:
\[
\begin{aligned}
    \Psi(\bm{r}) &= \int_{\Omega} \dd^3{\bm{s}} \, \delta^{(3)}(\bm{s}-\bm{r}) \Psi(\bm{s}) \\
    &= \int_{\Omega} \dd^3{\bm{s}} \, \sum_{\bm{n}} \psi_{\bm{n}}^*(\bm{s}) \psi_{\bm{n}}(\bm{r}) \Psi(\bm{s}) \\
    &= \sum_{\bm{n}} \psi_{\bm{n}}(\bm{r}) \int_{\Omega} \dd^3{\bm{s}} \, \psi_{\bm{n}}^*(\bm{s}) \Psi(\bm{s}) = \sum_{\bm{n}} \psi_{\bm{n}}(\bm{r}) \braket{\psi_{\bm{n}}}{\Psi} = \sum_{\bm{n}} a_{\bm{n}} \psi_{\bm{n}}(\bm{r}) \quad,
\end{aligned}
\]
con $a_{\bm{n}}$ i coefficienti dello sviluppo di Fourier della funzione, definiti come di consueto. Possiamo ancora verificare l'identità di Parseval, essendo la base ortonormale:
\[
    \begin{aligned}
    \norm{\Psi}^2 = \braket{\Psi}{\Psi} &= \int_{\Omega} \dd^3{\bm{r}} \, \Psi(\bm{r})^* \Psi(\bm{r}) = \int_{\Omega} \dd^3{\bm{r}} \, \Psi(\bm{r})^* \sum_{\bm{n}} a_{\bm{n}} \psi_{\bm{n}}(\bm{r})\\
    &= \sum_{\bm{n}} a_{\bm{n}} \int_{\Omega} \dd^3{\bm{r}} \, \Psi(\bm{r})^*  \psi_{\bm{n}}(\bm{r}) = \sum_{\bm{n}} a_{\bm{n}} \braket{\Psi}{\psi_{\bm{n}}} = \sum_{\bm{n}} |a_{\bm{n}}|^2 \quad.
    \end{aligned}
\]

A livello formale, è necessario notare che tramnite la base delle onde piane stazionarie è possibile costruire qualsiasi stato $\Psi$ tramite serie di Fourier, anche se definito \emph{al di fuori} di $\Omega$. Questo ci porta a chiederci se la base in questione è effettivamente \emph{la migliore} per rappresentare una generica funzione d'onda. La risposta, abbastanza controintuitivamente, risiede nella geometria specifica del sistema e, per rendere ciò il più evidente possibile, analizzeremo nel paragrafo che segue un particolare caso in cui la base delle onde piane è quanto mai problematica. 

\subsection{Il caso della particella libera}
Benché la soluzione del caso della particella libera sia già stata proposta, è assolutamente istruttivo analizzare la base formata dagli autostati dell'hamiltoniana:
\[
    \mathcal{B}_{free} = \left\{ \phi_{\bm{k}}(\bm{r}) : \h{K}\phi_{\bm{k}}(\bm{r}) = \frac{\hbar^2k^2}{2m} \phi_{\bm{k}}(\bm{r}) \right\}, \quad \phi_{\bm{k}}(\bm{r}) = \frac{e^{i\bm{k}\cdot \bm{r}}}{(2\pi)^{\frac{3}{2}}}, \quad \bm{k} \in \mathbb{R}^3 \quad.
\]
Questa base corrisponde a quanto si troverebbe nel caso della scatola con $L \to \infty$ e dunque $\Omega \to \mathbb{R}^3$. Si noti che in questo caso non abbiamo più a che fare con onde puramente \emph{sinusoidali} bensì contiene anche una componente coseno... Chissà se questo fatto ci creerà problemi.

\subsubsection{Relazione di completezza}
Iniziamo questa volta con la valutazione della relazione di completezza perché ci tornerà utile in breve. Come nel caso della scatola, ricaviamo la relazione di completezza ricordando che, trattandosi di uno spettro intrinsecamente vettoriale, stiamo cercando una delta di Dirac tridimensionale:
\begin{equation}
        \delta^{(3)}(\bm{r}-\bm{s}) = \int_{\mathbb{R}^3} \dd^3{\bm{k}} \, \phi_{\bm{k}}(\bm{s})^* \phi_{\bm{k}}(\bm{r}) = \frac{1}{(2\pi)^{3}} \int_{\mathbb{R}^3} \dd^3{\bm{k}} \, e^{i\bm{k}\cdot (\bm{s}-\bm{r})} \quad,
        \label{eq:com_rel}
\end{equation} 
dove abbiamo integrato dato lo spettro continuo dei possibili valori di $\bm{k}$.

\subsubsection{Ortonormalità della base}
Questa condizione ci è in realtà garantita \emph{a priori} dal Teorema di Sturm-Liouville. Presentiamo comunque un calcolo concreto per covincere anche i più miscredenti:
\[
    \braket{\phi_{\bm{k}}}{\phi_{\bm{q}}} = \int_{\mathbb{R}^3} \dd^3{\bm{r}} \, \frac{e^{-i\bm{k}\cdot\bm{r}}}{(2\pi)^{\frac{3}{2}}} \frac{e^{i\bm{q}\cdot\bm{r}}}{(2\pi)^{\frac{3}{2}}} = \frac{1}{(2\pi)^3} \int_{\mathbb{R}^3} \dd^3{\bm{r}} \, e^{i\bm{k}\cdot (\bm{q}-\bm{k})} \overset{(1)}{=} \delta^{(3)}(\bm{k}-\bm{q}) \quad,
\]
dove in $(1)$ abbiamo riconosciuto la $\eqref{eq:com_rel}$. Questo "inusuale" risultato appena trovato, benché risponda correttamente alla condizione di ortogonalità quando $\bm{k} \neq \bm{q}$, esplode nel \emph{banale} caso della norma quadra:
\[
    \braket{\phi_{\bm{k}}}{\phi_{\bm{k}}} = \int_{\mathbb{R}^3} \frac{\dd^3{\bm{r}}}{(2\pi)^3} \to \infty \quad.
\]
Le onde piane, come avevamo già scoperto, \textbf{non sono normalizzabili} in $\mathbb{R}^3$! La base $\mathcal{B}_{free}$ non è dunque ortonormale. In ogni caso, questo caso patologico non dovrebbe apparire "strano", essendo $|\phi_{\bm{k}}|^2 = (2\pi)^{-3}$, quindi costante in ogni punto di $\mathbb{R}^3$. 

Questo, apparente, enorme problema non deve però essere visto come una limitazione della teoria vista sin'ora, bensì come una conseguenza intrinseca della scelta di ambientare la nostra particella in uno spazio \emph{illimitato}. In questo contesto ritorna quanto accennato nel caso della scatola riguardo come geometria del sistema influenzi, enormemente, la comodità della base scelta per rappresentari gli stati. Addirittura siamo riusciti a "trasformare" una comoda e semplice base ortonormale, in una mostruosità \emph{ortogonale nel senso della delta di Dirac}, con elementi neppure normalizzabili. 

\subsubsection{Interpretazione probabilistica}
L'interpretazione di Copenhagen è ancora valida: è semplicemente necessario rivederla. Date due funzioni d'onda normalizzate $\phi_1(\bm{r})$ e $\phi_2(\bm{r})$, definiamo la probabilità di trovare la particella entro i due volumi $\Omega_1, \Omega_2 \in \mathbb{R}^3$ come di consueto:
\[
\begin{aligned}
    \mathcal{P}(\Omega_1) &= \int_{\Omega_1} \dd^3{\bm{r}} \, |\phi_1|^2 = \frac{\text{Volume}\,(\Omega_1)}{(2\pi)^3} \\
    \mathcal{P}(\Omega_2) &= \int_{\Omega_2} \dd^3{\bm{r}} \, |\phi_2|^2 = \frac{\text{Volume}\,(\Omega_2)}{(2\pi)^3}
\end{aligned}
\implies
\frac{\mathcal{P}(\Omega_1)}{\mathcal{P}(\Omega_2)} = \frac{\text{Volume}\,(\Omega_1)}{\text{Volume}\,(\Omega_2)} \quad,
\]
dove abbiamo infine suggerito una nozione di \textbf{densità relativa di probabilità}.

\subsubsection{Espansione di Fourier}
Il fatto che $\mathcal{B}_{free}$ sia una base per $\mathcal{L}^2$ garantisce che una qualsiasi funzione d'onda possa essere espressa tramite gli elementi della base stessa. Sfruttando la notazione integrale essendo lo spettro dei vari $\bm{k}$ continuo:
\[
    \Phi(\bm{r}) = \int_{\mathbb{R}^3} \dd^3{\bm{k}} \, a(\bm{k})\phi_{\bm{k}}(\bm{r})
\]
e, sfruttando la relazione di completezza, otteniamo l'espressione per i coefficienti $a(\bm{k})$:
\[
    \begin{aligned}
        \Phi(\bm{r}) = \int_{\mathbb{R}^3} \dd^3{\bm{s}} \, \Phi(\bm{s}) \delta^{(3)}{(\bm{s}-\bm{r})} &= \int_{\mathbb{R}^3} \dd^3{\bm{s}} \, \Phi(\bm{s}) \int_{\mathbb{R}^3} \dd^3{\bm{k}} \, \phi_{\bm{k}}(\bm{s})^*\phi_{\bm{k}}(\bm{r}) \\
        &= \int_{\mathbb{R}^3} \dd^3{\bm{k}} \, \phi_{\bm{k}}(\bm{r}) \underbrace{\int_{\mathbb{R}^3} \dd^3{\bm{s}} \, \Phi(\bm{s}) \phi_{\bm{k}}(\bm{s})^*}_{a(\bm{k})} \quad.
        \end{aligned}
\]
Possiamo, in utltimo, verificare se l'identità di Parseval regge anche in questo caso, nonostante la base analizzata non sia ortonormale\footnote{Attenzione! L'identità di Parseval vale solo quando si tratta con basi \emph{ortonormali}, condizione evidentemente non soddisfatta da $\mathcal{B}_{free}$. Come però si vede dal calcolo, è rigorosamente consentito estendere l'identità al continuo tramite il \textbf{Teorema di Plancherel}.}:
\[
\begin{aligned}
    \braket{\phi_{\bm{k}}}{\phi_{\bm{k}}} &= \int_{\mathbb{R}^3} \dd^3{\bm{r}} \, \phi_{\bm{k}}(\bm{r})^*\phi_{\bm{k}}(\bm{r}) \\
    &= \int_{\mathbb{R}^3} \dd^3{\bm{r}} \, \underbrace{\left(\int_{\mathbb{R}^3} \dd^3{\bm{k}} \, a(\bm{k})\phi_{\bm{k}}(\bm{r})\right)^*}_{\phi_{\bm{k}}(\bm{r})^*} \underbrace{\int_{\mathbb{R}^3} \dd^3{\bm{q}} \, a(\bm{q})\phi_{\bm{q}}(\bm{r})}_{\phi_{\bm{k}}(\bm{r})} \\
    &= \underbrace{\int_{\mathbb{R}^3} \dd^3{\bm{r}} \, \phi_{\bm{k}}(\bm{r})^* \phi_{\bm{q}}(\bm{r})}_{\delta^{(3)}({\bm{k}-\bm{q}})} \int_{\mathbb{R}^3} \dd^3{\bm{k}} \, a(\bm{k})^* \int_{\mathbb{R}^3} \dd^3{\bm{q}} \, a(\bm{q}) = \int_{\mathbb{R}^3} \dd^3{\bm{k}} \, |a(\bm{k})|^2 \quad. \\
\end{aligned}
\]

\section{Alcuni valori di aspettazione}
Siamo pronti per immergerci nella computazione di alcuni valori di aspettazione basilari. Quando abbiamo definito cosa significasse la notazione $ \langle \h{A} \rangle $, abbiamo proposto come unico esempio fondamentale per convincerci del fatto che \emph{ha un significato sufficientemente classico}: quello dell'operatore posizione. La conoscenza delle basi e dell'espansione in Fourier di un generico stato $\psi$ ci permette ora di fare conti più complessi come, ad esempio, il calcolo del valore di aspettazione della quantità di moto, dell'energia cinetica e dell'hamiltoniano. 

\subsection{Valore di aspettazione dell'operatore quantità di moto}
Immergiamoci nell'algebra. Preso un corpuscolo in un volume $\Omega \in \mathbb{R}^3$ descritto dalla funzione d'onda $\psi(\bm{r})$\footnote{Per semplicità trascuriamo la parte temporale in quanto contribuisce unicamente in un termine di fase $e^{-i\frac{E}{\hbar}t}$} espressa nella base delle onde piane che diagonalizzano $\hb{p}$ ($\mathcal{B}_{\hb{p}}$), abbiamo dalla definizione:
\[
    \begin{aligned}
        \avg{\hat{\bm{p}}} = (\psi, \hat{\bm{p}} \psi) = \braket{\psi}{\hb{p}\psi} &= \int_\Omega \dd^3{\bm{r}} \, \psi(\bm{r})^*\hb{p}\psi(\bm{r}) \\
        &\overset{(1)}{=} \int_\Omega \dd^3{\bm{r}} \, \psi(\bm{r})^* \hb{p} \int_{\mathbb{R}^3} \dd^3{\bm{k}} \, a(\bm{k}) \phi_{\bm{k}}(\bm{r}) = \int_\Omega \dd^3{\bm{r}} \, \psi(\bm{r})^* \int_{\mathbb{R}^3} \dd^3{\bm{k}} \, a(\bm{k}) \hb{p} \phi_{\bm{k}}(\bm{r}) \\
        &\overset{(2)}{=} \int_\Omega \dd^3{\bm{r}} \, \psi(\bm{r})^* \int_{\mathbb{R}^3} \dd^3{\bm{k}} \, a(\bm{k}) \hbar \bm{k} \phi_{\bm{k}}(\bm{r}) \\
        &= \int_{\mathbb{R}^3} \dd^3{\bm{k}} \, a(\bm{k}) \hbar \bm{k} \underbrace{\int_\Omega \dd^3{\bm{r}} \, \psi(\bm{r})^* \phi_{\bm{k}}(\bm{r})}_{\braket{\psi}{\phi_{\bm{k}}}^* = a(\bm{k})^*} = \int_{\mathbb{R}^3} \dd^3{\bm{k}} \, |a(\bm{k})|^2 \hbar \bm{k} \quad,
    \end{aligned}
\]
dove in $(1)$ abbiamo espanso $\psi$ in Fourier e in $(2)$ abbiamo commuatto l'operatore quantità di moto, agente sulle coordinate spaziali, con l'integrale sugli impulsi. 

Specifichiamo che la base in questione è quelle delle onde piane \emph{progressive}:
\[
   \mathcal{B}_{\hb{p}} = \left\{ \phi_{\bm{k}}(\bm{r}) = \frac{e^{i\bm{k}\cdot\bm{r}}}{(2\pi)^{\frac{3}{2}}} \right\} \quad.
\]

\subsubsection{Interpretazione fisica del risultato}
Si ricordi quanto avevamo trovato nel caso dell'operatore posizione: la perfetta corrispondenza quantomeccanica del concetto di centro di massa, ovvero la somma ponderata di tutte le possibili posizioni $\bm{r}_i$ assumibili dalla particella soppesate dalla probabilità di trovarla effettivamente in $\bm{r}_i$. Nel precedente calcolo rivediamo lo stesso schema: il valore di aspettazione di $\hb{p}$ corrisponde alla "media" dei possibili momenti $\bm{p}_i = \hbar \bm{k}_i$ associati al corpuscolo pesata dalla probabilità che, all'atto della misurazione, $\psi$ si trovi nello stato $\phi_{\bm{k}}$\footnote{Il concetto appena introdotto, meglio noto come \textbf{collasso della funzione d'onda}, verrà ripreso nel discutere la teoria della misura.}. 

In questo senso $|a(\bm{k})^2|$ è uguale alla probabilità che $\psi$ assuma come momento $\bm{p}_i$.

\subsection{Valore di aspettazione dell'operatore energia cinetica}
Nelle stesse ipotesi di prima (base derivante dalla diagonalizzazione di $\h{K}$ ($\mathcal{B}_{\h{K}}$), volume $\Omega \in \mathbb{R}^3$, trascuriamo la parte temporale), calcoliamo:
\[
    \begin{aligned}
        \avg{\h{K}} = (\psi, \h{K} \psi) = \braket{\psi}{\h{K}\psi} &= \int_\Omega \dd^3{\bm{r}} \, \psi(\bm{r})^*\h{K}\psi(\bm{r}) \\
        &= \int_\Omega \dd^3{\bm{r}} \, \psi(\bm{r})^* \h{K} \int_{\mathbb{R}^3} \dd^3{\bm{k}} \, a(\bm{k}) \phi_{\bm{k}}(\bm{r}) \\
        &= \int_\Omega \dd^3{\bm{r}} \, \psi(\bm{r})^* \int_{\mathbb{R}^3} \dd^3{\bm{k}} \, a(\bm{k}) \h{K} \phi_{\bm{k}}(\bm{r}) \\
        &= \int_\Omega \dd^3{\bm{r}} \, \psi(\bm{r})^* \int_{\mathbb{R}^3} \dd^3{\bm{k}} \, a(\bm{k}) \underbrace{\frac{\hbar^2 k^2}{2m}}_{E_{\mathcal{K}}} \phi_{\bm{k}}(\bm{r}) \\
        &= \int_{\mathbb{R}^3} \dd^3{\bm{k}} \, a(\bm{k}) E_{\mathcal{K}} \int_\Omega \dd^3{\bm{r}} \, \psi(\bm{r})^* \phi_{\bm{k}}(\bm{r}) = \int_{\mathbb{R}^3} \dd^3{\bm{k}} \, |a(\bm{k})|^2 E_{\mathcal{K}} \quad,
    \end{aligned}
\]
nella base delle onde piane (qualsiasi combinazione lineare sarebbe ammessa ma si scelgono unicamente le onde progressive, così come per la quantità di moto):
\[
   \mathcal{B}_{\h{K}} = \left\{ \phi_{\bm{k}}(\bm{r}) = \frac{e^{i\bm{k}\cdot\bm{r}}}{(2\pi)^{\frac{3}{2}}} \right\} \quad.
\]
L'interpretazione è in questo contesto la medesima di quella proposta nel caso della quantità di moto.

\subsection{Valore di aspettazione dell'operatore hamiltoniano}
In questo caso dobbiamo prestare più attenzione alla definizione di "autofunzioni" per l'operatore $\h{H}$. Nella soluzione della S.E. stazionaria, ci eravamo ritrovati a risolvere l'equazione agli autovalori per l'operatore hamiltoniano, ottenendo un'insieme di funzioni d'onda che ora sappiamo essere una base per $\mathcal{L}^2$. Supponendo lo spettro di energia come discreto, abbiamo $\mathcal{B}_{\h{H}} = \{ \psi_{E_n} \}$ e possiamo finalmente portare a termine il conto:
\[
    \begin{aligned}
        \avg{\h{H}} = (\Psi, \h{H} \Psi) = \braket{\Psi}{\h{H}\Psi} &= \int_\Omega \dd^3{\bm{r}} \, \Psi(\bm{r})^*\h{H}\Psi(\bm{r}) \\
        &= \int_\Omega \dd^3{\bm{r}} \, \Psi(\bm{r})^* \h{H} \sum_j C(E_n) \psi_{E_n}(\bm{r}) \\
        &= \sum_n C(E_n)\int_\Omega \dd^3{\bm{r}} \, \Psi(\bm{r})^* \h{H} \psi_{E_n}(\bm{r}) \\
        &= \sum_n C(E_n)\int_\Omega \dd^3{\bm{r}} \, \Psi(\bm{r})^* E_n \psi_{E_n}(\bm{r}) \\
        &= \sum_nC(E_n) E_n \underbrace{\int_\Omega \dd^3{\bm{r}} \, \Psi(\bm{r})^* \psi_{E_n}(\bm{r})}_{\braket{\Psi}{\psi_{E_n}}^* = C(E_n)^*} = \sum_n|C(E_n)|^2 E_n \quad.
    \end{aligned}
\]
Anche in questo caso, abbiamo implicitamente sfruttato il fatto che le autofunzioni dell'operatore hamiltoniano formano una base completa per $\mathcal{L}^2$. Benché una giustficazione possa venire dal teorema di Sturm-Liouville, dimostreremo questo fatto nella prossima sezione in maniera più dettagliata. L'importante da notare è quanto ottenuto è ancora una volta la media ponderata, \emph{quantistica}, dei possibili valori di energia del sistema. 

Si noti che, anche considerando l'evoluzione temporale del generico sistema $\Psi (\bm{r},t)$, il risultato sarebbe stato indipendente dal tempo, mimando la stessa indipendenza classica della funzione di Hamilton.
%!TEX root = ../dispense_QP.tex

\chapter{Proprietà degli operatori hermitiani e equazioni agli autovalori}
Come intuito nel precedente capito, non abbiamo ancora tutte le necessarie informazioni sugli operatori hermitiani per riuscire a rispondere a dovere ad alcune domande sorte nel corso della discussione come, ad esempio, perché l'insieme formato dagli autostati di un operatore hermitiano sia una \emph{base completa} per $\mathcal{L}^2$. Daremo ora altre informazioni sui nostri tanto amati (e odiati) operatori e, sperabilmente, riusciremo a venire a capo dei nostri problemi.

\section{Proprietà dello spettro di un operatore hermitiano}
Partiamo ora da un presupposto ormai appurato e deriviamo alcune considerazioni. Gli operatori fisici \emph{devono} essere hermitiani... ma quali sono le proprietà algebriche di questi fantomatici operatori autoaggiunti? Beh, due fra tutte, sono contraddistinti da autovalori \textbf{reali} e autostati (autofunzioni) \textbf{ortogonali} fra loro. 

\begin{teorema}[Proprietà degli autovalori di un operatore hermitiano]
    Sia $\h{A}$ un operatore hermitiano (ovvero per cui vale $\h{A} = \h{A}^\dagger$). Allora lo spettro degli autovalori $\alpha_m$ di $\h{A}$ è reale: 
    \[
        \alpha_m \in \mathbb{R}, \quad \forall m \quad.
    \]
\end{teorema}
\begin{proof}
    L'equazione agli autovalori per l'operatore $\h{A}$, dette $\psi_m$ le sue autofunzioni e $\alpha_m$ gli autovalori, risulta:
    \[
        \h{A}\psi_m = \alpha_m \psi_m \quad.
    \]
    Calcoliamo ora il prodotto scalare tra $\psi_m$ e $\h{A}\psi_m$:
    \[
        (\psi_m, \h{A}\psi_m) = (\psi_m, \alpha_m\psi_m) \overset{(1)}{=} \alpha_m(\psi_m,\psi_m)
    \]
    dove in $(1)$ abbiamo usato il fatto che il prodotto scalare in $\mathcal{L}^2$ è sesquilineare. Essendo $\h{A}$, hermitiano, possiamo ancora scrivere:
    \[
    \begin{aligned}
        (\psi_m, \h{A}\psi_m) &= (\h{A}\psi_m, \psi_m) \\
        &= \alpha_m^*(\psi_m, \psi_m) \quad.
    \end{aligned}
    \]
    Sottraendo memebro a membro le due diferse formulazioni di $(\psi_m, \h{A}\psi_m)$ otteniamo infine:
    \[
    \begin{aligned}
        (\psi_m, \h{A}\psi_m) - (\psi_m, \h{A}\psi_m) &= \alpha_m(\psi_m, \psi_m) - \alpha_m^*(\psi_m, \psi_m) \\ &= (\alpha_m -\alpha_m^*)(\psi_m, \psi_m) =(\alpha_m -\alpha_m^*) \norm{\psi_m}^2 = 0 \quad.
    \end{aligned}
    \]
    Dalla legge di annullamento del prodotto abbiamo due possibili soluzioni all'ultima equazione:
    \begin{itemize}
        \item $\norm{\psi_m}^2 = 0$: essendo la norma una funzione definita positiva, resitutisce zero se e solo se il suo argomento è la funzione nulla stessa. Il vettore nullo non può però essere un autovettore, da cui serve escludere il caso $\psi_m = 0$. 
        \item $\alpha_m = \alpha_m^*$: unica possibile soluzione che implica che $\alpha_m \in \mathbb{R}$.
    \end{itemize} 
\end{proof}

Derivata questa fondamentale nozione, possiamo finalmente approcciarci all'analisi del prossimo risultato. 
\begin{teorema}[Proprietà degli autofunzioni di un operatore hermitiano]
    Sia $\h{A}$ un operatore hermitiano (ovvero per cui vale $\h{A} = \h{A}^\dagger$). Allora $\h{A}$ ammette uno spettro di autostati $\psi_m$ ortogonali fra loro: 
    \[
        (\psi_m,\psi_n) = 0, \quad \forall \, m,n \quad.
    \]
\end{teorema}
Prima di dimostrare il teorema, serve però distinguere due casi. Detti $\alpha_m$ gli autovalori di $\h{A}$ possiamo infatti avere: $(1)$ $\alpha_m \neq \alpha_n \, \forall\,  m,n$ (\textbf{caso non degenere}); $(2)$ $\exists \, m,n : \alpha_m = \alpha_n$ (\textbf{caso degenere}). Dimostreremo per primo il caso non degenere per poi addentrarci nella discussione quello degenere.

\begin{proof}
    (\textbf{1. Caso non degenere}). Assumiamo di avere uno spettro di autovalori tutti diveri fra loro. Allora, dette $\psi_m$ e $\psi_n$ due autofunzioni di $\h{A}$, abbiamo le seguenti equazioni agli autovalori:
    \[
        \begin{aligned}
            &\h{A} \psi_m = \alpha_m\psi_m \\
            &\h{A} \psi_n = \alpha_n\psi_n \quad.
        \end{aligned}
    \]
    Continuiamo calcolando i seguenti prodotti scalari:
    \[
    \begin{aligned}
        &(\psi_m, \h{A}\psi_n) = \alpha_n (\psi_m,\psi_n) \\
        &(\psi_m, \h{A}\psi_n) = (\h{A}\psi_m, \psi_n) = \alpha_m^*(\psi_m, \psi_n) \overset{(1)}{=} \alpha_m(\psi_m, \psi_n)
    \end{aligned}
    \]
    dove in $(1)$ abbiamo sfruttato il fatto che gli autovalori dell'operatore sono reali. Abbiamo quindi la seguente uguaglianza:
    \[
        \alpha_m(\psi_m, \psi_n) = \alpha_n(\psi_m, \psi_n) \implies (\alpha_m - \alpha_n) (\psi_m, \psi_n) = 0 \quad.
    \]  
    Avendo assunto autovalori distinti e quindi $\alpha_m \neq \alpha_m$, serve che il prodotto scalare sia identicamente nullo e che quindi $\psi_m$ sia ortogonale a $\psi_n$.  
\end{proof}

\begin{proof}
    (\textbf{2. Caso degenere}). Sappiamo adesso che, per ipotesi, esistono autovalori uguali associati ad autostati diversi. Senza perdita di generalità, assumiamo che $\alpha_1 = \alpha_2 = \dots = \alpha_m \neq \alpha_{m+1} \dots$, ovvero che i primi $m$ autovalori siano uguali fra loro e diversi dai successivi. Definiamo l'insieme $\tilde{\mathcal{B}}$ come il contenente delle $m$ autofunzioni associati ai suddetti $m$ autovalori:
    \[
        \tilde{\mathcal{B}} \coloneq \left\{ \psi_i \in \mathcal{L}^2, i \in [1, m] : \h{A}\psi_i = \alpha_i\psi_i \right\} \quad.
    \]
    Chiamiamo ora $c_{ab}$ il prodotto scalare fra due elementi di $\tilde{\mathcal{B}}$:
    \[
    c_{ab} = (\psi_a, \psi_b), \quad \psi_a, \psi_b \in \tilde{\mathcal{B}} \quad.
    \]
    Per costruzione, deve valere $c_{ab} = c_{ba}^*$, infatti abbiamo che $(\psi_a, \psi_b) = (\psi_b, \psi_a)^*$. Se raggruppiamo i cavi $c_{ab}$ in una matrice $\mathcal{C}$, quest'ultima risulta essere quadrata ($\mathcal{C} \in \mathbb{M}^{m\times m}$) e \emph{simmetrica nel senso hermitiano} (o banalmente hermitiana), infatti:
    \[
    \begin{aligned}
        \mathcal{C} &= \mathcal{C}^\dagger \\
        \implies ({\mathcal{C}})_a^b = c_{ab} &= c_{ba}^* = ({\mathcal{C}^\dagger})_b^a \quad.
    \end{aligned}
    \]  
    Il \textbf{Teorema spettrale} ci garantisce che le matrici hermitiane sono sempre diagonalizzabili in campo reale e ammettono matrici diagonalizzanti ortogonali $\mathcal{O}$ e $\mathcal{O}^{\dagger}$: 
    \[
        \mathcal{D} = \mathcal{O}^\dagger \mathcal{C} \mathcal{O} 
    \]
    con $\mathcal{D}$ diagonale. In componenti, l'uguaglianza assume la forma di:
    \[
        \alpha_i \delta_{ij} = \sum_a \sum_b (\mathcal{O}^\dagger)_i^a (\mathcal{C})_a^b (\mathcal{O})_b^j \quad.
    \]
    Data la definizione del termine $[a,b]$ di $\mathcal{C}$, riarrangiamo i termini e calcoliamo:
    \[
        \begin{aligned}
            \alpha_i \delta_{ij} &= \sum_a \sum_b (\mathcal{O}^\dagger)_i^a (\psi_a,\psi_b) (\mathcal{O})_b^j \overset{(1)}{=} \sum_a \sum_b [(\mathcal{O})_i^a]^* (\psi_a,\psi_b) (\mathcal{O})_b^j \\
            &= \sum_a \sum_b ((\mathcal{O})_i^a\psi_a,(\mathcal{O})_b^j\psi_b) = \left( \sum_a(\mathcal{O})_i^a\psi_a, \sum_b (\mathcal{O})_b^j\psi_b \right) = (\phi'_i, \phi'_j) \quad. \\
        \end{aligned}
    \]
    dove in $(1)$ abbiamo sfruttato la proprietà di simmetria hermitiana della matrie diagonalizzante mentre i nuovi stati $\psi'_{i,j}$ sono stati costruiti come \emph{combinazione lineare} degli elementi di $\tilde{\mathcal{B}}$. Dividendo per $\alpha_i$ ottendiamo una utile relazione di ortonormalità:
    \[
       \frac{(\phi'_i, \phi'_j)}{\alpha_i} = \left( \frac{\phi'_i}{\sqrt{\alpha_i}}, \frac{\phi'_j}{\sqrt{\alpha_i}} \right) = (\phi_i, \phi_j) = \delta_{ij}, \quad
       \begin{dcases}
            \phi'_i \to \phi_i = \frac{\phi'_i}{\sqrt{\alpha_i}} \\
            \phi'_j \to \phi_j = \frac{\phi'_j}{\sqrt{\alpha_i}} \quad.
       \end{dcases}
    \]
    Possiamo ora costruire l'insieme:
    \[
        \mathcal{B} \coloneq \left\{ \phi_j, j \in [1, m]; \psi_i, i \in [m+1, \infty) \right\}
    \]
    che risulta evidentemente formato da elementi ortogonali fra loro\footnote{Si noti infatti che i vari $\phi_j$ sono costruiti a partire dagli autostati $\psi_j$ che condividevano gli stessi autovalori. Dal momento che i $\psi_i, i \in [m+1, \infty)$ hanno per ipotesi autovalori diversi rispeti agli $\psi_j, j \in [1, m]$, per la dimostrazione del caso non degenere abbiamo $\psi_i \perp \psi_j \implies \psi_i \perp \phi_j$.}.
\end{proof}

In generale, quindi, è possibile dimostrare che un insieme formato da autostati ortogonali associati ad un dato operatore hermitiano rappresenta una \emph{base} per $\mathcal{L}^2(\Omega)$ (per un determinato volume $\Omega$) e, pertanto, è un set \emph{completo}. 

\subsection{Interpretazione geometrica di alcuni risultati fisici}
Il secondo teorema visto ci permette di far evolvere in senso \emph{geometrico} un'intuizione derivante dalla discussione della soluzione generale della S.E.: un qualsiasi stato $\Psi$ può essere espresso sfruttando il Principio di sovrapposizione applicato alle infinite soluzione dell'equazione di Schrödinger, gli \emph{stati stazionari} $\psi_{E_n}$, modulate per un fattore di fase associato all'evoluzione temporale. Adesso capiamo che gli stati stazionari non sono altro che gli autostati dell'operatore hamiltoniano\footnote{Infatti avevamo assunto che fossero le funzioni d'onda che diagonalizzano $\h{H}$.} che, in quanto tali, formano una base \emph{discreta} per $\mathcal{L}^2(\Omega)$ tramite cui è così possibile descrivere una generica $\Psi$ tramite serie di Fourier. 

\subsubsection{Il principio di sovrapposizione}
Per un generico operatore $\h{A}$ agente su $\mathcal{L}^2(\Omega)$, con spettro discreto, abbiamo che una qualsiasi $\Psi \in \mathcal{L}^2(\Omega)$ può essere scritta come:
\[
    \Psi = \sum_n \beta_n\psi_n, \quad \beta_n = \braket{\psi_n}{\Psi}, \quad  \h{A} \psi_n = \alpha_n \psi_n \quad.
\]
Ancora, essendo $\h{A}$ hermitiano e quindi la base dei suoi autostati ortonormale, abbiamo l'identità di Parseval:
\[
    \braket{\Psi}{\Psi} = \sum_n |\beta_n|^2 = 1 \quad,
\]
giustificata dalla condizione di normalizzazione. In questo senso, vale la \textbf{regola di Born generalizzata} e $|\beta_n|^2$ rappresenta la probabilità che il sistema si trovi nello stato $\beta_n \to \psi_n$.

\subsubsection{Il valore di aspettazione di un operatore}
Cerchiamo ora di capire cosa significa trovare il valore di aspettazione di un operatore $\h{A}$. Supponendo inizialmente che abbia uno spettro discreto di autostati $\psi_n$ e autovalori $\alpha_n$, otteniamo dalla definizione:
\[
\begin{aligned}
    \avg{\h{A}} = \braket{\Psi}{\h{A}\Psi} = \braket{\sum_n \beta_n\psi_n}{\h{A}\sum_m \beta_m\psi_m} &= \sum_n \sum_m \beta_n^* \beta_m \braket{\psi_n}{\h{A}\psi_m} \\
    &= \sum_n \sum_m \beta_n^* \beta_m \braket{\psi_n}{\alpha_m\psi_m}\\
    &= \sum_n \sum_m \beta_n^* \beta_m \alpha_m \delta_{mn} = \sum_n |\beta_n|^2 \alpha_n \quad.
\end{aligned}
\]
Il valore di aspettazione di un operatore, che non è altro che il valore atteso a seguito di una misura effettuata sulla grandezza fisica associata ad $\h{A}$, è dunque la somma ponderata degli autovalori dell'operatore stesso pesata dalla probabilità che il generico stato misurato si trovi nello stato $\psi_n$. Si presti ben attenzione che il valore atteso di una misura $\neq$ il risultato di una misurazione che assumeremo corrispondere a \emph{uno degli autovalori dell'operatore}!

Ancora, abbiamo dal primo teorema discusso in precedenza, il fatto che necessariamente $\avg{\h{A}} \in \mathbb{R}$, in piena coerenza con l'interpretazione \emph{misurista} suggerita in precedenza. In ultimo, la relazione di completezza risulta:
\[
    \delta^{(3)}(\bm{s}-\bm{r}) = \sum_n \psi_n(\bm{s})^*\psi_n(\bm{r}) \quad.
\]
La discussione riguardo operatori con spettro continuo è assolutmente affine e non sarà condotta dettagliatamente.

\section{Gli assiomi della teoria quantistica}
Siamo adesso pronti a enumerare gli assiomi su cui si basa la teoria quantistica. 
\begin{enumerate}[label=\roman*.]
    \item Gli stati fisici di un sistema sono descritti da \textbf{funzioni d'onda} $\psi(\bm{r},t)$ e vale  $ \psi: \mathbb{R} \to \mathbb{C}$.
    \item Le osservabili fisiche (grandezze) corrispondono a operatori hermitiani.
    \item Il valore di aspettazione di un generico operatore $\h{A}$ è dato da:
        \[
            \avg{\h{A}} = \braket{\psi}{\h{A}\psi} = \int_\Omega \dd^3{\bm{r}} \psi(\bm{r})^*\h{A}\psi(\bm{r}) \quad.
        \]
    \item L'evoluzione temporale degli stati fisici è governata dall'equazione di Schrödinger:
    \[
        i\hbar \pdv{\psi(\bm{r},t)}{t} = \h{H}\psi(\bm{r},t) \quad.
    \]
    \item (\textbf{Principio del collasso di una funzione d'onda}). La misurazione di una grandezza fisica $\mathcal{A}$ a cui è associato un operatore hermitiano $\h{A}$ di autovalori $\alpha_m$ ha come risultato un $\alpha_m$. A seguito della misura, lo stato fisico si trova nell'autostato associato all'autovalore misurato $\psi \to \psi_m$.
\end{enumerate}
%!TEX root = ../dispense_QP.tex

\chapter{Il caso dell'oscillatore armonico}
Siamo ora pronti ad analizzare uno dei pilastri su cui si fonda non solo la fisica classica, bensì anche la meccanica quantistica: il caso dell'\textbf{oscillatore armonico}. Le potenzialità che contraddistinguono questo concetto sono pressochè infinite, dall'analisi di semplici moti armonici alla discussione sulla quantizzazione del campo magnetico. 

L'hamiltoniana del sistema in una dimensione\footnote{La restrizione unidimensionale non toglie generalità al sistema, eventualmente estendibile al caso tridimensionale.}, classicamente, risulta:
\[
    \mathcal{H} = \frac{{p}^2}{2m} + \frac{1}{2}m\omega^2x^2 \quad, 
\] 
con $\omega$ la pulsazione del sistema e $m\omega^2$ la sua \emph{costante elastica}. Effettivamente, un sistema con un potenziale dalla forma funzionale generica può essere ben approssimato, nell'intorno di un suo minimo (ovvero un punto di equilibrio stabile), da un polinomio di secondo grado:
\[
\begin{aligned}
    U(x) &= U(x_0) + \underbrace{U'(x_0)(x-x_0)}_{\text{0 nel punto di equilibrio}} + \frac{1}{2}U''(x_0)(x-x_0)^2 + o(x-x_0)^2 \\
    &\approx U(x_0) + \frac{1}{2}U''(x_0)(x-x_0)^2 \quad.
\end{aligned}
\]
È dunque fondamentale proporre una riformulazione quantistica che, per di più, è sufficientemente banale. Assumendo conto di un hamiltoniana formata dall'operatore energia cinetica e dall'operatore energia potenziale elastica. La S.E. stazionaria risulta:
\[
    \h{H}\psi(x) = \left( -\frac{\hbar^2}{2m}\dv[2]{}{x} + \frac{1}{2}m\omega^2x^2 \right)\psi(x) = E\psi(x) \quad.
\]
La soluzione di questa equazione differenziale è particolarmente complicata è coinvolge la \emph{teoria dei polinomi ortogonali di Hermite}. Noi sfrutteremo un metodo non puramente algebrico, detto \textbf{metodo dell'algebra generatrice dello spettro}, che ci semplificherà notevolmente la vita... il risultato che otterremo sarà comunque identico a quello previsto dalla teoria dei polinomi di Hermite, della quale daremo una breve introduzione al termine di questo capitolo.

\section{Il metodo dell'algebra generatrice dello spettro}
L'idea che implementeremo è quella di sfruttare le proprietà algebriche dell'hamiltoniana dell'oscillatore armonico per definire due nuovi operatori: l'\textbf{operatore di creazione} $\hat{a}$ e l'\textbf{operatore di annichilazione} $\hat{a}^\dagger$. Nell'algebra di base, la somma di due quadrati è fattorizzabile come segue:
\begin{equation}
    x^2 + y^2 = x^2 - (iy)^2 = (x+iy)(x-iy) \quad,
    \label{eq:somma_quadrati}
\end{equation}
classica e intramontabile regola della \emph{differenza di quadrati = somma per differenza}. Cerchiamo ora di fare lo stesso con la nostra hamiltoniana: gli operatori coinvolti, posizione e momento, si trovano infatti al grado 2 ma, in generale, dipendenze lineari sono più semplici da analizzare. Prima di partire, notiamo che il prodotto fra operatori, in generale, non è commutativo... dobbiamo riformulare la \eqref{eq:somma_quadrati}:
\begin{equation}
\begin{aligned}
    (\h{A} - i\h{B})(\h{A} + i\h{B}) = \h{A}^2 + \h{B}^2 - i\h{B}\h{A} + i\h{A}\h{B} = \h{A}^2 + \h{B}^2 +i[\h{A},\h{B}] \\
    \implies \h{A}^2 + \h{B}^2 = (\h{A} - i\h{B})(\h{A} + i\h{B}) - i[\h{A},\h{B}] \quad.
\end{aligned}
\label{eq:somma_diff_operatori}
\end{equation}
Vediamo adesso cosa succede provando fattorizzare $\h{H}$, assumendo $\hat{p} = \hat{p}_x$:
\[
\begin{aligned}
    \h{H} = \frac{\hat{p}^2}{2m} + \frac{1}{2}m\omega^2\hat{x}^2 &=  \frac{m\omega^2}{2}\left( \hat{x}^2 + \frac{\hat{p}^2}{m^2\omega^2}  \right) \\
    &\overset{(1)}{=}  \frac{m\omega^2}{2}\left[\left( \hat{x} - i\frac{\hat{p}}{m\omega}  \right)\left( \hat{x} + i\frac{\hat{p}}{m\omega}  \right)-i \frac{[\hat{x}, \hat{p}]}{m\omega}\right] \\
    &= \frac{m\omega^2}{2}\left[\left( \hat{x} - i\frac{\hat{p}}{m\omega}  \right)\left( \hat{x} + i\frac{\hat{p}}{m\omega}  \right)+ \frac{\hbar}{m\omega}\right] \\
    &= \omega\hbar\bigg[\frac{m\omega}{2\hbar}\left( \hat{x} - i\frac{\hat{p}}{m\omega}\right)   \left( \hat{x} + i\frac{\hat{p}}{m\omega}  \right)+ \frac{1}{2}\bigg] \\
    &= \omega\hbar\bigg[\underbrace{\sqrt{\frac{m\omega}{2\hbar}}\left( \hat{x} - i\frac{\hat{p}}{m\omega}  \right)}_{\hat{a}^\dagger} \underbrace{\sqrt{\frac{m\omega}{2\hbar}}\left(\hat{x} + i\frac{\hat{p}}{m\omega}  \right)}_{\hat{a}}+ \frac{1}{2}\bigg] \quad.
\end{aligned}
\]
dove in $(1)$ abbiamo sfruttato la \eqref{eq:somma_diff_operatori}. Definiamo così gli operatori di creazione e annichilazione:
\[
\begin{dcases}
    \hat{a} = \sqrt{\frac{m\omega}{2\hbar}}\left(\hat{x} + i\frac{\hat{p}}{m\omega}\right) \\
    \hat{a}^\dagger = \sqrt{\frac{m\omega}{2\hbar}}\left(\hat{x} - i\frac{\hat{p}}{m\omega} \right)
\end{dcases}
\implies 
\begin{dcases}
    \hat{x} = \sqrt{\frac{\hbar}{2m\omega}}(\hat{a}+ \hat{a}^\dagger) \\
    \hat{p} = -i\sqrt{\frac{m\omega\hbar}{2}}(\hat{a} - \hat{a}^\dagger) 
\end{dcases}
\]
con l'hamiltoniana assume la forma finale di:
\[
    \h{H} = \omega \hbar\left( \hat{a}^\dagger\hat{a} + \frac{1}{2} \right) \quad. 
\]

\subsection{Caratteristiche di \texorpdfstring{$\bm{\hat{a}}$}{a} e \texorpdfstring{$\bm{\hat{a}^\dagger}$}{a dagger}}
La riformulazione a cui siamo arrivati dell'hamiltoniana è tanto più utile quanto più capiamo con cosa, \emph{effettivamente}, abbiamo a che fare. La prima cosa che possiamo fare è scriverli in forma più compatta, introducendo il parametro $\lambda$ detto \textbf{lunghezza caratteristica dell'oscillatore armonico}:
\[
    \hat{a} = \frac{1}{\sqrt{2}}\left( \frac{\hat{x}}{\lambda} + \frac{i\lambda}{\hbar}\hat{p} \right), \quad \hat{a}^\dagger = \frac{1}{\sqrt{2}}\left( \frac{\hat{x}}{\lambda} - \frac{i\lambda}{\hbar}\hat{p} \right), \quad \lambda = \sqrt{\frac{\hbar}{m\omega}} \quad.
\]

\subsubsection{Il commutatore}
La prima caratteristica che possiamo verificare dei due operatori è se \emph{commutano}:
\[
\begin{aligned}
    [\hat{a}, \hat{a}^\dagger] = \frac{1}{2}\left[\left( \frac{\hat{x}}{\lambda} + \frac{i\lambda}{\hbar}\hat{p} \right), \left( \frac{\hat{x}}{\lambda} - \frac{i\lambda}{\hbar}\hat{p} \right) \right] &= \frac{1}{2}\left\{ \left[ \frac{\hat{x}}{\lambda},  -\frac{i\lambda}{\hbar}\hat{p}\right] + \left[\frac{i\lambda}{\hbar}\hat{p}, \frac{\hat{x}}{\lambda}\right]\right\} \\
    &= \left[ \frac{\hat{x}}{\lambda},  -\frac{i\lambda}{\hbar}\hat{p}\right] = -\frac{i}{\hbar}[\hat{x},\hat{p}] = \mathds{1} \quad.
\end{aligned}
\]
Il risultato che abbiamo appena ottenuto è incredibilmente importante ma il motivo sarà chiaro solo tra poco. Limitiamoci, per adesso, a visualizzare questo risultato \emph{puramente quantistico}: gli operatori di annichilazione e creazione\footnote{La risposta alla domanda: creazione e annichilazione \emph{di cosa?} sarà presto rivelata.} \textbf{non commutano}.

\subsubsection{L'operatore numero}
Dal momento che nell'hamiltoniana compare esplicitamente il termine $\hat{a}^\dagger\hat{a}$, definiamo questo prodotto come \textbf{operatore numero} $\hat{n}$. L'hamiltoniana si scrive:
\begin{equation}
    \h{H} = \omega \hbar \left( \hat{n} + \frac{1}{2} \right)
    \label{eq:ham_osc}, \quad \hat{n} = \hat{a}^\dagger\hat{a}
\end{equation}
con $\hat{n}$ che ci darà informazioni del tipo: \emph{l'oscillatore si trova nel livello energetico n}. In questo senso, possiamo iniziare a intravedere il significato di $\hat{a}$ e $\hat{a}^\dagger$: se l'interesse è quello di scoprire il livello energetico del sistema, l'operatore di annichilazione \emph{distrugge} uno stato energetico, abbassando l'energia totale; l'operatore di creazione lavora il contrario. L'evidenza appena dimostrata per cui $[\hat{a}, \hat{a}^\dagger] \neq 0$, poi, ci testimonia il fatto che \emph{abbassare e alzare il livello energetico di un sistema sono operazioni non commutative}\footnote{In soldoni, creare $\to$ distruggere $\neq$ distruggere $\to$ creare.} Ci premureremo di verificare questa ipotesi nel prosieguo della trattazione.

Si noti, in ultimo, che risolvere l'equazione agli autovalori di $\h{H}$ implica risolvere la medesima per $\hat{n}$! Ci concentreremo infatti sulla particolare equazione:
\[
    \hat{n}\psi_n = {n} \psi_n
\]
per poi estendere la soluzione all'hamiltoniana tramite la \eqref{eq:ham_osc}.

\begin{approfondimento}[title= {L'algebra di Weyl-Heisenberg}]
    Da qualche tempo a questa parte abbiamo lavorato con operatori, operatori, operatori... ma algebricamente cosa sono? 

    Esistono alcuni \emph{set} di operatori che sono in qualche modo \emph{chiusi} rispetto all'operazione di commutazione. Gli insiemi di questo tipo definiscono una cosiddetta \textbf{algebra di Weyl-Heisenberg} (un particolare tipo di \emph{Algebra di Lie}), con gli elementi dell'insieme detti \emph{generatori dell'algebra}. Nel caso in esame, abbiamo che 
    \[
        \mathcal{P} = \left\{ \mathds{1}, \hat{a}, \hat{a}^\dagger, \hat{n} \right\}
    \]
    definisce un'algebra di Weyl-Heisenberg. Per vederlo, verifichiamo innanzitutto i commutatori degli elementi di $\mathcal{P}$. Poiché l'operatore identità $\mathds{1}$ commuta banalmente con qualsiasi altro operatore ($[\mathds{1}, \hat{A}] = 0$), e dato che ogni operatore commuta con se stesso, gli unici commutatori non nulli che definiscono la struttura dell'algebra sono tre:
    
    \begin{enumerate}
        \item \textbf{Il commutatore fondamentale tra gli operatori di scala:}
        \begin{equation}
            [\hat{a}, \hat{a}^\dagger] = \mathds{1}
        \end{equation}
        \item \textbf{L'azione dell'operatore numero sull'operatore di distruzione:}
        \begin{equation}
            [\hat{n}, \hat{a}] = [\hat{a}^\dagger\hat{a}, \hat{a}] = \hat{a}^\dagger \underbrace{[\hat{a}, \hat{a}]}_{0} + [\hat{a}^\dagger, \hat{a}]\hat{a} = -\mathds{1}\hat{a} = -\hat{a}
        \end{equation}
        \item \textbf{L'azione dell'operatore numero sull'operatore di creazione:}
        \begin{equation}
            [\hat{n}, \hat{a}^\dagger] = [\hat{a}^\dagger\hat{a}, \hat{a}^\dagger] = \hat{a}^\dagger [\hat{a}, \hat{a}^\dagger] + \underbrace{[\hat{a}^\dagger, \hat{a}^\dagger]}_{0}\hat{a} = \hat{a}^\dagger\mathds{1} = \hat{a}^\dagger
        \end{equation}
    \end{enumerate}
    Come desideravamo, risultato di ciascuna di queste commutazioni non fa uscire il sistema dal set iniziale: otteniamo sempre un'espressione che è una \textbf{combinazione lineare} (c.e.) degli elementi di $\mathcal{P}$. In termini formali, una generica c.e. così definita:
    \[  
        \h{A} = c\hat{a} + c^* \hat{a}^\dagger + r \hat{n} + s\mathds{1}, \quad c \in \mathbb{C}, \quad r,s \in \mathbb{R}
    \]
    mantiene il commutatore di $\h{A}$ con un qualsiasi altro elemento di $\mathcal{P}$ nell'insieme. I termini $c, r, s$ prendono il nome di \textbf{costanti di struttura} dell'algebra e ne definiscono interamente la geometria intrinseca. 

    Ma perché è utile notare questa particolare impalcatura nascosta dietro gli operatori che utilizziamo? In fisica teorica, identificare questa struttura algebrica è una pietra miliare: significa che tutte le proprietà dinamiche dello spettro dell'oscillatore armonico quantistico, comprese le transizioni energetiche e la discretizzazione dei livelli, possono essere dedotte unicamente manipolando queste relazioni di commutazione, senza mai dover risolvere esplicitamente l'equazione differenziale di Schrödinger nello spazio delle coordinate... \emph{esattamente la semplicazione che cercavamo!}.
\end{approfondimento}

\subsection{Applicazioni di \texorpdfstring{$\bm{\hat{a}}$}{a} e \texorpdfstring{$\bm{\hat{a^\dagger}}$}{a dagger}}
Una volta finito tutto questo \emph{pippone}, cerchiamo di comprendere cosa effettivamente fanno gli operatori di creazione e annichilazione. Dal momento che quanto ci interessa è diagonalizzare $\hat{n}$, vediamo cosa succede se applichiamo l'operatore numero prima a un generico stato $\hat{a}^\dagger \psi_n$, poi a $\hat{a} \psi_n$. \emph{Si dia il via alle danze!}

\subsubsection{Generazione verso l'alto}
Iniziamo i calcoli, assumendo $\psi_n$ autofunzione di $\hat{n}$:
\[
    \begin{aligned}
        \hat{n}(\hat{a}^\dagger \psi_n) = (\hat{n}\hat{a}^\dagger) \psi_n &= (\hat{n}\hat{a}^\dagger - \hat{a}^\dagger\hat{n} + \hat{a}^\dagger\hat{n} ) \psi_n \\
        &= ([\hat{n},\hat{a}^\dagger]+\hat{a}^\dagger\hat{n})\psi_n = (\hat{a}^\dagger+\hat{a}^\dagger\hat{n})\psi_n \\
        &= \hat{a}^\dagger\psi_n + \hat{a}^\dagger n\psi_n = (n+1)\hat{a}^\dagger\psi_n \quad.
    \end{aligned}
\]
La funzione d'onda $\hat{a}^\dagger\psi_n$ è dunque autofunzione per $\hat{n}$ associata all'autovalore $n+1$. Facciamo adesso un piccolo passo di astrazione: se l'operatore numero "estrae" il valore $n$ associato alla sua generica autofunzione $\psi_n$, allora essendo $\hat{a}^\dagger \psi_n$ la sua autofunzione associata all'autovalore $n+1$, serve che:
\[
    \hat{a}^\dagger\psi_n = c_n \psi_{n+1} \implies \hat{n}\hat{a}^\dagger\psi_n = \hat{n}c_n \psi_{n+1} \propto n+1 \quad.
\]
Per ricavare il coefficiente $c_n$ imponiamo la normalizzazione di $\psi_{n+1}$:
\[
    \begin{aligned}
        (\psi_{n+1}, \psi_{n+1}) = \left( \frac{\hat{a}^\dagger\psi_n}{c_n}, \frac{\hat{a}^\dagger\psi_n}{c_n} \right) &= \frac{1}{|c_n|^2} (\psi_n, \hat{a}\hat{a}^\dagger\psi_n) \\
        &\overset{(1)}{=} \frac{1}{|c_n|^2} (\psi_n, (\hat{n}+1)\psi_n) \overset{(2)}{=} \frac{1}{|c_n|^2} + \frac{1}{|c_n|^2} (\psi_n, \hat{n}\psi_n)\\ &= \frac{n+1}{|c_n|^2} = 1 \quad,
    \end{aligned}
\]
dove in $(1)$ abbiamo sfruttato il commutatore $[\hat{a}, \hat{a}^\dagger] = \hat{a}\hat{a}^\dagger - \hat{a}^\dagger\hat{a} = \hat{a}\hat{a}^\dagger - \hat{n} = 1$ e in $(2)$ il fatto che la richiestache $\psi_{n+1}$ sia normalizzata si estende ricorsivamente a tutti gli altri $n$. Scegliendo, per convenienza, un fattore di fase nullo, otteniamo:
\[
    c_n = \sqrt{n+1}
\]
e l'equazione agli autovalori per l'operatore di creazione risulta:
\[
    \hat{a}^\dagger\psi_n = \sqrt{n+1} \, \psi_{n+1} \quad.
\]
La relazione appena trovata è potentissima: conoscendo il generico stato $\psi_{n}$, siamo in grado di generare a cascata tutti quelli a energia superiore. Chissà se si potrà iterare questo procedimento per partire da uno stato "base" per trovare tutti gli altri...

\subsubsection{Generazione verso il basso}
Ripetiamo la stessa derivazione, questa volta sullo stato $\hat{a}\psi_n$. Ipotizzando sempre $\psi_n$ autofunzione di $\hat{n}$ abbiamo:
\[
    \begin{aligned}
        \hat{n}(\hat{a} \psi_n) = (\hat{n}\hat{a}) \psi_n &= (\hat{n}\hat{a} - \hat{a}\hat{n} + \hat{a}\hat{n} ) \psi_n \\
        &= ([\hat{n},\hat{a}]+\hat{a}\hat{n})\psi_n = (-\hat{a}+\hat{a}\hat{n})\psi_n \\
        &= -\hat{a}\psi_n + \hat{a}n\psi_n = (n-1)\hat{a}\psi_n \quad.
    \end{aligned}
\]
Con un ragionamento simile a quello condotto in precedenza, deduciamo che deve essere:
\[
    \hat{a}\psi_n = d_n \psi_{n-1} \quad,
\]
con il coefficiente $d_n$ trovato normalizzando $\psi_{n-1}$:
\[
    \begin{aligned}
        (\psi_{n-1}, \psi_{n-1}) = \left( \frac{\hat{a}\psi_n}{d_n}, \frac{\hat{a}\psi_n}{d_n} \right) &= \frac{1}{|d_n|^2} (\psi_n, \hat{a}^\dagger\hat{a}\psi_n) \\
        &= \frac{1}{|d_n|^2} (\psi_n, \hat{n}\psi_n) = \frac{n}{|d_n|^2} = 1 \implies d_n = \sqrt{n} \quad,
    \end{aligned}
\]
assumendo, come fatto prima, un fattore di fase nullo. L'equazione agli autovalori per l'operatore di annichilazione risulta quindi:
\begin{equation}
    \hat{a}\psi_n = \sqrt{n} \, \psi_{n-1}.
    \label{eq:eigenvalue_annic}
\end{equation}
Ancora una volta, abbiamo ottenuto un operatore che genera tutti gli stati \emph{inferiori} rispetto a uno di partenza. 

Siamo finalmente in grado di \textbf{generare tutto lo spettro di autofunzioni} di $\hat{n}$ (e quindi di $\h{H}$) tramite $\hat{a}$ e $\hat{a}^\dagger$.

\section{Generazione dello spettro degli autostati}
Sfruttiamo la discussione appena condotta per costruire lo spettro dell'operatore numero. Notiamo innanzitutto che la \eqref{eq:eigenvalue_annic} ci permette di assumere che lo stato \emph{più basso} a disposizione per l'oscillatore armonico è contraddistinto da $ n = 0$, infatti
\begin{equation}
    \hat{a} \psi_0 = 0 \label{eq:ground_state}
\end{equation}
e procedendo iterativamente otteniamo che $\psi_n = 0 \quad \forall \, n \in \mathbb{Z}^-$. Forti di questa considerazione, abbiamo naturalmente che l'applicazione dell'operatore di creazione a $\psi_0$ (detto \emph{ground state}, dall'inglese \emph{stato base}) genera ricorsivamente tutti gli $\psi_n, \, n > 0$. 

\subsection{Partiamo dal fondo}
Verifichiamo cosa succede applicando ripetutamente $\hat{a}^\dagger$ a partire dallo stato fondamentale $\psi_0$, ricordando l'equazione agli autovalori $\hat{a}^\dagger\psi_n = \sqrt{n+1} \, \psi_{n+1}$:

\begin{itemize}
    \item \textbf{Livello $\bm{n = 1}$:}
    Applicando l'operatore di creazione a $\psi_0$, otteniamo direttamente il primo stato eccitato:
    \[
        \hat{a}^\dagger \psi_0 = \sqrt{1} \, \psi_1 \implies \psi_1 = \hat{a}^\dagger \psi_0 \quad.
    \]
    
    \item \textbf{Livello $\bm{n = 2}$:}
    Applicando nuovamente $\hat{a}^\dagger$ allo stato $\psi_1$ appena trovato ed esplicitando $\psi_1$ in funzione di $\psi_0$:
    \[
        \hat{a}^\dagger \psi_1 = \sqrt{2} \, \psi_2 \implies \psi_2 = \frac{1}{\sqrt{2}}\hat{a}^\dagger \psi_1 = \frac{1}{\sqrt{2}} (\hat{a}^\dagger)^2 \psi_0 \quad.
    \]
    
    \item \textbf{Livello $\bm{n = 3}$:}
    Iterando il processo sullo stato $\psi_2$:
    \[
        \hat{a}^\dagger \psi_2 = \sqrt{3} \, \psi_3 \implies \psi_3 = \frac{1}{\sqrt{3}}\hat{a}^\dagger \psi_2 = \frac{1}{\sqrt{3\cdot 2}} (\hat{a}^\dagger)^3 \psi_0 \quad.
    \]
    
    \item \textbf{Iterazione $\bm{n}$-esima:}
    A questo punto è immediato dedurre la formula generale per ricorrenza. L'azione progressiva genera a denominatore il prodotto di tutti i numeri interi da $1$ fino a $n$, ricostruendo esattamente il fattoriale:
    \[
        \hat{a}^\dagger \psi_{n-1} = \sqrt{n} \, \psi_n \implies \psi_n = \frac{1}{\sqrt{n}} \hat{a}^\dagger \psi_{n-1} = \frac{1}{\sqrt{n!}} (\hat{a}^\dagger)^n \psi_0 \quad.
    \]
\end{itemize}

Dall'ultima equazione abbiamo ricavato la fondamentale relazione ricorsiva:
\begin{equation}
    \psi_n = \frac{(\hat{a}^\dagger)^n}{\sqrt{n!}}\psi_0 \quad.
\end{equation}
In questo senso, l'algebra utilizzata sinora genera tutto lo spettro di $\hat{n}$ partendo \emph{esclusivamente} dallo stato iniziale $\psi_0$. L'unica \emph{task} che rimane da espletare per poter classificare come risolto il problema dell'oscillatore armonico è quella di trovare una formulazione funzionale esplicita per il ground state. Per riuscire in ciò, sfruttiamo la \eqref{eq:ground_state}.

\subsubsection{Il ground state $\bm{\psi_0}$}
Ricordando la definizione di operatore di distruzione, abbiamo:
\[
\begin{aligned}
    \hat{a}\psi_0 = \left(\hat{x} + i\frac{\hat{p}}{m\omega}\right)\psi_0 &= 0 \\
    \left({x} + \frac{\hbar}{m\omega} \dv{}{x}\right)\psi_0 &\overset{(1)}{=} \left({x} + \lambda^2 \dv{}{x}\right)\psi_0 = 0
\end{aligned}
\]
dove in $(1)$ abbiamo tradotto gli operatori in gioco nella loro forma funzionale e abbiamo riconosciuto la definizione di $\lambda$. Svolgendo la parentesi, otteniamo la seguente equazione differenziale:
\[
    \lambda^2\dv{\psi_0}{x} + \psi_0 x = 0 \quad.
\]
Separando le variabili otteniamo:
\[
    \begin{aligned}
        \int\frac{1}{\psi_0}\dd{\psi_0} &= -\frac{1}{\lambda^2} \int x\dd{x} \\
        \ln{\psi_0} &= -\frac{x^2}{2\lambda^2} + K \\
        \psi_0(x) &= Ce^{-\frac{x^2}{2\lambda^2}}, \quad C \in \mathbb{C}
    \end{aligned}
\]
e per determinare $C$ imponiamo la normalizzazione:
\[
\begin{aligned}
    1 = (\psi_0, \psi_0) = |C|^2 \int_\mathbb{R} e^{-\frac{x^2}{\lambda^2}} \dd{x} &= \lambda\sqrt{ \pi} \\
    \implies C &= \left(\frac{1}{\lambda\sqrt{\pi}}\right)^{\frac{1}{2}}
\end{aligned}
\]
dove, come sempre, abbiamo scelto fase nulla per $C$. Il ground state assoume quindi la forma di:
\begin{equation}
    \psi_0(x) = \left(\frac{1}{\lambda\sqrt{\pi}}\right)^{\frac{1}{2}}e^{-\frac{x^2}{2\lambda^2}} = \left(\frac{2m\omega}{h}\right)^{\frac{1}{4}} e^{-\frac{x^2}{2\lambda^2}}\quad.
\end{equation}
Si osservi che la densità di probabilità associata a $\psi_0$ ha il profilo di una gaussiana centrata nell'origine: così come nel caso della particella libera, anche per l'oscillatore armonico esiste un'interpretazione \emph{semiclassica} dei risultati quantistici e nel caso del ground state questa considerazione è ancora più evidente: $\psi_0$ è infatti uno stato \textbf{localizzato}.

\subsection{Costruzione della soluzione e Formula di Rodrigues}
Abbiamo tutti i mattoni necessari a costruire la generica autofunzione $\psi_n$. A partire dalla relazione ricorsiva e dal \emph{ground state} appena trovato, ricordando che ogni operatore $\hat{a}^\dagger$ introduce un coefficiente $1/\sqrt{2}$, possiamo scrivere:
\[
    \psi_n(x) = \frac{1}{\sqrt{n!}}(\hat{a}^\dagger)^n \psi_0(x) = \frac{1}{\sqrt{2^n n!}}\left( \frac{x}{\lambda} - \lambda \dv{}{x} \right)^n \left(\frac{1}{\lambda\sqrt{\pi}}\right)^{\frac{1}{2}}e^{-\frac{x^2}{2\lambda^2}} \quad.
\]
Per procedere in modo sistematico e svelare la natura funzionale di questi stati, operiamo il cambio di variabile adimensionale $y = x/\lambda$. L'operatore di creazione diventa proporzionale a $(y - \partial_y)$, e la funzione d'onda del ground state assume la forma $\psi_0(y) = C e^{-y^2/2}$, con $C = (\lambda\sqrt{\pi})^{-1/2}$.

Il segreto per iterare questo operatore senza perdersi nei calcoli risiede in una notevole identità operatoriale. Notiamo infatti che l'azione dell'operatore $(y - \partial_y)$ su una generica funzione $f(y)$ differenziabile può essere riscritta come:
\[
    \left( y - \dv{}{y} \right) f(y) = -e^{\frac{y^2}{2}} \dv{}{y} \left( e^{-\frac{y^2}{2}} f(y) \right) \quad.
\]
Applicando questa identità ripetutamente alla nostra $\psi_0$, l'iterazione svela immediatamente la sua intima struttura matematica:

\begin{itemize}
    \item \textbf{Livello $\bm{n = 1}$:}
    \[
        (\hat{a}^\dagger) \psi_0 = \frac{1}{\sqrt{2}}\left( y - \dv{}{y} \right) Ce^{-\frac{y^2}{2}} = \frac{C}{\sqrt{2}} \left[ -e^{\frac{y^2}{2}} \dv{}{y} \left( e^{-\frac{y^2}{2}} e^{-\frac{y^2}{2}} \right) \right] = \frac{C}{\sqrt{2}} (-1)^1 e^{\frac{y^2}{2}} \dv{}{y} \left( e^{-y^2} \right) \quad.
    \]

    \item \textbf{Livello $\bm{n = 2}$:} 
    Applicare nuovamente l'operatore significa agire sul risultato appena ottenuto, sfruttando di nuovo l'identità operatoriale:
    \[
    \begin{aligned}
        (\hat{a}^\dagger)^2 \psi_0 &= \frac{1}{\sqrt{2}}\left( y - \dv{}{y} \right) \left[ \frac{C}{\sqrt{2}} (-1) e^{\frac{y^2}{2}} \dv{}{y} \left( e^{-y^2} \right) \right] \\
        &= \frac{C}{2} \left\{ -e^{\frac{y^2}{2}} \dv{}{y} \left[ e^{-\frac{y^2}{2}} \left( -e^{\frac{y^2}{2}} \dv{}{y} e^{-y^2} \right) \right] \right\} \quad.
    \end{aligned}
    \]
    Il termine $e^{-y^2/2}$ e il termine $e^{y^2/2}$ all'interno della parentesi quadra si elidono in modo esatto, permettendo all'operatore derivata di agire direttamente sulla derivata precedente:
    \[
        (\hat{a}^\dagger)^2 \psi_0 = \frac{C}{2} (-1)^2 e^{\frac{y^2}{2}} \dv[2]{}{y} \left( e^{-y^2} \right) \quad.
    \]

    \item \textbf{Iterazione $\bm{n}$-esima:}
    A questo punto il \emph{pattern} è perfettamente cristallizzato. Ad ogni applicazione di $\hat{a}^\dagger$, l'esponenziale generato cancella in modo esatto quello residuo del passaggio precedente, incrementando l'ordine di derivazione e aggiungendo un segno meno. In generale, avremo:
    \[
        (\hat{a}^\dagger)^n \psi_0 = \frac{C}{\sqrt{2^n}} (-1)^n e^{\frac{y^2}{2}} \dv[n]{}{y} \left( e^{-y^2} \right) \quad.
    \]
\end{itemize}
Moltiplicando e dividendo l'espressione per $e^{-y^2/2}$ per isolare il termine gaussiano originale della funzione d'onda, otteniamo:
\[
    (\hat{a}^\dagger)^n \psi_0 = \frac{C}{\sqrt{2^n}} \underbrace{\left[ (-1)^n e^{y^2} \dv[n]{}{y} \left( e^{-y^2} \right) \right]}_{\coloneq \, H_n(y)} e^{-\frac{y^2}{2}} \quad.
\]
Il termine tra parentesi quadre è l'esatta definizione matematica dei \textbf{Polinomi di Hermite} $H_n(y)$, espressa nella sua forma più elegante e compatta: la \textbf{Formula di Rodrigues}. 

Inserendo la costante di normalizzazione e tornando alla variabile spaziale $x$, l'autofunzione $n$-esima dell'oscillatore armonico è rigorosamente dimostrata essere:
\begin{equation}
    \psi_n(x) = \left( \frac{1}{2^n n! \lambda \sqrt{\pi}} \right)^{\frac{1}{2}} H_n\left(\frac{x}{\lambda}\right) e^{-\frac{x^2}{2\lambda^2}} \quad.
\end{equation}
Abbiamo così dimostrato che il metodo algebrico e la complessa teoria delle equazioni differenziali restituiscono, per profonda necessità strutturale, il medesimo risultato fisico. Si noti ancora un'incredibile peculiarità: gli autostati sono \textbf{tutti puramente reali}, $\psi_n: \mathbb{R} \to \mathbb{R}$!

\subsubsection{Interpretazione fisica}
Quanto appena ricavato è, come detto, il set di autofunzioni normalizzate (tramite i coefficienti $c_n$ e $d_n$) che diagonalizzano l'operatore numero. Verifichiamo che una generica $\psi_n$ sia autofunzione anche per $\h{H}$:
\[
    \h{H} = \hbar \omega \left(\hat{n} + \frac{1}{2}\right) \psi_n = \hbar \omega \left(\hat{n}\psi_n + \frac{\psi_n}{2}\right) = \underbrace{\hbar \omega \left({n} + \frac{1}{2}\right)}_{E_n} \psi_n \quad.
\]
L'$n$-esimo autovalore dell'hamiltoniana, $E_n$, appartiene a uno spettro discreto! L'evidenza, intuibile dal fatto che lo spettro dell'operatore numero è per costruzione discreto, ci porta a comprendere che i livelli energetici concessi a un oscillatore armonico sono \textbf{discreti}. Ancora, dal momento che le $\psi_n$ appartengono allo spettro di un operatore hermitiano\footnote{Per completezza, si noti che anche $\hat{n}$, oltre ad $\h{H}$, è hermitiano. Infatti: $\hat{n}^\dagger = (\hat{a}^\dagger\hat{a})^\dagger = \hat{a}^\dagger (\hat{a}^\dagger)^\dagger = \hat{a}^\dagger \hat{a} = \hat{n}$.}, formano una base per $\mathcal{L}^2(\Omega)$ e garantiscono autovalori \textit{reali}. 

Possiamo in ultimo verificare che, presi due generici autostati dell'hamiltoniana $\psi_m$ e $\psi_n$ con $m \neq n$, abbiamo:
\[
    (\psi_n,\psi_m) = \int_{\mathbb{R}} \dd{x} \psi_n(x)^*\psi_m(x) \overset{(1)}{=} \int_{\mathbb{R}} \dd{x} \psi_m(x)\psi_n(x) = (\psi_n,\psi_m) \quad,
\]
dove in $(1)$ abbiamo sfruttato il fatto che gli autostati sono reali. Allora:
\[
    \begin{aligned}
        (\psi_m,\psi_n) &= \int_{\mathbb{R}} \dd{x}\psi_m(x)\psi_n(x) \\
        &= \int_{\mathbb{R}} \dd{x} \, \left( \frac{1}{2^m m! \lambda \sqrt{\pi}} \right)^{\frac{1}{2}} H_m\left(\frac{x}{\lambda}\right) e^{-\frac{x^2}{2\lambda^2}} \left( \frac{1}{2^n n! \lambda \sqrt{\pi}} \right)^{\frac{1}{2}} H_n\left(\frac{x}{\lambda}\right) e^{-\frac{x^2}{2\lambda^2}} \\
        &\overset{(2)}{=} \frac{1}{\sqrt{2^n2^m n! m! \pi}}\int_{\mathbb{R}} \dd{y} \, e^{-y^2}H_m(y)H_n(y) \quad,
    \end{aligned}
\]
dove in $(2)$ abbiamo eseguito il cambio di variabile $x \to y = x/\lambda$. Per risolvere esplicitamente l'integrale appena ottenuto, definiamolo come $I_{m,n}$ e assumiamo, senza perdita di generalità, che $m \leq n$\footnote{Nel caso opposto è sufficiente invertire gli autostati nel prodotto scalare essendo, in questo caso specifico, \textit{simmetrico}.}. Sfruttiamo la formula di Rodrigues per riscrivere il polinomio di grado maggiore $H_n(y)$:
\[
    H_n(y) = (-1)^n e^{y^2} \dv[n]{}{y} \left( e^{-y^2} \right) \quad.
\]
Sostituendola nell'integrale, il termine $e^{-y^2}$ e il fattore $e^{y^2}$ della formula di Rodrigues si elidono in modo esatto, lasciando:
\[
    I_{m,n} = \int_{-\infty}^{+\infty} \dd{y} \, e^{-y^2} H_m(y) \left[ (-1)^n e^{y^2} \dv[n]{}{y} \left( e^{-y^2} \right) \right] = (-1)^n \int_{-\infty}^{+\infty} \dd{y} \, H_m(y) \dv[n]{}{y} \left( e^{-y^2} \right) \quad.
\]
A questo punto, procediamo integrando per parti la funzione. Eseguendo la prima iterazione otteniamo:
\[
    I_{m,n} = (-1)^n \underbrace{\left[ H_m(y) \dv[n-1]{}{y}\left(e^{-y^2}\right) \right]_{-\infty}^{+\infty}}_{0} - (-1)^n \int_{-\infty}^{+\infty} \dd{y} \, \dv{H_m(y)}{y} \dv[n-1]{}{y} \left( e^{-y^2} \right) \quad.
\]
Il termine di bordo si annulla identicamente ai limiti di integrazione: l'esponenziale decrescente $e^{-y^2}$ (e tutte le sue derivate, che contengono sempre il fattore gaussiano) domina asintoticamente su qualsiasi espressione polinomiale per $y \to \pm\infty$. 

Iterando l'integrazione per parti per un totale di $n$ volte, accade una reazione a catena: a ogni passaggio il termine di bordo si annulla, si raccoglie un segno meno (che alla fine restituisce un fattore $(-1)^n$ andandosi a cancellare con quello iniziale), e l'operatore derivata viene "trasferito" progressivamente dal termine gaussiano al polinomio $H_m$:
\[
    I_{m,n} = \int_{-\infty}^{+\infty} \dd{y} \, \left( \dv[n]{H_m(y)}{y} \right) e^{-y^2} \quad.
\]
Questa espressione è la chiave di volta analitica del problema. Essendo $H_m(y)$ un polinomio di grado $m$, si delineano due scenari distinti:
\begin{itemize}
    \item \textbf{Se $m < n$:} l'operatore sta derivando un polinomio un numero di volte superiore al suo grado. La derivata $n$-esima è di conseguenza identicamente nulla su tutto il dominio, il che implica inequivocabilmente $I_{m,n} = 0$.
    \item \textbf{Se $m = n$:} la derivata $n$-esima di un polinomio di grado $n$ si riduce a una costante determinata unicamente dal coefficiente del suo termine di grado massimo. Poiché il termine direttore del polinomio di Hermite è noto essere $(2y)^n$, derivandolo $n$ volte si ottiene esattamente $n! 2^n$. L'integrale collassa così nel calcolo del celebre integrale gaussiano:
    \[
        I_{n,n} = 2^n n! \int_{-\infty}^{+\infty} \dd{y} \, e^{-y^2} = 2^n n! \sqrt{\pi} \quad.
    \]
\end{itemize}
Possiamo compattare in modo elegante entrambi i risultati sfruttando la definizione di delta di Kronecker:
\[
    I_{m,n} = \int_{\mathbb{R}} \dd{y} \, e^{-y^2}H_m(y)H_n(y) = 2^n n! \sqrt{\pi} \, \delta_{mn} \quad.
\]
Sostituendo finalmente questo risultato calcolato nel nostro prodotto scalare originale in $\mathcal{L}^2$, il cerchio si chiude in modo perfetto:
\[
\begin{aligned}
    (\psi_m,\psi_n) &= \frac{1}{\sqrt{2^n2^m n! m! \pi}} \left( 2^n n! \sqrt{\pi} \, \delta_{mn} \right) \\
    &= \frac{2^n n! \sqrt{\pi}}{\sqrt{(2^n)^2 (n!)^2 \pi}} \, \delta_{mn} = \delta_{mn} \quad.
\end{aligned}
\]
L'ortonormalità della base dell'oscillatore armonico è così rigorosamente provata anche in ambito funzionale. Otteniamo così gratuitamente la relazione di completezza nel caso discreto:
\[
    \delta(y-x) = \sum_{i=0}^\infty \psi_n(y)^* \psi_n(x) = \sum_{i=0}^\infty \psi_n(y) \psi_n(x) \quad.
\]

\section{Valori di aspettazione di posizione e momento}
Una volta "risolto" il problema dell'oscillatore dal punto di vista matematico, ci dedichiamo a caratterizzarlo dal punto di vista fisico. Classicamente la "posizione media" e la "quantità di moto media" assunte da un occillatore armonico (una molla, ad esempio) di periodo $T$ e descritto dalla seguente evoluzione temporale:
\[
    x(t) = A\sin(\omega t + \theta) \implies p(t) = m\dot{x} = m\omega A \cos(\omega t + \theta) \quad,
\] 
sono date da:
\[
    \begin{aligned}
        &\avg{x} \to \frac{1}{T}\int_0^T x(t)\dd{t} = 0\\
        &\avg{p} \to \frac{1}{T}\int_0^T p(t)\dd{t} = 0 \quad,
    \end{aligned}
\]
che sono risultati assolutamente sensati notando che il moto è centrale avviene attorno alla posizione $x = 0$ con velocità media nulla. Curiosamente, nel mondo quantistico, il risultato è il medesimo.

\subsection{Calcolo dei valori di aspettazione quantistici}
Applicando la definizione abbiamo per la posizione:
\[
    \begin{aligned}
        \avg{\hat{x}} \to (\psi_n, \hat{x}\psi_n) &= \left( \psi_n, \sqrt{\frac{\hbar}{2m\omega}} (\hat{a}+\hat{a}^\dagger)\psi_n \right) = \sqrt{\frac{\hbar}{2m\omega}} [(\psi_n, \hat{a}\psi_n) + (\psi_n, \hat{a}^\dagger\psi_n)]\\
        &= \sqrt{\frac{\hbar}{2m\omega}} [\sqrt{n}(\psi_n, \psi_{n-1}) + \sqrt{n+1}(\psi_n, \psi_{n+1})] \overset{(1)}{=} 0
    \end{aligned}
\]
e per la quantità di moto:
\[
    \begin{aligned}
        \avg{\hat{p}} \to (\psi_n, \hat{p}\psi_n) &= \left( \psi_n, -i\sqrt{\frac{m\hbar\omega}{2}} (\hat{a}-\hat{a}^\dagger)\psi_n \right) = -i\sqrt{\frac{m\hbar\omega}{2}} [(\psi_n, \hat{a}\psi_n) - (\psi_n, \hat{a}^\dagger\psi_n)]\\
        &= -i\sqrt{\frac{m\hbar\omega}{2}} [\sqrt{n}(\psi_n, \psi_{n-1}) - \sqrt{n+1}(\psi_n, \psi_{n+1})] \overset{(1)}{=} 0 \quad,
    \end{aligned}
\]
dove in $(1)$ abbiamo sfruttato il fatto che le autofunzioni di un operatore hermitiano sono ortonomali e quindi $(\psi_n, \psi_{n+1}) = \delta_{n,n+1} = 0$. Anche quantisticamente parlando, la posizione media e la quantità di moto media sono nulle. \emph{Eppure qualcosa dovrà pur cambiare... no?}

\subsection{Altri dati interessanti sulle misure di posizione e quantità di moto}
In generale, ogni qual volta che si conduce una misura in laboratorio, è fondamntale studiare la \textit{deviazione (standard)} $\sigma$ dei risultati ottenuti rispetto alla media. Classicamente è definita come il quadrato della \textit{varianza}, oppure:
\[
    \sigma = \sqrt{\frac{1}{N}\sum_{i=1}^N (e_i - \mu)^2}
\]
dove abbiamo chiamato $e_i$ l'$i$-esimo risultato di misura, $\mu$ la media 
\[
    \mu = \frac{1}{N}\sum_{i=1}^N e_i
\]
e $N$ il numero di processi di misurazione eseguiti. Se espandiamo il quadrato, otteniamo:
\[
    \begin{alignedat}{2}
        \sigma = \sqrt{\frac{1}{N}\sum_{i=1}^N (e_i - \mu)^2} = \sqrt{\frac{1}{N}\sum_{i=1}^N (e_i^2 + \mu^2 - 2e_i \mu)} &= \sqrt{\frac{1}{N}\sum_{i=1}^N e_i^2 + \frac{1}{N}\sum_{i=1}^N \mu^2 - \frac{1}{N}\sum_{i=1}^N 2e_i \mu} \\
        &= \sqrt{\avg{e^2} + \mu^2 - \frac{2\mu}{N}\sum_{i=1}^N e_i} \\
        &= \sqrt{\avg{e^2} + \mu^2 - 2\mu^2} = \sqrt{\avg{e^2} - \mu^2} \\
        & = \sqrt{\avg{e^2} - \avg{e}^2}\quad.
    \end{alignedat}
\]
La deviazione su una misura è quindi data dalla radice della media dei quadrati, indicata con $\avg{e^2}$, meno il quadrato della media $\mu = \avg{e}$. 

La traduzione quntistica è abbastanza semplice: preso un operatore generico $\h{A}$ sulla quale grandezza associata si conduce una misura, si ha che l'\textbf{operatore di deviazione}, o errore, $\hat{\sigma}$ risulta\footnote{Si noti come in questo caso l'operatore coincida con un numero reale e, pertanto, si confonderà la notazione $\hat{\sigma} \equiv \sigma$.}:
\[
    \hat{\sigma} = \sqrt{\avg{\h{A}^2} - \avg{\h{A}}^2} \quad.
\] 
Applicando questa nozione al caso della posizione e della quantità di moto nell'oscillatore armonico otteniamo:
\[
    \begin{dcases}
        \hat{\sigma}_{\hat{x}} = {\sigma}_{\hat{x}} = \sqrt{\avg{\hat{x}^2} - \avg{\hat{x}}^2} = \sqrt{\avg{\hat{x}^2}} \\
        \hat{\sigma}_{\hat{p}} = {\sigma}_{\hat{p}} = \sqrt{\avg{\hat{p}^2} - \avg{\hat{p}}^2} = \sqrt{\avg{\hat{p}^2}} \quad,
    \end{dcases}
\]
essendo nulli i loro valori di aspettazione. Continuiamo nella computazione: iniziamo con la posizione.
\[
    \begin{aligned}
        \sigma_{\hat{x}} = \sqrt{\avg{\hat{x}^2}} = \left[(\psi_n,\hat{x}^2\psi_n)\right]^{\frac{1}{2}} &= \sqrt{\frac{\hbar}{2m\omega}}\left[(\psi_n,(\hat{a}+ \hat{a}^\dagger)^2\psi_n)\right]^{\frac{1}{2}} = \frac{\hbar}{2m\omega}\left[(\psi_n,(\hat{a}^2+ (\hat{a}^\dagger)^2 + \hat{a}\hat{a}^\dagger + \hat{a}^\dagger\hat{a})\psi_n)\right]^{\frac{1}{2}} \\
        &\overset{(1)}{=}\sqrt{\frac{\hbar}{2m\omega}} \left[(\psi_n,(\hat{a}^2 + (\hat{a}^\dagger)^2 + 2\hat{n} + 1)\psi_n)\right]^{\frac{1}{2}} \\ &=\sqrt{\frac{\hbar}{2m\omega}} \left[(\psi_n,\hat{a}^2\psi_n) + (\psi_n,(\hat{a}^\dagger)^2\psi_n) + (\psi_n,2\hat{n}\psi_n) + 1\right]^{\frac{1}{2}} \\
        &=\sqrt{\frac{\hbar}{2m\omega}} \left[\sqrt{(n+1)(n+2)}(\psi_n,\psi_{n+2}) + \sqrt{n(n-1)}(\psi_n,\psi_{n-2}) + 2n + 1\right]^{\frac{1}{2}} \\
        &= \sqrt{\frac{\hbar}{2m\omega}(2n+1)} = \lambda\sqrt{\left( n+\frac{1}{2} \right)} \quad,
    \end{aligned}
\]
dove in $(1)$ abbiamo sfruttato il commutatore $[\hat{a},\hat{a}^\dagger] = 1$. Procedendo similmente otteniamo la deviazione per la quantità di moto:
\[
    \sigma_{\hat{p}} = \frac{\hbar}{\lambda}\sqrt{\left( n+\frac{1}{2} \right)} \quad.
\]

\subsubsection{Uno sguardo al Principio di indeterminazione}
Anticipiamo ora un risultato fondamentale che dimostraremo entro breve: segue infatti una conseguenza dell'arcinoto \textbf{Prinipio di indeterminazione di Heisenberg}. Moltiplicando le due varianze otteniemo:
\[
    \sigma_{\hat{x}}\sigma_{\hat{p}} = \hbar \left( n+\frac{1}{2} \right) \quad,
\]
il cui minimo si raggiunge per $n = 0$ (in occorrenza di $\psi_0$, il ground state):
\[
    \min_{n\in\mathbb{N}} \left( \sigma_{\hat{x}}\sigma_{\hat{p}} \right) = \frac{\hbar}{2} \quad.
\]
L'interpretazione è tanto preziosa quanto lampante: se, come nel caso classico, la "posizione media" e la "quantità di moto media" sono quantisticamente nulle, la \emph{conoscenza} riguardo la precisione con cui si sono condotte le due misure segue tutt'altra logica. Infatti, tanto è minore l'incertezza su una quanto è maggiore l'incertezza sull'altra. Formalizzeremo meglio questa intuzione entro breve, per ora cerchiamo di capire come questa informazione influenzi concretamente il problema dell'oscillatore armonico.

\section{La comparsa della delocalizzazione}
Ricordate il caso della particella libera? L'analisi quantistica del problema ci aveva condotto a scoprire che il \emph{determinismo} della fisica classica si perde con il passare del tempo per lasciare il passo ad una delocalizzazione puramente quantistica. Ora, a seguito della discussione riguardo l'incertezza sulla misura di posizione e momento, abbiamo tutti gli strumenti necessari per rispondere alla domanda: -\emph{Esiste indeterminazione nel caso dell'oscillatore armonico?}-.

Una particolare proprietà dei polinomi di Hermite è tale da garantire che gli zeri del generico $H_n$ sono confinati in un intervallo simmetrico rispetto all'origine di ampiezza $r$, ovvero $[-r, r]$... e questo a cosa serve? Beh, innanzitutto la generica autofunzione di $\h{H}$ è proporzionale ad $H_n$. Supponiamo che gli zeri più distanti dall'origine di $H_n$, e quindi di $\psi_n$, si trovino in corrispondenza di $x = \pm r$: per valori maggiori di ascissa, $|x| \gg r$, l'autofunzione è ben approssimata da:
\[
    \psi_n(x) \approx x^n e^{-\frac{x^2}{2\lambda^2}}
\] 
dal momento che in $H_n$ domina il monomio di grado $n$. È evidente a questo punto come lo smorzamento esponenziale domini e che a breve distanza dagli "ultimi" zeri $\psi_n$ si schiacci rapidamente a zero. Otteniamo così che la densità di probabilità associata a $\psi_n$ è non nulla solo nella regione in cui $H_n$ ammette zeri:
\[
    \rho_n(x) = |\psi_n(x)|^2 \neq 0 \implies x \in [-r,r] \quad.
\]
Per ricavare il valore di $r$ procediamo con un argomentazione tanto euristica quanto \emph{semiclassica}. Classicamente, la posizione della \emph{massa oscillante} collegata all'oscillatore è confinata a ritrovarsi entro il doppio dell'ampiezza $A$ della sinusoide associata al moto. In questo senso la densità "classica" di probabilità è non nulla entro l'intervallo $[-A, A]$. Per costruire il nesso, scriviamo l'energia meccanica (classica) del sistema, assumendo che il moto sia descritto da una forma d'onda sinusoidale:
\[
    E = \frac{p^2}{2m} + \frac{1}{2}m\omega^2 x^2 = \frac{(m\overbrace{A\omega\cos(\omega t + \theta)}^{\dot{x}})^2}{2m} + \frac{1}{2}m\omega^2 (A\sin(\omega t + \theta))^2 = \frac{1}{2}m A^2 \omega^2 \quad.
\]
Gli autovalori dell'hamiltoniana sono, invece:
\[
    E_n = \hbar \omega \left(n+\frac{1}{2}\right)
\]
che nel limite semiclassico di "grandi" $n$ diventano $E_n \approx n\hbar\omega$. Comparando le due formule\footnote{Si presti attenzione che, benchè questa sia un'argomentazione euristica e quindi non teoricamente formale, l'uguaglianza è lecita dal momento che si sta passando al limite semiclassico di $n \to \infty$.}, otteniamo:
\[
    E = \frac{1}{2}m A^2 \omega^2 \approx E_n = n\hbar\omega \implies A_{(n)} \approx \sqrt{\frac{\hbar}{m\omega}}\sqrt{2n}  = \lambda \sqrt{2n} \quad. 
\]
In questo senso, la regione in cui possiamo trovare la massa oscillante corrisponde all'intervallo 
\[
    \mathcal{I} = \left[ -\lambda \sqrt{2n}, \lambda \sqrt{2n} \right] \quad.
\]
Abbiamo quindi mappato l'intervallo entro cui sono compresi gli zeri del generico polinomio di Hermite in questo nuovo insieme, $[-r,r] \to \mathcal{I}$. L'incertezza sulla posizione calcolata in precedenza declinata nel limite semiclassico risulta:
\[
    \sigma_{\hat{x}} = \lambda \sqrt{\left(n+\frac{1}{2}\right)} \approx \lambda \sqrt{n}
\] 
ma cosa rappresenta effettivamente questa deviazione? Ricordando che l'incertezza di misura si puù esprimere come $\mu \pm \sigma$, abbiamo banalmente che nel caso della posizione vale:
\[
    x = \avg{\hat{x}} \pm \underbrace{\sigma_{\hat{x}}}_{\Delta x}
\]
e in questo contesto $\sigma$ assume il significato di "variazione" rispetto alla media\footnote{Ovvero \emph{incertezza!}, $\sigma_{\hat{x}} \to \Delta x$.}. Il rapporto fra incertezza e ampiezza dell'intervallo dove posso trovare il corpo oscillante è data da:
\[
    \frac{\Delta x}{2A} \approx \frac{\lambda \sqrt{n}}{2\lambda \sqrt{2n}} = \frac{1}{2\sqrt{2}} \approx 35 \% \quad.
\]
I due termini sono dello stesso ordine di grandezza e l'incertezza sulla posizione corrisponde a più di $1/3$ dell'intervallo in cui \emph{è possibile} trovare il corpuscolo. Ecco la nostra amata delocalizzazione quantistica!

\begin{approfondimento}[title={Densità classica di probabilità}]
    \noindent
    Nel fare conti di meccanica quantistica, ci serviamo molto sovente dell'interpretazione di Born che identifica $\rho \to |\psi|^2$. Ma qual è la controparte classica? Facciamo un po' di \emph{reverse enegeneering}. Partiamo dal caso dell'oscillatore armonico, in cui abbiamo:
    \[
        x(t) = A\sin(\omega t + \theta) \quad.
    \]
    Differenziando ad ambro i membri otteniamo:
    \[
    \begin{aligned}
        \dd{x} = A\omega \cos(\omega t + \theta) \dd{t} = A\omega \sqrt{1- \sin^2(\omega t + \theta)} \dd{t} &= A\omega \sqrt{1- \frac{x^2}{A^2}} \dd{t} \\
        &= \omega\sqrt{A^2-x^2} \dd{t} \quad.
    \end{aligned}
    \]
    Essendo $\omega$ la pulsazione caratteristica del sistema, possiamo ancora scrivere:
    \[
        \dd{x} = \frac{2\pi}{T} \sqrt{A^2-x^2} \dd{t} = \pi \sqrt{A^2-x^2} \frac{\dd{t}}{T/2} \quad.
    \]
    Cosa rappresenta il fattore $\dd{t}/(T/2)$? L'infinitesimo di tempo corrisponde a quanto la particella impiega per passare dal punto $x$ a $x + \dd{x}$ mentre $T/2$ è il lasso di tempo necessario a fare mezza oscillazione, ovvero a muoversi da $-A$ sino ad $A$ (e viceversa). 

    Possiamo così identificare 
    \[
        \frac{\dd{t}}{T/2} = \dd{\mathcal{P}} \quad,
    \]
    con $\mathcal{P}$ la probabilità classica. Abbiamo dunque:
    \[
        \dd{x} = \pi \sqrt{A^2-x^2} \dd{\mathcal{P}}
    \]
    da cui possiamo definire la densità classica di probabilità $\rho_{{cl}}$ come:
    \[
        \rho_{cl}(x) = \dv{\mathcal{P}}{x} = \frac{1}{\pi \sqrt{A^2-x^2}} \quad.
    \]
    È interessante notare come $\rho : \text{dom}(\rho) \to \mathbb{R}^+$, con $\text{dom}(\rho) \equiv [-A, A]$, ovvero l'intervallo in cui è possibile trovare la massa oscillante (coerentemente con quanto si poteva intuire). Si presti attenzione al fatto che la densità diverge per $x \to \pm A$, pur mantenendo un integrale convergente rispetto alle singolarità. La sua primitiva, a meno di una costante, corrisponde alla probabilità di avere la particella in un intervallo $[x_1, x_2] \in [-A, A]$. 
    \begin{center}
\begin{tikzpicture}[>=stealth, scale=1.3]
    % Definizione dei colori (in tinta con il grafico precedente)
    \colorlet{myblue}{blue!70!black}
    \colorlet{myred}{red!70!black}

    % Assi cartesiani
    \draw[->, thick] (-3.5, 0) -- (3.5, 0) node[below right] {$x$};
    \draw[->, thick] (0, -0.2) -- (0, 3.5) node[above left] {$\rho_{cl}(x)$};

    % Punti di inversione classici (scelgo A = 2.5 per il disegno)
    \draw[dashed, gray, thick] (2.5, 0) -- (2.5, 3.2);
    \draw[dashed, gray, thick] (-2.5, 0) -- (-2.5, 3.2);
    \node[below, font=\small] at (2.5, 0) {$A$};
    \node[below, font=\small] at (-2.5, 0) {$-A$};
    \node[below left, font=\small] at (0, 0) {$0$};

    % DENSITÀ DI PROBABILITÀ CLASSICA: rho = C / sqrt(A^2 - x^2)
    % Il dominio si ferma a 2.45 per evitare la divisione per zero in x=2.5
    \draw[domain=-2.45:2.45, samples=200, smooth, myblue, thick] 
        plot (\x, {1.2 / sqrt(6.25 - \x*\x)});

\end{tikzpicture}
\end{center}
\end{approfondimento}

\subsection{Densità di probabilità quantistica}
Cerchiamo di dare una giustificazione formale alla nostra precedente considerazione euristica per la quale abbiamo assunto $|\psi_n(x)|^2 \approx 0$ per $|x| > A$. Dalla forma della generica autofunzione trovata in precedenza, esplicitando il modulo quadro, abbiamo:
\[
    \rho_n(x) = |\psi_n(x)|^2 = \frac{1}{2^n n! \lambda \sqrt{\pi}} H_n^2\left(\frac{x}{\lambda}\right) e^{-\frac{x^2}{\lambda^2}} \quad.
\]
La vera magia racchiusa da questa relazione emerge, come anticipato, quando studiamo il sistema per grandi numeri quantici ($n \to \infty$), ovvero quando l'energia dell'oscillatore diventa macroscopica. Dal momento che lo stato quantistico presenta $n$ nodi, all'aumentare di $n$ la funzione d'onda oscilla in modo sempre più frenetico. Tracciamo un grafico qualitativo (con $n$ arbitrariamente grande) che sovrapponga la densità quantistica $\rho_n(x)$ a quella classica $\rho_{cl}(x) = 1/(\pi\sqrt{A_n^2-x^2})$:
\begin{center}
\begin{tikzpicture}[>=stealth, scale=1.5]
    \colorlet{qcolor}{blue!70!white}
    \colorlet{ccolor}{red!80!black}
    
    % Assi
    \draw[->, thick] (-3.5, 0) -- (3.5, 0) node[below right] {$x$};
    \draw[->, thick] (0, -0.2) -- (0, 3.2) node[above left] {$\rho(x)$};
    
    % Punti di inversione classici (A = 2.8)
    \draw[dashed, gray, thick] (2.8, 0) -- (2.8, 2.5);
    \draw[dashed, gray, thick] (-2.8, 0) -- (-2.8, 2.5);
    \node[below, font=\small] at (2.8, 0) {$A_n$};
    \node[below, font=\small] at (-2.8, 0) {$-A_n$};
    \node[below left, font=\small] at (0, 0) {$0$};

    % Funzione Quantistica (Inviluppo doppio della classica * oscillazione)
    % Usiamo 1500*\x in gradi: crea un'oscillazione fitta ma sotto i limiti di memoria di LaTeX.
    % Il dominio si ferma a 2.7 per evitare la divergenza esatta in 2.8.
    \draw[domain=-2.7:2.7, samples=250, smooth, qcolor, thick] 
        plot (\x, { (1.8 / sqrt(7.84 - \x*\x)) * (cos(1500 * \x))^2 });

    % Sconfinamento quantistico (le code esponenziali raccordate)
    % A x=2.7 l'inviluppo vale circa 2.42, raccordiamo l'esponenziale da lì.
    \draw[domain=2.7:3.3, samples=30, smooth, qcolor, thick] 
        plot (\x, { 2.42 * exp(-15*(\x-2.7)) });
    \draw[domain=-3.3:-2.7, samples=30, smooth, qcolor, thick] 
        plot (\x, { 2.42 * exp(15*(\x+2.7)) });

    % Funzione Classica (Media)
    % Esattamente la metà dell'inviluppo quantistico
    \draw[domain=-2.7:2.7, samples=150, smooth, ccolor, thick, dashed] 
        plot (\x, { 0.9 / sqrt(7.84 - \x*\x) });

    % Legenda
    \node[ccolor, right, font=\footnotesize] at (-0.5, 2.7) {Classica ($\rho_{cl}$)};
    \node[qcolor, right, font=\footnotesize] at (1.5, 1.8) {Quantistica ($\rho_n$)};
\end{tikzpicture}
\end{center}
Osservando il grafico, possiamo trarre considerazioni fisiche formidabili:
\begin{itemize}
    \item \textbf{Compensazione integrale:} La curva quantistica è fortemente oscillatoria: tocca continuamente lo zero (i nodi) per poi rimbalzare verso picchi massimi. Tuttavia, i picchi quantistici giacciono su una curva di \emph{inviluppo} che è esattamente il doppio della densità classica. Poiché il valore medio di un seno al quadrato su molti periodi è $1/2$, calcolando l'integrale della probabilità quantistica su un intervallino macroscopico $\Delta x$ (che comprende molte oscillazioni) si ottiene l'esatta compensazione: le aree coincidono perfettamente con l'integrale della controparte classica.
    \item \textbf{Picchi ai bordi:} Man mano che ci si avvicina ai punti di inversione $\pm A_n$, i picchi quantistici diventano sempre più alti e distanziati (la lunghezza d'onda di de Broglie aumenta perché l'impulso cinetico diminuisce). L'accumulo di probabilità quantistica ai bordi mima perfettamente la divergenza asintotica classica causata dalla velocità nulla della particella nei punti di inversione.
\end{itemize}
Per $n \to \infty$, gli strumenti di misura macroscopici non hanno la risoluzione sufficiente per distinguere i singoli nodi quantistici e misurano esclusivamente la media locale della densità di probabilità, restituendo l'esatta distribuzione classica di un oscillatore a molla.

%!TEX root = ../dispense_QP.tex

\chapter{Alcuni esempi di problemi quantistici}
Ci dedichiamo adesso alla soluzione di alcuni "problemi" quantistici monodimensionali\footnote{L'estensione a più dimensioni è fattibile ma esula dai fini di questo corso.}, ovvero all'analisi di situazioni tanto semplici quanto istruttive e \emph{concretamente utilizzabili}. Ci concentreremo su casistiche in cui una singola particella incontra \emph{potenziali differenti}: i.e. il potenziale a gradino, la barriera di potenziale e il pozzo di potenziale. Iniziamo dal primo in elenco.

\section{Il caso del potenziale a gradino}
Si immagini un corpo di massa $m$ che viene lanciato lungo un piano orizzontale a quota $y = 0$ contro un pendio di altezza $y = h$ e inclinato di un angolo $\varphi$ rispetto all'orizzontale. Classicamente, è possibile che il corpo \emph{risalga} l'ostacolo \textbf{se e solo se} la sua energia meccanica $E$ è maggiore della variazione di energia potenziale fra quota iniziale e finale:
\[
    E > \Delta V = mgh \quad.
\]
Sorge ora spotaneo il quesito: -\emph{la meccanica quantistica, fra tutte le sue stranezze, preserva questa legge apparentemente invalicabile?}-. La risposta alla domanda sarà sviscerata in questo paragrafo e risulterà, con tutte le probabilità, sorprendente.

\subsection{Approccio al problema}
Iniziamo con il definire un potenziale a gradino. Nel limite tale per cui l'angolo di inclinazione del pendio classico di cui si discuteva poc'anzi tenda a diventare retto, $\varphi \to \pi/2$, la barriera di potenziale che deve affrontare il corpo diventa a tutti gli effetti un \textbf{gradino}, nullo prima dell'inizio della salita e di valore $V_0 = mgh$ dopo:
\[
    V(x) = V_0 \, \theta(x-a) = \begin{cases}
        0 \quad x < a \\
        V_0 \quad x \geq a \quad,
    \end{cases}
\]
dove abbiamo definito con $\theta(x-a)$ la funzione \textit{gradino di Heaviside} centrata in $x = a$:
\[
    \theta(x-a) = \begin{cases}
    0 \quad x < a \\
    1 \quad x \geq a \quad .
    \end{cases}
\]
definita a piacere in $x = a$. Assumeremo d'ora in poi che, nel nostro problema quantistico, sia coinvolto un potenziale del tutto affine a questo, con $a = 0$. 

\subsubsection{L'hamiltoniana del sistema e S.E. stazionaria}
Il primo passo che dobbiamo fare per tradurre questo contesto classico in uno puramente quantistico è scrivere $\h{H}$:
\[
    \h{H} = \frac{\hat{p}^2}{2m} + V(x) = \frac{\hat{p}^2}{2m} V_0\theta(x) \quad,
\]
dove abbiamo già tenuto in considerazione la natura moltiplicativa dell'operatore potenziale. La struttura funzionale di $V(x)$ ci permette scrivere l'operatore di Hamilton in due modi differenti in base a se ci si trova prima o dopo il gradino:
\begin{equation}
    \h{H} = \begin{dcases}
        \frac{\hat{p}^2}{2m} \quad x < 0 \\
        \frac{\hat{p}^2}{2m} + V_0 \quad x \geq 0 \quad.
    \end{dcases}
    \label{eq:H_gradino}
\end{equation}
L'equazione di Schrödinger in forma stazionaria allora si scrive come:
\[
    \h{H}\psi = E\psi
\]
e tenendo conto della \eqref{eq:H_gradino}, possiamo suddidere la soluzione in due parti distinte caratterizzate dalle proprie funzioni d'onda $\psi_1$ e $\psi_2$:
\[
    \begin{aligned}
        & x < 0 \quad \frac{\hat{p}^2}{2m} \psi_1(x) = E \psi_1(x)\\
        & x \geq 0 \quad \left(\frac{\hat{p}^2}{2m} + V_0\right)\psi_2(x) = E \psi_2(x) \quad.
    \end{aligned}
\]
La soluzione generale $\Psi$ del problema deve essere data da \emph{entrambe} le funzione d'onda, dovutamente raccordate in occorrenza del gradino per garantire continuità e differenziabilità:
\[
    \Psi(x) = \psi_1(x)\theta(-x) + \psi_2(x)\theta(x) \quad.
\]

\subsection{Soluzione delle equazioni agli autovalori}
In entrambi i casi si tratta con equazioni simili al caso della particella libera, di energia pari a $E$ nella regione $(-\infty, 0)$ e $E-V_0$ in $[0,+\infty)$. La soluzione del problema agli autovalori deve quindi essere una combinazione lineare di onde piane, una regressiva e una progressiva\footnote{Si ricordi che il fatto che un'onda sia \emph{progressiva} o \emph{regressiva} dipende esclusivamente da se la parte temporale e spaziale della fase sono rispettivamente \emph{discordi} o \emph{concordi}. Per renderlo più chiaro, si prenda un'onda di fase $\Phi(x,t) = kx \pm \omega t$: con il passare del tempo (all'aumento di $t$), la funzione d'onda $f(\Phi)$ trasla rigidamente in \emph{avanti} (letteralmente progredisce) se la parte temporale è preceduta da segno negativo, regredisce altrimenti.}:
\[
    \psi_1(x) = Ae^{ikx} + Re^{-ikx}, \quad \psi_2(x) = Te^{iqx} + Be^{-iqx} \quad,
\]
con $q$ e $k$ i vettori d'onda della forma:
\[
    k = \frac{\sqrt{2mE}}{\hbar}, \quad q = \frac{\sqrt{2m(E-V_0)}}{\hbar} \quad.
\]  
Dal momento che il generico fattore temporale associato alle onde è dato da $\exp{-iE t/\hbar}$, abbiamo che le componenti progressive sono contraddistinte da momenti preceduti dal segno $+$, le regressive altrimenti. Analizziamo quindi termine a termine per identificarne il carattere: 
\begin{itemize}
    \item $Ae^{ikx}$: progressiva, parte da $-\infty$ e si dirige verso la barriera con quantità di moto $k$, coerentemente positiva dal momento che si muove lungo il verso di crescita delle $x$.
    \item $Re^{-ikx}$: regressiva, parte da $x = 0$ e continua sino a $-\infty$. Corrisponde a un'onda \textbf{riflessa}.
    \item $Te^{iqx}$: progressiva, vive nella regione $x \geq 0$. Corrisponde pertanto a una componente \textbf{trasmessa} della radiazione incidente sulla barriera.
    \item $Be^{-iqx}$: regressiva, vive nella regione $x \geq 0$. Se esistesse, implicherebbe l'esistenza, fisicamente insensata, di un'onda che parte da $+\infty$ per dirigersi verso la barriera. Procedermo dunque assumendo 
    \[
    B = 0 
    \]
    per garantire corenza fisica ai risultati.
\end{itemize}
Complessivamente, dunque, la soluzione assume la forma di:
\[
    \Psi(x) = \begin{cases}
    \psi_1(x) = Ae^{ikx} + Re^{-ikx} \quad x < 0 \\
    \psi_2(x) = Te^{iqx} \quad x \geq 0 \quad,
    \end{cases}
\]
con $A, R, T$ costanti d'integrazione in $\mathbb{C}$. In questo caso non possono essere ricavate tramite un processo di normalizzazione di $\Psi$, in quanto la base delle {onde piane} costruisce funzioni d'onda \emph{non normalizzabili}. In questo senso, i coefficienti in discussione rappresentano unicamente il "peso" delle singole componenti delle onde. 

Quanto possiamo fare è imporre condizioni al contorno tali da garantire continuità e derivabilità in occorrenza del gradino, dal momento che per ogni altro valore di $x$ le $\psi$ sono addirittura $\mathcal{C}^\infty$. Si noti che queste condizioni, oltre a essere necessarie ai fini di $\Psi$, sono anche indispensabili per garantire che la densità di probabilità possa soddisfare l'equazioni di continuità e per evitare discontinuità non fisiche di $\rho = |\Psi|^2$.

\subsubsection{Condizioni al contorno}
Iniziamo dalla continuità: vogliamo che $\psi_1(0) = \psi_2(0)$:
\[
    A + R = T \quad.
\]  
Per quanto riguarda la derivabilità, serve richiedere che $\psi'_1(0) = \psi'_2(0)$:
\[
    \begin{aligned}
        ik A - ik R &= iqT \\
        k(A-R) &= qT \quad.
    \end{aligned}
\]
Abbiamo quindi due equazioni equazioni e tre incognite: possiamo dunque sceglierne una e esprimere le altre due in funzione della prima. Sfruttando $A$ come parametro\footnote{-\emph{Perché proprio $A$?}- Il coefficiente $A$ altro non è che l'ampiezza dell'onda incidente. Se spariamo un fascio di $n$ elettroni contro il gradino, l'intensità del fascio è proporzionale a $n$ e ad $|A|^2$ per l'ipotesi di Born. Il rapportare tutto ad $A$ implica eseguire una sorta di \emph{normnalizzazione} che riconduce tutti i risultati a \emph{quanto si ha all'inizio}.}, otteniamo:
\begin{equation}
\begin{aligned}
    R &= A \, \frac{k-q}{k+q} = A \, \frac{\sqrt{E}-\sqrt{E-V_0}}{\sqrt{E}+\sqrt{E-V_0}} \implies \frac{R}{A} = \frac{k-q}{k+q}\\
    T &= A\, \frac{2k}{k+q} = A \, \frac{2\sqrt{E}}{\sqrt{E}+\sqrt{E-V_0}} \implies \frac{T}{A} = \frac{2k}{k+q} \quad.
\end{aligned}
\label{eq:RA}
\end{equation}
Si noti come il rapporto $R/A$ è strettamente positivo, con ciò a implicare che la componente incidente è in fase rispetto a quella riflessa. La soluzione generale della S.E. stazionaria è quindi data da:
\[
\begin{aligned}
        \Psi(x) &= A \left[ \bigg(e^{ikx} + \frac{R}{A} \,e^{-ikx}\bigg)\theta(-x) + \frac{T}{A}e^{iqx}\,\theta(x) \right] \\
        &= A \left[ \bigg(e^{ikx} + \frac{k-q}{k+q}e^{-ikx}\bigg)\theta(-x) + \frac{2k}{k+q} \,e^{iqx}\,\theta(x) \right] \quad.
\end{aligned}
\]

\subsection{Densità di corrente di probabilità}
Come \emph{prova del nove} del fatto che abbiamo trovato i corretti parametri $R$ e $T$, possiamo ora verificare che la densità di corrente di probabilità associata a $\Psi$ sia continua nell'origine, ovvero:
\[
    \psi_1(0) \rightarrow j_1(0) = j_2(0) \leftarrow \psi_2(0) \quad.
\]
Dalla definizione abbiamo:
\[
\begin{aligned}
    j_1(x) &= \frac{\hbar}{2mi}\left( \psi_1^*(x)\pdv{\psi_1(x)}{x} - \psi_1(x)\pdv{\psi_1^*(x)}{x} \right) \\
    &= \frac{\hbar}{2mi}\left[ \psi_1^*(x)\pdv{\psi_1(x)}{x} - \bigg(\psi_1^*(x)\pdv{\psi_1(x)}{x} \bigg)^*\right] \overset{(1)}{=} \frac{\hbar}{2mi} \, 2i \, \Im\left( \psi_1^*(x)\pdv{\psi_1(x)}{x} \right) \\
    &= \frac{\hbar}{m} \,\Im\left[ A^*\left( e^{-ikx}+\frac{R}{A}\,e^{ikx} \right) A \left( ike^{ikx}-ik\frac{R^*}{A^*}\,e^{-ikx} \right)  \right] \\
    &= \frac{\hbar}{m} \,\Im\left[ ik|A|^2\bigg(1 - \frac{|R|^2}{|A|^2}+ \frac{R^*}{A^*}e^{2ikx} - \frac{R}{A}e^{-2ikx}
    \bigg)  \right]
\end{aligned}
\]
dove in $(1)$ abbiamo usato la relazione $z-z^* = 2i \Im(z)$. Per continuare, notiamo che il rapporto $R/A$ è, nonchè strettamente positivo, \emph{puramente reale}. Segue che $R^*/A^* = R/A$:
\[
    \begin{aligned}
    j_1(x) &= \frac{\hbar}{m} \,\Im\left[ ik|A|^2\bigg(1 - \frac{|R|^2}{|A|^2}+ \frac{R}{A}\left(e^{2ikx} - e^{-2ikx}
    \right)\bigg)  \right] \\
    &= \frac{\hbar}{m} \,\Im\left[ ik|A|^2\bigg(1 - \frac{|R|^2}{|A|^2}+ 2i\frac{R}{A}\sin(2kx)\bigg)  \right]\\
    &= \frac{\hbar}{m} \,\Im\left[ ik\big(|A|^2 - {|R|^2}\big)  -2k|A|^2\frac{R}{A}\sin(2kx)  \right] = \frac{\hbar k}{m} \big(|A|^2 - {|R|^2}\big) \quad.
\end{aligned}
\]
Per semplificare i conti che seguono, senza perdita di generalità, assumiamo d'ora in poi che $A$ e $R$ (e anche $T$) siano parametri puramente reali. Benché sembra che nulla garantisca che questa assunzione trovi fondamento fisico, o matematico, effettivamente non esiste neppure nulla che impedisca di farlo. Ciò che è osservabile è $|\psi|^2$, associato a u n generico termine costante, ad esempio $|A|^2$. Essendo $A \in \mathbb{C}$, possiamo scrivere:
\[
    A = |A|e^{i\beta}
\]
con il termine esponenziale di fase che scompare nella valutazione del modulo quadro del numero. Imponiamo quindi $\beta = 0$ \emph{senza modificare la fisica osservabile!}. Abbiamo allora:
\[
    j_1(x) = \frac{\hbar k}{m}(A^2-R^2) = \text{costante!} 
\]
\begin{approfondimento}[title={Densità di corrente di probabilità costante}]
    \noindent
    La densità di corrente appena calcolata è una costante... ma ha senso? 
Si pensi all'equazione di continuità:
\[
    \pdv{\rho}{t} + \nabla \cdot \bm{j} = 0 \implies \pdv{\rho}{t} +  \pdv{j}{x} = 0
\]
in una dimensione. Dal momento che stiamo risolvendo la S.E. stazionaria, $\rho$ necessariamente è indipendente dal tempo e quindi $\partial_x j = 0$: $j_x$ deve essere una costante.
\end{approfondimento}

\noindent
Procedendo similmente si ricava $j_2(x)$
\[
    j_2(x) = \frac{\hbar}{m} \Im\left( Te^{-iqx}Tiqe^{iqx} \right) = \frac{\hbar q}{m}T^2 \quad,
\]
giustamente costante. Verifichiamo adesso, usando le \eqref{eq:RA}, che $j_1 = j_2$:
\[
\begin{aligned}
    &\frac{\hbar k}{m}(A^2-R^2) = \frac{\hbar q}{m}T^2 \\
    &\frac{k}{q}\left( 1-\frac{R^2}{A^2} \right) = \frac{T^2}{A^2} \implies \frac{k}{q}\left( 1-\frac{(k-q)^2}{(k+q)^2} \right) = \frac{4k^2}{(k+q)^2} \\ &\implies \frac{1}{q}\left( \frac{(k+q)^2-(k-q)^2}{(k+q)^2} \right) = \frac{4k}{(k+q)^2} \implies \frac{4k}{(k+q)^2} = \frac{4k}{(k+q)^2} \quad,
\end{aligned}
\] 
come cercavamo.

È ancora possibile dare un'ulteriore caratterizzazione ai rapporti $R/A$ e $T/A$ e, per fare ciò, introduciamo la densità di corrente di probabilità associata alle onde incidente, riflessa e trasmessa seguendo lo stesso procedimento usato sinora. Per completezza, presentiamo il conto esplicito solo per la componente incidente, $\psi_{in}(x) = Ae^{ikx}$:
\[
\begin{aligned}
    j_{in}(x) = \frac{\hbar}{2mi}\left( \psi_{in}^*(x) \pdv{\psi_{in}(x)}{x} - \psi_{in}(x) \pdv{\psi_{in}^*(x)}{x} \right) &= \frac{\hbar}{m}\Im\left( \psi_{in}^*(x) \pdv{\psi_{in}(x)}{x} \right) \\
    &= \frac{\hbar}{m}\Im\left( Ae^{-ikx} \, (ik A e^{ikx}) \right) \\
    &= \frac{\hbar}{m}\Im\left( i k A^2 \right) = \frac{\hbar k}{m} A^2 \quad,
\end{aligned}
\]
coerentemente costante e positiva come il verso di propagazione dell'onda. Scriviamo quindi:
\[
    j_{in} = \frac{\hbar k}{m} A^2, \quad j_{R} = \frac{\hbar k}{m} R^2, \quad j_T = \frac{\hbar q}{m}T^2
\]
con 
\[
    j_1 = j_{in} - j_R, \quad j_2 = j_T \quad.
\]
Rapportando le varie densità di corrente si ottiene:
\[
    \frac{j_R}{j_{in}} = \frac{R^2}{A^2} = \left( \frac{k-q}{k+q} \right)^2, \quad \frac{j_T}{j_{in}} = \frac{qT^2}{kA^2} = \frac{4kq}{(q+k)^2}, \quad \frac{j_R}{j_{in}} + \frac{j_T}{j_{in}} = 1 \quad,
\]
mostrando chiaramente come $R/A$ e $T/A$ non sono altro che le frazioni dell'onda incidente che vengono rispettivamente \emph{riflessa} e \emph{trasmessa}. Ancora, si noti come questo risultato conferma la coerenza logica del costrutto quantistico: la somma delle probabilità relative a tutti gli eventi fisicamente possibili (riflessione o attraversamento) deve necessariamente restituire la certezza assoluta ($100 \%$).

\section{Il caso della barriera di potenziale}
A differenza del gradino di potenziale, si immagini la barriera come un qualcosa di \emph{finito}: il potenziale parte da un livello zero, sale istantaneamente a un valore costante $V_0$ per poi tornare a zero. L'hamiltoniana del sistema deve constare allora di un potenziale del tipo:
\[
    V(x) = \begin{cases}
        0 \quad x < -a \\
        V_0 \quad -a \leq x \leq a \\
        0 \quad x > a \quad.
    \end{cases}
\]
L'equazione di Schrödinger si scrive:
\[
    \h{H}\psi(x) = \left[ -\frac{\hbar^2}{2m} + V(x) \right]\psi(x) = E\psi(x) \quad.
\]
La geometria del problema permette di decomporre $\psi$ in tre parti distinte, ciascuna vivente in una delle tre zone di definizione di $V(x)$: 
\begin{equation}
        \psi(x) = \psi_1(x)\, \theta(a-x) + \psi_2(x)\, \theta(a-|x|)+\psi_3(x)\, \theta(x-a) \quad,
        \label{eq:hp_psi}
\end{equation}
portando a tre diverse S.E., una per funzione d'onda. Di seguito investighiamo il caso in cui la particella è dotata di energia maggiore della barriera, $E > V_0$, mentre il viceversa è oggetto di discussione della sezione successiva sull'\emph{effetto tunnel}.

\subsection{Analisi del problema}
Abbiamo una particella libera di massa $m$ e di energia $E$ nella regione $|x| > a$ e $E-V_0$ in $|x| \leq a$. Le soluzioni delle S.E. sono quindi date da onde piane con vettori d'onda determinati dall'energia:
\begin{equation}
    \psi_1(x) = Ae^{ikx} + Re^{-ikx}, \quad \psi_2(x) = Ce^{ikx} + De^{-ikx}, \quad \psi_3(x) = Te^{ikx} \quad,
    \label{eq:psi_barriera}
\end{equation}
dove non abbiamo considerato termini regressivi non fisici in $\psi_3$. Le interpretazioni dei vari termini sono identiche al caso del gradino di potenziale e non vengo pertanto riproposte. Si noti come entro la regione della barriera la più generale espressione per $\psi_2$ comprende un termine progressivo (\emph{trasmesso} dopo l'incidenza al gradino in $x = -a$) e uno regressivo (derivante dalla \emph{riflessione} sul secondo gradino in $x = a$). 

I vettori d'onda sono dati dalle solite espressioni:
\[
    k = \frac{\sqrt{2mE}}{\hbar}, \quad q = \frac{\sqrt{2m(E-V_0)}}{\hbar} \quad.
\]
A questo punto bisogna analizzare le condizioni al contorno che risultano necessarie per garantire che la forma di $\psi(x)$ ipotizzata nella \eqref{eq:hp_psi} sia fisicamente plausibile.

\subsubsection{Condizione di continuità di $\bm{\psi(x)}$ e $\bm{\partial_x\psi(x)}$}
Imponiamo che la funzione d'onda e la sua derivata siano continue\footnote{Si ricordi che una generica funzione d'onda che descrive uno stato fisico è soluzione di un'equazione di Schrödinger e dunque deve essere \emph{almeno} $\mathcal{C}^2$ nello spazio e $\mathcal{C}^1$ nel tempo.} per $x = \pm a$. La condizione di continuità di $\psi(x)$ vuole:
\[
    \begin{aligned}
        &\psi_1(-a) = \psi_2(-a) \implies Ae^{-ika} + Re^{ika} = Ce^{-iqa} + De^{iqa} \\
        &\psi_2(a) = \psi_3(a) \implies Ce^{iqa} + De^{-iqa} = Te^{ika} \quad.
    \end{aligned}
\]
La continuità della derivata, poi:
\[
    \begin{aligned}
        &\pdv{\psi_1(-a)}{x} = \pdv{\psi_2(-a)}{x} \implies ik\left(Ae^{-ika} - Re^{ika}\right) = iq\left(Ce^{-iqa} - De^{iqa}\right) \\
        &\pdv{\psi_2(a)}{x} = \pdv{\psi_3(a)}{x} \implies iq\left(Ce^{iqa} - De^{-iqa}\right) = ikTe^{ika} \quad.
    \end{aligned}
\]
Il sistema delle quattro equazioni appena trovate nelle cinque incognite ($A, R, C, D, T$) ammette un grado di libertà. Come nel caso precedente, esprimiamo tutto come rapporto del tipo \emph{variabile}/$A$. Nel caso di $R$ e $T$ abbiamo:
\[
   \frac{R}{A} = e^{-i2ka}\frac{i (q^2-k^2) \sin(2qa)}{2kq\cos(2qa) - i{k^2+q^2}\sin(2qa)} \quad \frac{T}{A} = e^{-i2ka}\frac{2kq}{\cos2kq(2qa) - i{k^2+q^2}\sin(2qa)} \quad.
\]
Le correnti di densità di probabilità devono essere costanti per l'equazione di continuità e si ricavano come nel caso precedente.  Abbiamo quindi:
\[
\begin{aligned}
    &j_1(x) = \frac{\hbar}{2mi} \left( \psi_1^*(x)\pdv{\psi_1(x)}{x} - \psi_1(x)\pdv{\psi_1^*(x)}{x} \right) = \dots = \frac{\hbar k}{m}\left( A^2 - |R|^2 \right) \\
    &j_2(x) = \frac{\hbar q}{m}\left( |C|^2 - |D|^2 \right) \\
    &j_3(x) = \frac{\hbar k}{m} |T|^2 \quad,
\end{aligned} 
\]
con $j_1(-a) = j_2(-a)$ e $j_2(a) = j_3(a)$ per via della conservazione della densità di probabilità. Definiamo la densità incidente, riflessa e trasmessa come a partire da $j_1 = j_{in} - j_R$ e $j_3 = j_T$:
\[
    j_{in} = \frac{\hbar k}{m} A^2 \quad j_R = \frac{\hbar k}{m} |R|^2 \quad j_T = \frac{\hbar k}{m} |T|^2 \quad. 
\]
Definiamo in questo contesto il \textbf{coefficiente di riflessione} $\mathcal{R}$ e il \textbf{coefficiente di trasmissione} $\mathcal{T}$:
\[
\begin{aligned}
    \mathcal{R} &= \frac{j_R}{j_{in}} = \frac{|R|^2}{A^2} = \frac{\left( \frac{k^2-q^2}{2kq} \right)^2 \sin^2(2qa)}{1 + \left( \frac{k^2-q^2}{2kq} \right)^2 \sin^2(2qa)} \\
    \mathcal{T} &= \frac{j_T}{j_{in}} = \frac{|T|^2}{A^2} = \frac{1}{1 + \left( \frac{k^2-q^2}{2kq} \right)^2 \sin^2(2qa)} \quad,
\end{aligned}
\]
con $\mathcal{R} + \mathcal{T} = 1$. Vediamo ora cosa succede in diversi regimi energetici.

\subsubsection{Limite di alta energia}
Nel caso in cui l'energia della aprticlla sia molto maggiore dell'altezza della barriera, $E \gg V_0$, abbiamo che al primo ordine $q \approx k$ infatti:
\[
    q = \frac{\sqrt{2m(E-V_0)}}{\hbar} = \frac{\sqrt{2mE(1-V_0/E)}}{\hbar} \to \frac{\sqrt{2mE}}{\hbar} = k 
\]
per grandi energie. Abbiamo quindi che $\mathcal{R} \approx 0$: l'alta energia rende il gradino invisibile alla particella, che riesce a superarlo senza riflessioni. Vale infatti che la componente trasmessa è la totalità di quella incidente, con $\mathcal{T} \approx 1$. Parallelamente, notiamo che la funzione d'onda complessiva è data da un'unica onda piana di numero d'onda $k$. Si verifica che:
\[
    \psi_1(x) = Ae^{ikx} + Re^{-ikx} = A\left( e^{ikx} + \frac{R}{A} e^{-ikx} \right) \to Ae^{ikx} 
\] 
dal momento che $\mathcal{R} \to 0 \implies R \to 0$. Allo stesso modo si vede che $\mathcal{T} \to 1 \implies T \to A$ e:
\[
    \psi_3(x) = Te^{ikx} \to Ae^{ikx}.
\] 
Nella regione intermedia, poi, si puù dimostrare che si ottiene $D \to 0$ e $C \to A$ portando a un risultato complessivo di:
\[
    \psi(x) = Ae^{ikx} \quad,
\]
coerentemente con l'ipotesi, informale, di \emph{cecità} della particella rispetto alla barriera garantita dalla grande energia relativa.

\subsubsection{Limite di bassa energia}
Se si considerano energia nell'ordine dell'altezza della barriera, $E \approx V_0$, il numero d'onda nella regione interna tende a zero, $q \to 0$ e i rapporti $R/A$ e $T/A$ tendono a:
\[
    \frac{R}{A} \approx \frac{ika}{ika - 1} \quad \frac{T}{A} \approx \frac{1}{1-ika} \quad,
\] 
mentre i coefficienti di trasmissione e riflessione risultano:
\[
    \mathcal{R} \approx \frac{k^2a^2}{1+k^2a^2} \quad \mathcal{T} \approx \frac{1}{1+k^2a^2} \quad.
\]  
Dal momento che il numero d'onda è inversamente proporzionale alla lunghezza d'onda, possiamo ancora distinguere due subregimi a queste basse energie, contraddistinti rispettivamente da una barriera \emph{larga} o \emph{stretta}. Si ha una barriera larga quando $ka \gg 1$, ovvero se $\lambda / a \gg 2\pi$ (la lunghezza d'onda è molto maggiore dello spessore fisico della barriera), stretta con $\lambda / a \ll 2\pi$. In queste condizioni, abbiamo:
\[
    \begin{cases}
    \text{barriera larga} \implies \mathcal{R} \approx 1, \quad \mathcal{T} \approx 0 \\
    \text{barriera stretta} \implies \mathcal{T} \approx 1, \quad \mathcal{R} \approx 0 \quad.
    \end{cases}
\]
Si noti come questi risultati siano fisicamente coerenti: una barriera larga è, letteralmente, più complicata da superare a basse energie, viceversa se l'ostacolo è stretto. Nel primo caso, infatti, il coefficiente di riflessione tende a uno, nel secondo quello di trasmissione si avvicina ad essere unitario.

\subsubsection{Trasmissione totale}
Un effetto \textbf{puramente quantistico} e difficilmente interpretabile in un framework classico è quello dato dalla \emph{trasmissione totale dell'onda} in occorrenza di alcuni particolari valori dei numeri d'onda. Per $2qa = 2\pi n, n \in \mathbb{Z}$ i coefficienti diventano $\mathcal{R} = 0$ e $\mathcal{T} = 1$: trasmissione totale dell'onda attraverso la barriera.

\section{L'effetto tunnel quantistico}
Questa sezione apre a una trattazione di carattere strettamente quantistico: qualsiasi reminiscenza classica deve lasciare il passo ai risultati, piuttosto controintuitivi, che seguono. Nel caso in cui l'energia della particella sia minore dell'altezza della barriera, $E < V_0$, la regione $x > -a$ è \emph{classicamente proibita}. Se nelle regioni esterne alla barriera la funzione d'onda assume la stessa forma proposta in \eqref{eq:psi_barriera}, la soluzione della S.E. nella regione intermedia, $-a \leq x \leq a$, conduce a un risultato differente:
\[  
    -\frac{\hbar^2}{2m}\pdv[2]{\psi_2(x)}{x} + V_0\psi_2(x) = E \psi_2(x) \implies \psi_2(x) = Ce^{-qx} + De^{qx} \quad,
\]
con $q = \sqrt{2m(V_0-E)}/\hbar$. La presenza di una soluzione non nulla della S.E. nella region proibita dimostra matematicamente come il \emph{tunnel} classicamente impossibile sia una \emph{realtà} fisica quantistica. 

\subsection{Condizioni di continuità e coefficienti di trasmissione eriflessione}
Le condizioni di continuità in $x = \pm a$ portano al seguente sistema: 
\[
\begin{dcases}
    Ae^{-ika} + Re^{ika} = Ce^{qa} + De^{-qa} \\
    Ce^{-qa} + De^{qa} = Te^{ika} \\
    ik\big(Ae^{-ika} - Re^{ika}\big) = q\big( -Ce^{qa} + De^{-qa} \big) \\
    q\big(-Ce^{-qa}+De^{qa}\big) = ikTe^{ika} 
\end{dcases}
\]
che conduce alle seguenti soluzioni (espresse come di consueto come rapporto rispetto al coefficiente $A$):
\[
    \frac{T}{A} = \frac{e^{-i2ka}}{\cosh(2qa) - i\frac{k^2 - q^2}{2kq}\sinh(2qa)} \quad
    \frac{R}{A} = \frac{-i\frac{k^2 + q^2}{2kq}\sinh(2qa) e^{-i2ka}}{\cosh(2qa) - i\frac{k^2 - q^2}{2kq}\sinh(2qa)} 
\]
e ai seguenti coefficienti:
\[
    \mathcal{R} = \frac{|R|^2}{|A|^2} = \frac{\left( \frac{k^2 + q^2}{2kq} \right)^2 \sinh^2(2qa)}{1 + \left( \frac{k^2 + q^2}{2kq} \right)^2 \sinh^2(2qa)} \quad
    \mathcal{T} = \frac{|T|^2}{|A|^2} = \frac{1}{1 + \left( \frac{k^2 + q^2}{2kq} \right)^2 \sinh^2(2qa)} 
\]
con il coefficiente di trasmissione evidentemente non nullo. Esclusivamente nel caso in cui $E \ll V_0$ (e quindi $q \gg k$) si ritrova il limite classico della totale riflessione, con $\mathcal{R} \approx 1$ e $\mathcal{T} \approx 0$, $R/A \approx -e^{2iak}$ e $T/A \approx 0$. Per avere una visione più concreta dei risultati appena ottenuti, si possono ancora sostituire le relazioni per i numeri d'onda nel coefficiente di riflessione e ottenere:
\[
    \mathcal{T} = \frac{1}{1 + \frac{V_0^2}{4E(V_0 - E)} \sinh^2(2qa)} \quad. 
\]  

\begin{approfondimento}[title={Approssimazione di barriera spessa}]
    \noindent
Nel caso di una \emph{barriera opaca o spessa}, ovvero quando l'attenuazione dell'onda all'interno della regione classicamente proibita è molto forte ($2qa \gg 1$), è possibile ricavare una fondamentale espressione approssimata. Per argomenti grandi, il seno iperbolico è dominato dall'esponenziale crescente:
\[
    \sinh(2qa) = \frac{e^{2qa} - e^{-2qa}}{2} \approx \frac{e^{2qa}}{2} \implies \sinh^2(2qa) \approx \frac{e^{4qa}}{4} \,.
\]
Sostituendo questa approssimazione nell'espressione esatta di $\mathcal{T}$, si ha:
\[
    \mathcal{T} \approx \frac{1}{1 + \frac{V_0^2}{16E(V_0 - E)} e^{4qa}} \,.
\]
Essendo $e^{4qa} \gg 1$, il termine esponenziale al denominatore risulta preponderante rispetto all'unità, che può essere trascurata. Invertendo la frazione si ottiene:
\[
\begin{aligned}
    \mathcal{T} &\approx \frac{16E(V_0 - E)}{V_0^2} e^{-4qa} \\
                &= \frac{16E(V_0 - E)}{V_0^2} \exp\left(-4a\frac{\sqrt{2m(V_0 - E)}}{\hbar}\right) \quad .
\end{aligned}
\]
\end{approfondimento}
%!TEX root = ../dispense_QP.tex

\chapter{Il principio di indeterminazione e le sue conseguenze}
La discussione che segue presenta il suo cardine nel celeberrimo \textbf{Principio di indeterminazione di Heisenberg} che dimostreremo entro breve. Le conseguenze fisiche che ne conseguono sono \emph{devastanti}: un assaggio lo abbiamo già avuto durante l'analisi della posizione della massa oscillante nel caso dell'oscillatore armonico. In quel contesto avevamo scoperto che non è possibile conoscere simultaneamente con errore nullo sia la posizione sia la quantità di moto associate alla particella, a meno di un fattore di scarto $\hbar/2$. Nel seguito proveremo a formalizzare quella che era stata solo un'intuizione e riprenderemo gli autostati dell'oscillatore armonico, imparando a identificarli come \textbf{stati coerenti}.

\section{Il Principio di indeterminazione di Heisenberg}
L'idea di partenza è semplice: se classicamente è possibile conoscere \emph{tutto} e \emph{allo stesso tempo} su una particella in moto, cosa mi garantisce che lo stesso è fattibile nel mondo quantistico? Effettivamente nulla permette di avere tanta \emph{hybris} da avanzare una tale assunzione, per di più smentita da uno dei più grandi che la meccanica quantistica abbia conosciuto: Werner Karl Heisenberg. Discuteremo più avanti il suo approccio alla fisica dei quanti, sufficientemente diverso da quello dei compagni Schrödinger e Dirac: per ora ci \emph{limitiamo} a cercare di capire se e in quali condizioni possiamo mimare la certezza classica della conoscenza istantanea e assoluta di un sistema. 

\subsection{Alcune considerazioni preliminari}
Scegliamo inizialmente due osservabili fisiche, $A_1$ e $A_2$, a cui sono associate gli operatori hermitiani $\h{A}_1$ e $\h{A}_2$. Assumiamo che il commutatore dei due operatori sia non nullo e richiediamo che valga:
\[
    [\h{A}_1, \h{A}_2] = i\h{G} \quad,
\]
con $\h{G}$ un generico operatore. Questa apparentemente particolare forma non lede in alcun modo la generalità dell'approccio al problema: quanto ci siamo limitati a fare è stato imporre il commutatore fra due operatori, che in generale è proprio un operatore che possiamo chiamare $\h{C}$, uguale a $i\h{G}$. Questo escamotage, in ogni caso, permette di garantire che $\h{G}$ sia anch'esso hermitiano, infatti:
\[
\begin{aligned}
    \hd{G} = \left(-i[\h{A}_1, \h{A}_2]\right)^\dagger = i(\h{A}_1\h{A}_2 - \h{A}_2\h{A}_1)^\dagger &= i(\hd{A}_2\hd{A}_1 - \hd{A}_1\hd{A}_2) \\ &\giustifica[=]{1} i(\h{A}_2\h{A}_1 - \h{A}_1\h{A}_2) = -i[\h{A}_1, \h{A}_2] = \h{G} \quad,
\end{aligned}
\]
dove in $(1)$ abbiamo sfruttato il fatto che i due opratori sono hermitiani. Definiamo a questo punto due nuovi operatori costruiti a partire da quelli di partenza tramite cui mimare la definizione classica di varianza\footnote{Si ricordi che $Var(X) = E[(X-\mu_X)^2]$, con $X$ una generica variabile aleatoria e $\mu_X$ la sua media.}:
\[
    \h{D}_1 = \h{A}_1 - \avg{\h{A}_1}\mathds{1}, \quad \h{D}_2 = \h{A}_2 - \avg{\h{A}_2}\mathds{1} \quad,
\]
che sono a loro volta hermitiani perché $\h{A}_i$ e $\mathds{1}$ lo sono. Si noti che i $\h{D}_i$ sono satti costruiti come \emph{osservabile} meno la sua media: la media del loro quadrato deve dunque essere la controparte quantistica della varianza di $\h{A}_i$:
\begin{equation}
\begin{aligned}
        \avg{\h{D}_i^2} = \avg{\left( \h{A}_i - \avg{\h{A}_i} \mathds{1} \right)^2} &= \avg{\hat{A}_i^2 - 2\h{A}_i\avg{\h{A}_i} + \avg{\h{A}_i}^2 \mathds{1}} \\ &\giustifica[=]{2} \avg{\h{A}_i^2} - 2\avg{\h{A}_i}^2 + \avg{\h{A}_i}^2 = \avg{\h{A}_i^2} - \avg{\h{A}_i}^2 = \sigma_i^2 = (\Delta \h{A}_i)^2 \quad,
\end{aligned}
\label{eq:varianza_A}
\end{equation}
dove in $(2)$ abbiamo sfruttato la linearità del prodotto scalare. Si noti ancora che:
\[
\begin{aligned}
    [\h{D}_1, \h{D}_2] &= [\h{A}_1 - \avg{\h{A}_1}\mathds{1}, \h{A}_2 - \avg{\h{A}_2}\mathds{1}] \\ &= [\h{A}_1, \h{A}_2] - \avg{A_1}[\mathds{1}, \h{A}_2] - \avg{A_2}[\mathds{1}, \h{A}_1] + \avg{A_1}\avg{A_2}[\mathds{1}, \mathds{1}] = [\h{A}_1, \h{A}_2] \quad,
\end{aligned}
\]
dal momento che qualsiasi operatore commuta con l'identità. 

\subsection{Derivazione della disuguaglianza di Robertson}
Se nella meccanica classica le parentesi di Poisson fra due funzioni dello spazio delle fasi rappresentano un'indicatore di \emph{quanto l'evoluzione della prima influenzi l'evoluzione della seconda} e viceversa, la regola di quantizzazione ci permette di estendere questa descrizione geometrica alla meccanica quantistica. Il commutatore diventa così lo strumento tramite cui discernere quanto due operatori (e dunque le osservabili ad essi associate) siano \emph{compatibili}. 

Il modulo quadro del commuatatore fra due operatori deve quindi resitutire una \emph{misura} di questa fantomatica compatibilità. La quantità d'interesse diventa allora:
\[
    \left| \avg{\h{G}} \right|^2 = \left| \avg{-i[\h{A}_1, \h{A}_2]} \right|^2 \quad.
\]
Espandendo l'ugaglianza si ottiene:
\[
        \left| \avg{\h{G}} \right|^2 = \left| \avg{[\h{A}_1, \h{A}_2]} \right|^2 = \left| \avg{[\h{D}_1, \h{D}_2]} \right|^2 = \left| \avg{\h{D}_1\h{D}_2 - \h{D}_2\h{D}_1} \right|^2 = \left|\avg{\h{D}_1\h{D}_2} - \avg{\h{D}_2\h{D}_1} \right|^2 \quad.
\]
È possibile sfruttare l'hermitianità di $\h{D}_i$ per riscrivere l'ultimo termine. Vale infatti:
\[
\begin{aligned}
    \avg{\h{D}_2\h{D}_1} = \left( \psi, \h{D}_2\h{D}_1 \psi \right) &= \left(\hd{D}_2 \psi, \h{D}_1 \psi \right) = \left(\h{D}_2 \psi, \h{D}_1 \psi \right) \\
    &= \left(\h{D}_1 \psi, \h{D}_2 \psi \right)^* = \left(\psi, \hd{D}_1\h{D}_2 \psi \right)^* = \left(\psi, \h{D}_1\h{D}_2 \psi \right)^* = \avg{\h{D}_1\h{D}_2}^* \quad.
\end{aligned}
\]
Sfruttando questo risultato otteniamo:
\[
    \begin{aligned}
        \left| \avg{\h{G}} \right|^2 = \left|\avg{\h{D}_1\h{D}_2} - \avg{\h{D}_1\h{D}_2}^* \right|^2 = \left|2i\Im\left(\avg{\h{D}_1\h{D}_2}\right) \right|^2 = 4\left|\Im\left(\avg{\h{D}_1\h{D}_2}\right) \right|^2 \quad.
    \end{aligned}
\] 
Il valore di aspettazione del prodotto $\h{D}_1 \h{D}_2$ è un numero complesso del tipo $\alpha + i\beta$ da cui:
\[
    \begin{aligned}
        \left| \avg{\h{G}} \right|^2 = 4\left|\Im\left(\alpha + i\beta\right) \right|^2 = 4\beta^2 \leq 4(\alpha^2 + \beta^2) = 4\left| \avg{\h{D}_1\h{D}_2} \right|^2 \quad.
    \end{aligned}
\]
Si noti l'introduzione del segno di minore o uguale! Questa è l'origine della fantomatica \textbf{formulazione di Robertson} del principio di indeterminazione. Continuando con la definizione di valore di aspettazione abbiamo:
\[
    \begin{aligned}
        \left| \avg{\h{G}} \right|^2 \leq 4\left| (\psi, \h{D}_1\h{D}_2\psi) \right|^2 = 4\left| (\h{D}_1\psi, \h{D}_2\psi) \right|^2 = 4\left| (\phi_1, \phi_2) \right|^2 \quad,
    \end{aligned}
\]
dove si sono chiamati $\phi_1$ e $\phi_2$ i due stati prodotti dall'applicazioni di $\h{D}_i$ su $\psi$. La \emph{disuguaglianza di Cauchy-Schwarz} permette di continuare:
\[
    \begin{aligned}
        \left| \avg{\h{G}} \right|^2 \leq 4(\phi_1, \phi_1)(\phi_2, \phi_2) = 4(\h{D}_1\psi, \h{D}_1\psi)(\h{D}_2\psi, \h{D}_2\psi) = 4(\psi, \h{D}_1^2 \psi)(\psi, \h{D}_2^2 \psi) = 4\avg{\h{D}_1^2}\avg{\h{D}_2^2}
    \end{aligned}
\]
e grazie alla \eqref{eq:varianza_A}, con qualche rimaneggiamento di carattere algebrico, otteniamo la formula di Robertson:
\begin{equation}
    \begin{aligned}
    \left| \avg{\h{G}} \right|^2 &\leq 4(\Delta \h{A}_1)^2 \, (\Delta \h{A}_2)^2 \\
    \implies (\Delta \h{A}_1)^2 \, (\Delta \h{A}_2)^2 &\geq \frac{1}{4} \left| \avg{[\h{A}_1, \h{A}_2]} \right|^2 \quad.
    \end{aligned}
\end{equation}
Gli stati per cui la disuguaglianza diventa un'uguaglianza sono detti \textbf{stati coerenti}.

\subsection{interpretazione fisica}
La relazione appena trovata testimonia come la conoscenza di due osservabili associate a operatori non commutanti è intrinsecamente affetta da un errore che non permette di avere informazioni \emph{esatte} e \emph{simultanee} riguardo entrambe le grandezze studiate. L'incertezza in questione, ancora, non può essere ridotta arbitrariamente: nel limite in cui $(\Delta \h{A}_1)^2 \to 0$, necessariamente $(\Delta \h{A}_1)^2 \to \infty$. 

In soldoni, nel caso in esame di operatori non commutanti: diminuire l'incertezza sulla misura di una grandezza impone la perdita di informazioni sulla seconda. 

\subsubsection{L'esempio della particella libera}
Nel caso della particella libera ci siamo concentrati sullo studio degli operatori posizione $\hat{x}_k$ e momento $\hat{p}_j$, declinati nelle loro componenti. Cerchiamo di applicare il Principio di indeterminazione e vedere dove ci porta:
\[
    (\Delta \hat{x}_k)^2 (\Delta \hat{p}_j)^2 \leq \frac{1}{4} \left| \avg{[\hat{p}_j, \hat{x}_k]}\right|^2 = \frac{1}{4} \left| \avg{i\hbar \delta_{jk}\mathds{1}}\right|^2 = \frac{\hbar^2 \delta_{jk}}{4} \quad.
\]
Questo risultato ci conferma il fatto che, se è possibile conoscere con assoluta precisione e simultaneamente posizione e momento relativi a due \emph{direzioni differenti}, il collasso alla stessa componente porta a:
\[
    (\Delta \hat{x}_k)^2 (\Delta \hat{p}_j)^2 \leq \frac{\hbar^2}{4} \implies \Delta \hat{x}_k \Delta \hat{p}_j \leq \frac{\hbar}{2} 
\]
come conseguenza della natura prettamente non deterministica, bensì probabilistica, della teoria quantistica. Il percorso seguito da una particella libera, quindi, non è dato da una \emph{linea retta} in quanto questo implicherebbe la conoscenza della legge oraria,
\[
    x(t) = x_0 + \frac{p_x}{m} t \quad,
\]
in disacccordo con Heisenberg. Quanto si conosce, come discusso nella sezione dedicata, è il fatto che il massimo della probabilità è localizzato laddove classicamente si riuscirebbe a prevedere la posizione del corpuscolo che, almeno nei limiti quantistici, si muove entro una sorta di percorso \emph{ombreggiato}.


\section{Teorema sugli operatori commutanti}
Di grande interesse a questo punto della trattazione è capire se, e quando, due operatori commutano. Il seguente teorema ci viene in soccorso.

\begin{teorema}[Operatori hermitiani commutanti]
    Siano $\h{A}$ e $\h{B}$ due operatori hermitiani e \emph{commutanti}, ovvero per cui vale:
    \[
        [\h{A}, \h{B}] = \h{A}\h{B} - \h{A}\h{B} = 0 \quad,
    \]
    allora $\h{A}$ e $\h{B}$ condividono almeno un insieme di autofunzioni.
\end{teorema}
La dimostrazione si articola in due casistiche, quella relativa a operatori con spettro degenere e non. Iniziamo dall'evenienza più semplice.

\subsubsection{Caso non degenere}
Si ricordi che un operatore ha spettro non degenere se per ogni autostato esiste un unico autovalore, ovvero se
    \[
        \forall \alpha : \h{A}\psi = \alpha \psi \quad \exists! \, \psi \quad.
    \] 
\begin{proof}
    \textbf{(1. Caso non degenere)} 
    Sia $\h{A}$ un operatore con spettro non degenere e $\h{B}$ un operatore \textbf{qualunque}. Si noti come non è \emph{assolutamente necessario} imporre che anche lo spettro di $\h{B}$ sia non degenere (segue esempio al termine della dimostrazione). Sia $\psi$ un generico autostato di $\h{A}$ con associato l'autovalore $\alpha$, allora:
    \[
        \h{A}(\h{B}\psi) = (\h{A}\h{B})\psi \giustifica[=]{1} (\h{B}\h{A})\psi = \h{B}(\h{A}\psi) = \alpha \h{B}\psi \quad,
    \]
    dove in $(1)$ abbiamo sfruttato il fatto che $[\h{A}, \h{B}] = 0$. Notiamo allora che anche lo stato $\h{B}\psi$ è autostato di $\h{A}$. Essendo lo spettro di quest'ultimo non degenere, però, $\h{B}\psi$ non può che essere proporzionale a $\psi$:
    \[
        \h{B}\psi = \beta \psi \quad.
    \]
    Ne consegue che $\psi$ deve essere autostato anche di $\h{B}$ associato all'autovalore $\beta$. Dal momento che questo ragionamento si può estendere a tutte le autofunzioni di $\h{A}$, allora ciascuna di esse è anche autostato di $\h{B}$. 
\end{proof}
Come anticipato, la richiesta che $\h{B}$ abbia spettro non degenere è superflua. Si prenda infatti un qualsiasi $\h{A}$ non degenere, allora vale sempre:
\[
    [\h{A}, \mathds{1}] = \h{A}\mathds{1} - \mathds{1}\h{A} = \h{A} - \h{A} = 0 \quad,
\]
anche se $\mathds{1}$ ammette come unico autovalore $\alpha = 1$. 

\subsubsection{Caso degenere}
Se entrambi gli operatori analizzati avessero spettro degenere, ovvero se fossero caratterizzati da più autostati associati al medesimo autovalore, non sarebbe possibile concludere \emph{facilmente} la dimostrazione come in precedenza. 

\begin{proof}
    \textbf{(2. Caso degenere)} Siano $\h{A}$ e $\h{B}$ due operatori commutanti con spettro degenere. Concentriamoci su $\h{A}$ e scegliamo, senza perdita di generalità, il generico autovalore $\alpha$ a cui sono associate $n$ autofunzioni. L'insieme
    \[
        \mathcal{F}_n \coloneq \left\{ \psi_i, i \in \{1, n\} : \h{A}\psi_i = \alpha \psi_i  \right\}
    \]
    è un sottospazio vettoriale\footnote{Non è difficile verificarlo: contiene certamente l'elemento nullo ed è formato da vettori appartenenti a uno spazio vettoriale, lo spettro di $\h{A}$.}. Procedendo come nel caso non degenere, scriviamo:
    \[
        \h{A}(\h{B}\psi_i) = (\h{A}\h{B})\psi_i = (\h{B}\h{A})\psi_i = \h{B}(\h{A}\psi_i) = \alpha \h{B}\psi_i \quad.
    \]
    A questo punto, però, non è possibile concludere che lo stato $\h{B}\psi_i$ sia proporzionale a $\psi_i$: nell'evenienza precedente, il sottospazio generato dall'unico autostato associato ad $\alpha$ aveva dimensione 1, di conseguenza qualsiasi autofunzione che avrebbe condiviso $\alpha$ come autovalore sarebbe dovuta essere proporzionale a $\psi$. In questo caso, $\dim{\mathcal{F}} > 1$ per ipotesi: quanto possiamo certamente assumere è che $\h{B}\psi_i$ si possa esprimere come combinazione lineare degli elementi di $\mathcal{F}$:
    \[
        \h{B}\psi_i = \sum_j C_{ij}\psi_j,\quad C_{ij} = (\psi_j, \h{B}\psi_i), \quad \psi_j \in \mathcal{F} \quad.
    \]
    Se raccogliamo in una matrice $\mathcal{C}$ tutti gli elementi $C_{ij}$, ci rendiamo conto del fatto che $\mathcal{C}$ è hermitiana, infatti:
    \[
        C_{ij}^* = (\psi_j, \h{B}\psi_i)^* = (\h{B}\psi_i, \psi_j) =  (\psi_i, \hd{B}\psi_j) \giustifica[=]{1} (\psi_i, \h{B}\psi_j) = C_{ji}
    \]
    dove in $(1)$ abbiamo sfruttato il fatto che $\h{B}$ è hermitiano. Il \emph{Teorema spettrale} garantisce che $\mathcal{C}$ possa essere diagonalizzata tramite matrici diagonalizzanti unitarie (ovvero per cui $\mathcal{U}^{-1} =\mathcal{U}^{\dagger} $):
    \[  
        \mathcal{C} = \mathcal{U} \mathcal{D} \mathcal{U}^{\dagger}
    \]
    con $\mathcal{D}$ diagonale. Riscrivendo l'uguaglianza in componenti si ottiene:
    \[
        C_{ij} = \sum_m \sum_k \mathcal{U}_{im} \mathcal{D}_{mk} (\mathcal{U}^{\dagger})_{kj} \giustifica[=]{2} \sum_m \sum_k \mathcal{U}_{im} d_m \delta_{mk} \mathcal{U}^*_{jk} = \sum_m \mathcal{U}_{im} d_m \mathcal{U}^*_{jm} \quad,
    \] 
    dove in $(2)$ si è sfrutatto il fatto che $\mathcal{D}$ è diagonale. Sostituendo nella relazione per $\h{B}\psi_i$ si ottiene:
    \[
        \h{B}\psi_i = \sum_j\sum_m \mathcal{U}_{im} d_m \mathcal{U}^*_{jm} \psi_j = \sum_m \mathcal{U}_{im} d_m \underbrace{\sum_j \mathcal{U}^*_{jm} \psi_j}_{\phi_m} = \sum_m \mathcal{U}_{im} d_m \phi_m
    \]
    dove si è identificato con $\phi_m$ il nuovo stato ottenuto tramite C.L. delle $\psi_j \in \mathcal{F}$. Si noti come $\phi_m$ deve necessariamente essere autostato di $\h{A}$ essendo costruito a partire dfa autostati di $\h{A}$. L'ultima relazione può suggerire che $\phi$ sia autostato anche di $\h{B}$ dal momento che l'applicazione dell'operatore sullo stato $\psi_i$ genera una C.L. delle $\phi_m$. Verifichiamolo:
    \[
        \h{B}\phi_m = \h{B} \sum_j \mathcal{U}^*_{jm} \psi_j = \sum_j \mathcal{U}^*_{jm} \h{B}\psi_j = \sum_j \mathcal{U}^*_{jm} \sum_k \mathcal{U}_{jk} d_k \phi_k =\sum_k\underbrace{\sum_j \mathcal{U}^*_{jm}\mathcal{U}_{jk}}_{\delta_{mk}} d_k \phi_k = d_m \phi_m \quad,
    \]
    con $\phi_m$ che è così autofunzione sia di $\h{B}$ sia di $\h{A}$. Si noti che sono state trovate $m$ autostati $\phi$, nello stesso numero della cardinalità dell'insieme delle autofunzioni degenri di $\h{A}$. Dal momento che questo approccio può essere esteso per ogni insieme $\mathcal{F}$, segue l'ipotesi.
\end{proof}

\subsection{Insieme completo di operatori}
Nell'analizzare un problema quantistico, un'informazione di notevole importanza è data dalla conoscenza di tutti gli operatori associati a osservabili del sistema che commutanto fra loro: questa caratteristica garantisce di poter misurare con precisione assoluta e simultaneamente le grandezze ad essi associate e di descrivere, in questo modo, il sistema nel modo più \emph{completo} possibile. 

Definiamo in questo senso un generico insieme di operatori $\mathcal{I}$ come \textbf{insieme completo di operatori commutanti} (ICOC) di un sistema se contiene il numero massimo di operatori che commutano fra loro e se la base di autostati condivisa è non degenere\footnote{Questa richiesta non è un vezzo matematico, bensì una necessità fondata nella \emph{teoria della misura}.}. Si noti come è spesso possibile trovare più di un ICOC per ogni sistema: la scelta di quale usare spetta a colui che lo analizza e solo alcuni portano a lavorare con sole quantità conservate.

\section{Stati coerenti}
Passiamo adesso a discutere di quegli stati che minimizzano la formulazione di Robertson del Principio di Heinseberg: gli \emph{stati coerenti}. La loro rilevanza nello studio della materia è data dal fatto che il loro comportamento è quanto di più affine si possa trovare a quello proprio di un sistema classico. 

Gli stati coerenti possono essere ricavati in diversi modi, qui si presenta la definizione classica che si appoggia al già introdotto \textbf{operatore di annichilazione} $\hat{a}$. Si dice che uno stato $\psi_z$ è uno stato coerente se è un autostato di $\hat{a}$:
\begin{equation}
    \hat{a} \psi_z = z\psi_z \quad.
    \label{eq:stati_coerenti}
\end{equation}

\subsection{Soluzione dell'equazione differenziale}
Armati di pazienza, ci immergiamo adesso nella soluzione della \eqref{eq:stati_coerenti}, ricordando la definizione operazionale di $\hat{a}$:
\[
    \hat{a} = \sqrt{\frac{m\omega}{2\hbar}} \left( \hat{x} + i\frac{\hat{p}}{m\omega} \right) \quad.
\]
Abbiamo:
\[
    \sqrt{\frac{m\omega}{2\hbar}} \left( x + \frac{\hbar}{m\omega} \dv{}{x}\right)\psi_z(x) = z\psi_z(x) \quad.
\]
Si può dimostrare che la soluzione ha la forma di:
\[
    \psi_z(x) = Ce^{-\sigma x^2 + i\eta x}
\] 
che, sostituita nell'equazione, permette di trovare le costanti $\sigma$ e $\eta$:
\[
    \begin{aligned}
        \sqrt{\frac{m\omega}{2\hbar}} \left( x + \frac{\hbar}{m\omega} (-2\sigma x + i\eta)\right)&\psi_z(x) = \sqrt{\frac{m\omega}{2\hbar}} \left( x - \frac{2\hbar}{m\omega}\sigma  + \frac{\hbar}{m\omega}i\eta \right)\psi_z(x) = z\psi_z(x)\\
        &\implies \begin{dcases}
            \frac{2\hbar}{m\omega}\sigma = 1 \\
            \frac{\hbar}{m\omega}i\eta = z 
        \end{dcases}
        \implies \begin{dcases}
        \sigma = \frac{m\omega}{2\hbar} = \frac{1}{2\lambda^2} \\
        i\eta = \frac{2m \omega}{\hbar}z = \frac{\sqrt{2}}{\lambda}z 
        \end{dcases}
    \end{aligned}
\] 
dove abbiamo riconosciuto la lunghezza caratteristica dell'oscillatore armonico $\lambda$. Segue che lo stato coerente può essere scritto come:
\[
    \psi_z(x) = Ce^{-\frac{x^2}{2\lambda^2} + \frac{\sqrt{2}z}{\lambda}x} \quad.
\]
Per trovare la cosatante $C$ si impone la normalizzazione:
\[
    (\psi_z, \psi_z) = 1 \implies 1 = |C|^2 \int_{\mathbb{R}} \dd{x} \, \exp\left[ -\frac{x^2}{\lambda^2} + \frac{\sqrt{2}x}{\lambda}(z+z^*) \right] \quad,
\]
con $z + z^* = 2\Re(z) = 2v$. Completando il quadrato e effettuando il cambio di variabile $y = x/\lambda$:
\[
    1 = |C|^2 \lambda \int_{\mathbb{R}} \dd{y} \, e^{ -y^2 + 2\sqrt{2}yv } = |C|^2 \lambda \int_{\mathbb{R}} \dd{y} e^{-\left(y-\sqrt{2}v\right)^2+2v^2} = |C|^2 \lambda e^{2v^2} \int_{\mathbb{R}} \dd{y} e^{-\left(y-\sqrt{2}v\right)^2} \quad,
\]  
dove riconosciamo l'integrale gaussiano come uguale a $\sqrt{\pi}$. Imponendo un fattore di fase nullo per $C$ abbiamo infine:
\[
    C^2\lambda e^{2v^2}\sqrt{pi} = 1 \implies C = \frac{e^{-v^2}}{\sqrt{\lambda\sqrt{\pi}}} \quad.
\]
Otteniamo infine:
\begin{equation}
    \psi_z(x) = \frac{e^{-\frac{x^2}{2\lambda^2} + \frac{\sqrt{2}z}{\lambda}x-v^2}}{{\sqrt{\lambda\sqrt{\pi}}}} \quad.
    \label{eq:stati_coerenti}
\end{equation}

\subsubsection{Significato fisico dell'autovalore $\bm{z}$}
L'interessante significato fisico di $z$ diventa palese quando si proietta lo stato $\hat{a}\psi_z$ su $\psi_z$, infatti:
\[
    \sqrt{\frac{m\omega}{2\hbar}} \left( \hat{x} + i\frac{\hat{p}}{m\omega} \right)\psi_z = z\psi_z \implies \sqrt{\frac{m\omega}{2\hbar}} \left[ (\psi_z,\hat{x}\psi_z) + i\frac{(\psi_z,\hat{p}\psi_z)}{m\omega} \right]= z(\psi_z, \psi_z) \quad.
\]
Riconoscendo adesso la definizione di valore di aspettazione e sfruttando la normalizzazione di $\psi_z$ abbiamo:
\[
    z = \sqrt{\frac{m\omega}{2\hbar}} \left( \avg{\hat{x}} + i\frac{\avg{\hat{p}}}{m\omega} \right) \implies \begin{dcases}
    \Re(z) = \sqrt{\frac{m\omega}{2\hbar}}\avg{\hat{x}} = \frac{\avg{\hat{x}}}{\lambda\sqrt{2}}\\
    \Im(z) = \frac{\avg{\hat{p}}}{\sqrt{2\hbar m\omega}} = \frac{\lambda \avg{\hat{p}}}{\hbar \sqrt{2}} \quad.
    \end{dcases}
\]
Questo permette di riscrivere in modo più elegante e fisicamente significativo la \eqref{eq:stati_coerenti}:
\[
\begin{aligned}
        \psi_z(x) &= \frac{1}{\sqrt{\lambda\sqrt{\pi}}} \exp\left[ -\frac{x^2}{2\lambda^2} + \frac{\sqrt{2}x}{\lambda} \left(\frac{\avg{\hat{x}}}{\lambda\sqrt{2}} + i\frac{\lambda \avg{\hat{p}}}{\hbar \sqrt{2}}  \right) - \frac{\avg{\hat{x}}^2}{2\lambda^2} \right] \\
        &= \frac{1}{\sqrt{\lambda\sqrt{\pi}}} e^{i\frac{\avg{\hat{p}}}{\hbar}x} \exp\left[ -\frac{x^2 + \avg{\hat{x}}^2 - 2x\avg{\hat{x}}}{2\lambda^2} \right] = \frac{1}{\sqrt{\lambda\sqrt{\pi}}} e^{i\frac{\avg{\hat{p}}}{\hbar}x} e^{-\frac{(x-\avg{\hat{x}})^2}{2\lambda^2}} \quad,
\end{aligned}
\]
dimostrando come uno stato coerente non è altro che la combinazione di un'onda piana di numero d'onda $\avg{\hat{p}}/\hbar$ e di uno stato localizzato gaussiano.

\subsection{Sviluppo nella base di Fock}
Di seguito indagheremo la soluzione della \eqref{eq:stati_coerenti} sfruttando una tecnica lievemente diversa dalla \emph{bruta} soluzione dell'equazione differenziale. L'idea è quella di scrivere $\psi_z$ in serie di Fourier sfruttando la base di $\mathcal{L}^2$ fornita dall'hamiltoniana del sistema, ovvero:
\[
    \h{H} = \hbar \omega \left( \frac{1}{2} + \hat{n} \right) \quad,
\]
dove $\hat{n} = \hatd{a}\hat{a}$ è l'operatore numero. Lo spettro in questione è dunque $\mathcal{B} \coloneq \left\{ \psi_n : \hat{n}\psi_n = n\psi_n \right\}$. Iniziamo dall'espansione in serie:
\[
    \psi_z(x) = \sum_{n= 0}^\infty f_n \psi_n(x)
\]
dove con $f_n$ si identificano i coefficienti di Fourier. Sfruttando il fatto che $\psi_z$ è autostato di $\hat{a}$, lo applichiamo ad ambo i membri e otteniamo:
\[
    \begin{aligned}
        \hat{a}\psi_z(x) = z\psi_z(x) &= \hat{a}\sum_{n= 0}^\infty f_n \psi_n(x) = \sum_{n= 0}^\infty f_n \hat{a}\psi_n(x)\\ &= \sum_{n= 0}^\infty f_n \sqrt{n}\psi_{n-1}(x) \giustifica[=]{1} \sum_{n= 1}^\infty f_n \sqrt{n}\psi_{n-1}(x) \giustifica[=]{2} \sum_{n= 0}^\infty f_{n+1} \sqrt{n+1}\psi_{n}(x) \quad,
    \end{aligned}
\]
dove in $(1)$ abbiamo fatto partire la somma da $n=1$ trascurando il primo termine nullo ($n= 0$) e in $(2)$ abbiamo eseguito il cambio di variabile $n' = n+1$. Si noti come $z\psi_z$ altro non è che:
\[
    z\psi_z(x) = z \sum_{n= 0}^\infty f_n \psi_n(x) \quad.
\]
Uguagliando le due serie ottenute si ha:
\[
    \begin{aligned}
        z \sum_{n= 0}^\infty f_n \psi_n(x) = \sum_{n= 0}^\infty f_{n+1} \sqrt{n+1}\psi_{n}(x) \implies \sum_{n= 0}^\infty\left( f_{n+1}\sqrt{n+1} - zf_n \right)\psi_n(x) = 0 \quad.
    \end{aligned}
\]
Perché la serie sia nulla per ogni scelta di $\psi_z(x)$ serve che il termine fra parentesi sia identicamente nullo. Questo permette di definire una relazione ricorsiva per il coefficiente di Fourier $f_n$:
\[
    f_{n+1} =  z\frac{f_n}{\sqrt{n+1}} 
\]
che, iterata a partire da $n = 0$, resistuisce:
\[
    f_n = \frac{z^n}{\sqrt{n!}}f_0 \quad.
\]
Il generico stato coerente assume allora la forma di:
\[  
    \psi_z(x) = \sum_{n=0}^\infty \frac{z^n}{\sqrt{n!}}f_0\, \psi_n(x)
\]
e la costante $f_0$ si può trovare imponendo la normalizzazione di $\psi_z$:
\[
    \begin{aligned}
        (\psi_z,\psi_z) &= 1 = \left( \sum_{n=0}^\infty \frac{z^n}{\sqrt{n!}}f_0\, \psi_n, \sum_{m=0}^\infty \frac{z^m}{\sqrt{m!}}f_0\, \psi_m \right) = 1 \\
        &\implies \sum_{n=0}^\infty\sum_{m=0}^\infty \frac{z^n (z^m)^*}{\sqrt{n!m!}}f_0f_0^*\underbrace{(\psi_n,\psi_m)}_{\delta_{mn}} = |f_0|^2 \sum_{n=0}^\infty \frac{|z|^{2n}}{n!} = |f_0|^2 e^{|z|^2} = 1 \implies f_0 \giustifica[=]{3} e^{\frac{|z|^2}{2}} \quad,
    \end{aligned}
\]
dove abbiamo riconosciuto la delta di Kronecker derivante dal prodotto scalare di due autofunzioni appartenenti alla base normalizzabile dell'operatore numero e in $(3)$ abbiamo assunto, senza perdita di generalità, la fase di $f_0$ come nulla\footnote{Si ricordi che quanto si osserva è il modulo quadro della funzione d'onda, in questo caso proporzionale a $|f_0|^2$. La scleta della fase è dunque assolutamente arbitraria.}. 

La forma finale dello stato coerente è pertanto:
\begin{equation}
        \psi_z(x) = e^{\frac{|z|^2}{2}}  \sum_{n=0}^\infty \frac{z^n}{\sqrt{n!}}\, \psi_n(x) \quad, 
        \label{eq:stati_coerenti_2}
\end{equation}
con $\psi_n$ autofunzione dell'operatore numero:
\[
    \psi_n(x) = \frac{1}{\sqrt{2^n \lambda n!\sqrt{\pi}}}e^{-\frac{x^2}{2\lambda^2}}H_n\left({\frac{x}{\lambda}}\right) \quad.
\]
La base dell'ocillatore armonico viene detta \emph{base di Fock} pertanto la \eqref{eq:stati_coerenti_2} prende il nome di \textbf{espansione nella base di Fock.}

\subsubsection{Funzione generatrice dei polinomi di Hermite}
Abbiamo così trovato due espressioni differenti per definire uno stato coerente, la \eqref{eq:stati_coerenti} e la \eqref{eq:stati_coerenti_2}. Con manipolazioni algebriche dalla difficoltà non indifferente, è possibile dimostrare che la combinazione delle due porti a trovare la cosiddetta \emph{funzione generatrice} dei polinomi di Hermite definiti secondo la teoria dei polinomi ortogonali. 

\subsubsection{Base dell'oscillatore armonico}
Ricordando che le $\psi_n$ formano una base per $\mathcal{L}^2$ poiché compongono lo spettro di un operatore hermitiano ($\hat{n}$), l'espressione ottenuta per $\psi_z(x)$ che fa uso delle serie di Fourier non è altro che la scrittura dello stato coerente nella base dell'oscillatore armonico:
\[
    \psi_z(x) = \sum_n C_n(z)\psi_n(x), \quad (\psi_n, \psi_z) = C_n(z) = e^{-\frac{|z|^2}{2}}\frac{z^n}{\sqrt{n!}} \quad.
\]

Si presti attenzione al fatto che l'insieme degli stati coerenti \textbf{non} è una base per $\mathcal{L}^2$ dal momento che diagonalizza un operatore non hermitiano. In ogni caso, il prodotto scalare di due generici stati coerenti $\psi_z$ e $\psi_\xi$ con $z \neq \xi$ porta a:
\[
    (\psi_z, \psi_\xi) = e^{-\frac{|z|^2+|\xi|^2}{2}}\sum_m \sum_n \frac{(z^*)^m \xi^n}{\sqrt{m!n!}}\underbrace{(\psi_m, \psi_n)}_{\delta_{mn}} =  e^{-\frac{|z|^2+|\xi|^2}{2}} \sum_n \frac{(z^*\xi)^n}{n!} = e^{-\frac{|z|^2+|\xi|^2}{2}} e^{z^* \xi} \neq 0 \quad.
\]
In ogni caso, con $z = \xi$ il prodotto scalare restituisce l'unità, a testimonianza del fatto che la condizione di normalizzazione è preservata. 

\subsection{Evoluzione temporale degli stati coerenti}


%!TEX root = ../dispense_QP.tex

\chapter{Il momento angolare}
La trattazione del momento angolare in termini quantistici è di fondamentale importanza ai fini dello studio di innumerevoli modelli fisici, alcuni fra tutti quelli dell'\emph{atomo di idrogeno}, della \emph{molecola biatomica} e del \emph{moto di una carica in un campo magnetico}. Come anticipato nella sezione introduttiva sugli operatori, essendo a tutti gli effetti il momento angolare un'osservabile fisica, possiamo definire $\hb{L} = (\hat{L}_1, \hat{L}_2, \hat{L}_3)$ come \emph{operatore momento angolare}. 

\section{Richiami e definizione}
Senza ripetere quanto già detto, si definisce operatore momento angolare:
\[
    \hb{L} = \hb{x} \times \hb{p} \implies \hat{L}_i = \varepsilon_{ijk}x_jp_k = -i\hbar \varepsilon_{ijk} x_j \pdv{}{x_k} \quad.
\]
Si noti come il commutatore delle componenti dell'operatore momento angolare è \emph{ciclico}:
\[
    [\hat{L}_i, \hat{L}_j] = i\hbar \varepsilon_{ijk}\hat{L}_k \quad,
\]
a testimonianza del fatto che l'insieme $\mathcal{M} \coloneq \left\{ \hat{L}_i, i \in 1,2,3 \right\}$ genera un Algebra di Lie. Sfruttando questa peculiarità possiamo costruire un'ulteriore struttura algebrica con le medesime caratteristiche. A tale fine definiamo gli operatori:
\[
    \hat{L}_+ = \hat{L}_1 + i\hat{L}_2, \quad \hat{L}_- = \hat{L}_1 - i\hat{L}_2 
\]
che saranno molto utili nel definire lo spettro di $\hat{L}_3$\footnote{Si noti che la scelta di definire $\hat{L}_\pm$ tramire $\hat{L}_{1,2}$ nasce esclusivamente dall'esigenza di studiare lo spettro di $\hat{L}_3$ perché, complessivamente, decisamente più semplice rispetto agli altri due. }. 

\subsection{Caratteristiche di \texorpdfstring{$\bm{\hat{L}_{\pm}}$}{L pm}} 
Si preannuncia che la finalità di quanto segue è quella di arrivare a costruire, tramite il metodo dell'algebra generatrice di spettro, l'insieme degli autostati di $\hat{L}_3$. Si noti come, essendo:
\[
    [\hat{L}_3, \hat{L}^2] = 0 \quad,
\]
è sempre possibile definire una base comune di autostati. Assumeremo che lo spettro che troveremo per $\hat{L}_3$ diagonalizzi anche $\hat{L}^2$ senza porci particolari domande\footnote{Una giustificazione piuttosto euristica sarà fornita al termine della derivazione.}. 

Partiamo dal calcolare il commutatore di dei due nuovi operatori:
\[
    \begin{aligned}
        [\hat{L}_+, \hat{L}_-] = [\hat{L}_1 + i\hat{L}_2, \hat{L}_1 - i\hat{L}_2] = 2i[\hat{L}_2, \hat{L}_1] = - 2i (i\hbar)\hat{L}_3 = \hbar \hat{L}_3 \quad.
    \end{aligned}
\]
Il commutatore con $\hat{L}_3$, invece:
\[
    \begin{aligned}
        [\hat{L}_3, \hat{L}_\pm] = [\hat{L}_3, \hat{L}_1 \pm i\hat{L}_2] = [\hat{L}_3, \hat{L}_1] \pm i[\hat{L}_3, \hat{L}_2] = i\hbar\hat{L}_2 \mp i (i\hbar) \hat{L}_1 = \pm\hbar(\hat{L}_1 \pm i\hat{L}_2) = \pm \hbar \hat{L}_{\pm} \quad.
    \end{aligned}
\]  
L'insieme $\mathcal{N} \coloneq \left\{ \hat{L}_+, \hat{L}_-, \hat{L}_3 \right\}$ definisce quindi un'algebra di Lie e ben si presta, come nel caso dell'oscillatore armonico, all'applicazione del metodo dell'algebra generatrice di spettro\footnote{Il parallelo con il caso dell'oscillatore armonico vuole: $\hat{a} \to \hat{L}_-$, $\hatd{a} \to \hat{L}+-$, $\hat{n} \to \hat{L}_3$. Si noti che esiste però un'enorme differenza fra le due situazioni. Se, infatti, il commutatore fra gli operatori di creazione e annichilazione è una costante, per altro unitaria, nel caso del momento angolare il commutatore restituisce il terzo componente dell'algebra.}.

\section{Lo spettro di \texorpdfstring{$\bm{\hat{L}_3}$}{L3}}
L'importante di trovare lo spettro della terza componete (componente $z$) del momento angolare risiede nel fatto che, nei problemi a simmetria sferica come quello dell'atomo di idrogeno, si conoscono questi due commutatori fondamentali:
\[
    [\h{H}, \hat{L}^2] = 0 = [\hat{L}^2, \hat{L}_3] \quad.
\]
Diagonalizzare $\hat{L}_3$, quindi, significa diagonalizzare $\hat{L}^2$ e quindi, almeno in parte, l'hamiltoniana. 


\appendix
% \include{appendici/formularis}

\end{document}